% Options for packages loaded elsewhere
\PassOptionsToPackage{unicode}{hyperref}
\PassOptionsToPackage{hyphens}{url}
\documentclass[
]{article}
\usepackage{xcolor}
\usepackage{amsmath,amssymb}
\setcounter{secnumdepth}{-\maxdimen} % remove section numbering
\usepackage{iftex}
\ifPDFTeX
  \usepackage[T1]{fontenc}
  \usepackage[utf8]{inputenc}
  \usepackage{textcomp} % provide euro and other symbols
\else % if luatex or xetex
  \usepackage{unicode-math} % this also loads fontspec
  \defaultfontfeatures{Scale=MatchLowercase}
  \defaultfontfeatures[\rmfamily]{Ligatures=TeX,Scale=1}
  % unicode-math maps \setminus to U+29F5 (REVERSE SOLIDUS OPERATOR), which
  % latinmodern-math.otf lacks (tectonic/lualatex emit "Missing character
  % U+29F5"). Remap to U+2216 (SET MINUS) -- the glyph the font provides and
  % that matches the pdfTeX/amssymb rendering. Output-preserving; silences the
  % warning. Applies to all \setminus occurrences globally.
  \AtBeginDocument{\renewcommand{\setminus}{\mathbin{\char"2216}}}
\fi
\usepackage{lmodern}
\ifPDFTeX\else
  % xetex/luatex font selection
\fi
% Use upquote if available, for straight quotes in verbatim environments
\IfFileExists{upquote.sty}{\usepackage{upquote}}{}
\IfFileExists{microtype.sty}{% use microtype if available
  \usepackage[]{microtype}
  \UseMicrotypeSet[protrusion]{basicmath} % disable protrusion for tt fonts
}{}
\makeatletter
\@ifundefined{KOMAClassName}{% if non-KOMA class
  \IfFileExists{parskip.sty}{%
    \usepackage{parskip}
  }{% else
    \setlength{\parindent}{0pt}
    \setlength{\parskip}{6pt plus 2pt minus 1pt}}
}{% if KOMA class
  \KOMAoptions{parskip=half}}
\makeatother
\setlength{\emergencystretch}{3em} % prevent overfull lines
\providecommand{\tightlist}{%
  \setlength{\itemsep}{0pt}\setlength{\parskip}{0pt}}
\usepackage{bookmark}
\IfFileExists{xurl.sty}{\usepackage{xurl}}{} % add URL line breaks if available
\urlstyle{same}
\hypersetup{
  hidelinks,
  pdfcreator={LaTeX via pandoc}}

\author{}
\date{}

\begin{document}

\textbf{Ratushnyi Neural Repair Coding, Part III:}

\textbf{Operational Rate--Distortion Dispersion, the Gray Region, and the
Replica Spectral Inequality}

Pavlo Ratushnyi

\emph{Independent Research --- Draft v0.8, 2026}

\emph{Part III of the unified series ``Ratushnyi Neural Repair Coding''.}

\textbf{Abstract}

This Part develops the second-order (dispersion) theory of \emph{lossy}
Ratushnyi Neural Repair (RNR) coding. Working under Hamming distortion $D$
and the $d$-tilted information $\jmath_X(x;D)$, we first isolate what
transfers from the lossless second-order rate and what does not: for a
memoryless source the i.i.d.\ rate-distortion dispersion $V(D)=
\mathrm{Var}(\jmath_X(X;D))$ concentrates the repair stream at
$NR(D)+\sqrt{NV(D)}\,Q^{-1}(\varepsilon)$, whereas the \emph{cold-context}
random-access interplay of the lossless companion result does not survive
into the memoryless regime --- an obstruction we record honestly rather than
force. We then treat sources \emph{with memory}. The central object is the
\emph{Gray region}: the range of distortion, source, and Fourier-angle
parameters on which the operational lossy dispersion equals the lossless
varentropy, $V_{\mathrm{op}}=V_{\mathrm{lossless}}$. We prove a
source-agnostic converse through a dual spectral identity
$K_\nu K^{-1}=\gamma I$ (the tilted transfer operator is a scalar multiple
of the identity on the optimal reproduction law), giving
$V_{\mathrm{op}}\ge V_{\mathrm{lossless}}$ unconditionally on the region.
Around this we assemble the dispersion constant, the convexity wall that
forces strict positivity of the leading second-order coefficient, the
third-order $\tfrac12\log n$ refinement on the Gray grid via the operational
recentring $j_n=i_n-nc$, the beyond-Gray landscape, and a lossy deviation
hierarchy transferring the large-deviation, moderate-deviation,
extreme-value, and functional-CLT endpoints of the companion Part. Finally
we prove the replica spectral inequality $R_A(s;D)>3/2$ underlying the
convexity wall --- closed form at $A=2$, with exact certificates at
$A=3,4,5$.

\textbf{Keywords:} lossy source coding; rate-distortion dispersion;
$d$-tilted information; operational second-order rate; finite blocklength;
Gray region; varentropy; transfer operator; spectral radius convexity;
replica spectral inequality; neural repair coding.

\textbf{0. Numbering, companion parts, and how to read this document}

\emph{Numbering follows the unified RNR series; this Part contains
Section~7.34 (Remark 7.34 through Remark 7.34q) together with the
prerequisites restated in Section~7.3 below. Section, theorem, lemma, and
remark numbers are those of the unified series and are \textbf{not}
renumbered here. Results proved in a companion Part are cited as
[RNR-I]/[RNR-II] (Parts~I/II): [RNR-I] is Part~I (framework, type taxonomy,
achievability--converse matrix, and hardness); [RNR-II] is Part~II (random
access and deviation theory). A first mention of a companion result carries
the tag ``of [RNR-II]'' (or ``[RNR-II, \S7.15]''); subsequent mentions keep
the bare number, which is unambiguous because the numbering is shared across
the series. The one internal cross-reference label of the whole series
(\texttt{eq:G1chain}) lives in this Part.}

\textbf{7.3 Prerequisites (restated from the companion Parts)}

This section collects, statement-only and with provenance tags, the results
this Part builds on. Proofs are in the cited companion Part.

\emph{RNR recap ([RNR-I]).} The RNR \emph{repair principle} codes a source
stream $X=(X_1,\dots,X_N)$ not by a monolithic entropy coder but by a
predictor-driven \emph{repair stream}: a side-information model $\Theta$
(shared, source-independent) proposes each symbol and the repair stream
carries only the residual information $-\log_2 Q_i(X_i\mid\mathrm{ctx}_i)$
needed to correct it. The stream is partitioned into $m=N/K$ contiguous
\emph{sub-blocks} of spacing $K$, each decodable from its own payload plus a
short \emph{sync} record; this is the architecture that makes byte-granular
random access possible (a query landing in sub-block $i$ reads only that
block's repair bits). The \emph{repair length}
$L^{\mathrm{rep}}(X)=\sum_i -\log_2 Q_i(X_i\mid\mathrm{ctx}_i)$ (plus $O(1)$
per sub-block for sync/flush) is the quantity whose mean, tail, and
distribution the deviation theory analyses; its first-order value is $Nh+o(N)$
for a source of entropy rate $h$ under an ideal predictor. The lossy variant
([RNR-I], and Theorem 7.7 below) replaces the per-sub-block lossless coder by
a Shannon-optimal rate-distortion encoder at rate $R(D)$.

\emph{The lossy operational model (Theorem 7.7 of [RNR-I]/[RNR-II]).} Let
$X=(X_n)_{n\ge1}$ be stationary, exp-mixing, finite-alphabet, with sub-block
spacing $K=\mathrm{polylog}(N)$, and let $d_H:\Sigma\times\Sigma\to\{0,1\}$ be
Hamming distortion. For a per-position budget $D\in[0,1-1/|\Sigma|]$ let
$R(D):=\lim_n\tfrac1n\min_{P_{\hat X^n\mid X^n}:\,\mathbb E[d_H(X^n,\hat X^n)/n]\le D}
I(X^n;\hat X^n)$ be the Shannon rate-distortion function. RNR with lossy
per-sub-block rate-distortion coding achieves expected repair-stream length
$\mathbb E[L^{\mathrm{rep}}_{\mathrm{RNR},D\text{-lossy}}(X^N)]\le
N\,R(D)+m\,\delta_\infty^D+O(1)+\varepsilon_{\mathrm{ac}}^{\mathrm{total}}$
with per-position distortion $\mathbb E[d_H(X^N,\hat X^N)/N]\le D$, where
$\delta_\infty^D$ is the $D$-distortion analogue of the Crutchfield--Feldman
excess (finite under exp-mixing). Limit cases: $D\to0$ recovers the lossless
bound; $D\to1-1/|\Sigma|$ gives $R(D)\to0$. This is the first-order object
whose \emph{second-order} refinement is the subject of this Part.

\textbf{Theorem 7.28 ([RNR-II, Thm 7.28]; Non-asymptotic CLT for
$L^{\mathrm{rep}}$: Berry--Esseen rate under exp-mixing).} \emph{Let
$X=(X_1,\ldots,X_N)$ be stationary $\alpha$-mixing with
$\alpha_X(k)\leq C_\alpha e^{-c k}$; per-position cost
$\ell_i:=-\log_2 Q_i(X_i\mid\mathrm{ctx}_i)$ bounded,
$0\leq\ell_i\leq B=\log_2(1/\eta)$, depending on a window of $W_N$ coordinates
of $X$; with marginal variance $\sigma_\ell^2:=\mathrm{Var}(\ell_1)$ and
long-run variance $v^2:=\sigma_\ell^2+2\sum_{k\geq1}\mathrm{Cov}(\ell_0,\ell_k)$,
finite and strictly positive. Let $\Phi$ be the standard normal CDF and
$\kappa:=m_3/(B\,v^2)$ with $m_3:=\mathbb E|\ell_1-\mathbb E\ell_1|^3$. Then
there is a constant $C_{\mathrm{BE}}$, depending only on $(c,C_\alpha)$ and on
$\kappa$, such that for all $N\geq2$,}
\[
  \sup_{x\in\mathbb R}
  \left|
    \Pr\!\left(
      \frac{L^{\mathrm{rep}}(X)-\mathbb E[L^{\mathrm{rep}}(X)]}{\sqrt{N\,v^2}}
      \leq x
    \right)-\Phi(x)
  \right|
  \;\leq\;
  C_{\mathrm{BE}}\cdot\frac{B\,(\log N)^2}{v\,\sqrt N}.
\]
\emph{Equivalently, the Kolmogorov distance between the law of the normalised
repair length and the standard Gaussian is $O(N^{-1/2}(\log N)^2)$. (The
$(\log N)^2$ factor is genuine for strongly mixing summands --- Tikhomirov
1980 --- and cannot be removed in the Kolmogorov metric.)}

\textbf{Theorem 7.29 ([RNR-II, Thm 7.29]; Second-order coding rate of RNR;
the random-access / dispersion interplay).} \emph{Adopt the hypotheses of
Theorem 7.28 with the ideal predictor $Q_i=P(\cdot\mid X_{i-W_N:i-1})$ on a
context window $W_N=\Theta(\log N)$, so $\ell_i\to-\log_2 P(X_i\mid X_{<i})$
and the long-run variance $v^2$ equals the source \emph{varentropy}
$V:=\lim_N\mathrm{Var}(-\log_2 P(X^N))/N$ (finite, positive). Let the sub-block
spacing be $K$ and let $\delta_\infty>0$ be the per-sync-point excess overhead
(Crutchfield--Feldman), so the sync cost is $m\,\delta_\infty=(N/K)\delta_\infty$.
Write $Q^{-1}(\varepsilon):=\Phi^{-1}(1-\varepsilon)$ for the Gaussian
tail-quantile. Then:}
\begin{enumerate}
\item \emph{(Achievability.) For every $\varepsilon\in(0,1)$, with probability
at least $1-\varepsilon$,
$L^{\mathrm{rep}}_{\mathrm{RNR}}(X^N)\le
Nh+\tfrac NK\delta_\infty+\sqrt{NV}\,Q^{-1}(\varepsilon)+o(\sqrt N)$,}
\emph{the $\sqrt{NV}\,Q^{-1}(\varepsilon)$ term coming from the uniform
Berry--Esseen CLT of Theorem 7.28 and the mean from the first-order rate.}
\item \emph{(Dispersion floor --- converse.) For a memoryless or finite-state
irreducible aperiodic Markov source the varentropy $V$ is the source
dispersion: under the excess-length criterion the optimal fixed-to-variable
codelength obeys $L^\star(N,\varepsilon)\ge Nh+\sqrt{NV}\,Q^{-1}(\varepsilon)
-o(\sqrt N)$ (Kontoyiannis--Verd\'u 2014), so $\sqrt{NV}\,Q^{-1}(\varepsilon)$
cannot be removed.}
\item \emph{(The interplay.) At the balanced spacing $K=c\sqrt N$ the
random-access overhead and the dispersion are both $\Theta(\sqrt N)$,
$L^{\mathrm{rep}}_{\mathrm{RNR}}=Nh+\sqrt N[\tfrac{\delta_\infty}{c}
+\sqrt V\,Q^{-1}(\varepsilon)]+o(\sqrt N)$, with their ratio independent of
$N$.}
\end{enumerate}
\emph{The per-sync-point constant $\delta_\infty$ (the \emph{cold-context}
excess $I(\text{past};\text{future})$ paid at each block boundary because a
random-access decoder starts each sub-block without the preceding context) is
the quantity whose lossy analogue $\delta_\infty^D$ organises this Part: it
vanishes for a memoryless source and is the obstruction of Remark 7.34.}

\emph{Deviation-hierarchy pointers ([RNR-II]).} Beyond the second order, the
companion Part establishes the large-deviation overflow exponent (Theorem 7.31
of [RNR-II]; Markov endpoint $=$ max cycle mean), the moderate-deviation
random-access fade (Theorem 7.33 of [RNR-II]; the RA penalty is one
standardized shift $\beta$ that fades), the extreme-value law for the
worst-case random-access query (Theorem 7.35 of [RNR-II]; read-buffer margin
$Kh+\sqrt{2KV\ln m}$), and the functional CLT for the streaming buffer
(Theorem 7.36 of [RNR-II]; the repair process converges to Brownian motion).
Remarks 7.34p--q below transfer these four endpoints to the lossy regime.

\emph{End of prerequisites. The remainder of this document is the Section~7.34
family, verbatim from the unified series.}

\emph{Remark 7.34 (Lossy RNR at second order: the i.i.d.\ dispersion transfers, but the
\emph{cold-context} random-access interplay of Theorem 7.29 of [RNR-II] does not --- an honest
obstruction).} For a
\emph{memoryless} source coded at Hamming distortion $D$, the per-symbol lossy codelength
is the $d$-tilted information $\jmath_X(x;D)=-\lambda^\ast D-\log_2\mathbb E_{Y^\ast}[2^{-
\lambda^\ast d_H(x,Y^\ast)}]$ (all in bits: $\lambda^\ast=-R'(D)>0$ with $R$ in bits,
$Y^\ast$ the optimal reproduction marginal; Kostina--Verd\'u 2012, \emph{IEEE Trans.\
Inf.\ Theory}\ 58(6):3309--3338), with $R(D)=\mathbb E[\jmath_X(X;D)]$ and rate-distortion
dispersion $V(D)=\mathrm{Var}(\jmath_X(X;D))$. The lossy repair stream then concentrates,
with probability $1-\varepsilon$, as
\[
  L^{\mathrm{rep}}_{\mathrm{RNR},D\text{-lossy}}
  =\sum_n\jmath_X(X_n;D)+O\!\bigl(\tfrac NK\bigr)
  =N R(D)+\sqrt{N V(D)}\,Q^{-1}(\varepsilon)+O\!\bigl(\tfrac NK\bigr)+o(\sqrt N),
\]
the i.i.d.\ CLT for $\jmath_X$ being the lossy analogue of Theorem 7.29's
information-density CLT, and the $O(N/K)$ the $m=N/K$ sub-block flush/sync overhead. At
$K=\Theta(\sqrt N)$ this flush overhead and the dispersion are both $\Theta(\sqrt N)$ ---
a genuine (if milder) random-access/dispersion interplay even in the memoryless case ---
but it is the \emph{flush} term, not a cold-context $\delta_\infty^D$, that supplies it
(see below). For the binary memoryless source the
$d$-tilted information difference $\jmath_X(0;D)-\jmath_X(1;D)=\log_2\tfrac{p}{1-p}$ is
$D$-\emph{independent}, so $V(D)=p(1-p)\log_2^2\tfrac{1-p}{p}$ is constant in $D$ and equals
the lossless varentropy $V$ of Theorem 7.29 --- a clean degeneracy (lossy dispersion $=$
source varentropy), continuous with the lossless case as $D\to0$; general sources have
$V(D)$ genuinely $D$-dependent (a ternary source by Blahut--Arimoto; verification script).

\emph{The interplay does not transfer (honest).} Theorem 7.29's signature --- the
cold-context random-access penalty $\tfrac NK\delta_\infty^D$ at the \emph{same} $\sqrt N$
scale as the dispersion --- requires a source \emph{with memory}: for a memoryless source
the Crutchfield--Feldman excess $\delta_\infty^D=I(\text{past};\text{future})=0$, so there
is no cold-context penalty at all --- only the milder flush overhead $O(N/K)$ above, which
sits at the same $\sqrt N$ scale as the dispersion $\sqrt{N V(D)}\,Q^{-1}(\varepsilon)$.
But once the source has memory, $V(D)$ is no
longer the single-letter $\mathrm{Var}(\jmath_X)$: it becomes the variance of the
\emph{process} $d$-tilted information --- the rate-distortion dispersion of a source with
memory --- which (unlike the lossless long-run varentropy that Theorem 7.28 of [RNR-II] supplies
cleanly via the exp-mixing CLT) is delicate and not settled here. So the clean lossy
analogue of the Theorem 7.29 interplay is \emph{open}; we record the obstruction rather
than force a single-letter $V(D)$ into a memory setting where it fails. The missing
ingredient is precisely a \emph{mixing-source rate-distortion dispersion}, the lossy
counterpart of the Theorem 7.28 varentropy. Verified (the i.i.d.\ part) in
\texttt{scripts/verify/thm\_7\_34\_lossy\_rnr\_dispersion.py} (V1 $R(D)=\mathbb E[\jmath]
=h(p)-h(D)$; V2 BMS $V(D)$ constant $=$ varentropy; V3 ternary $V(D)$ $D$-dependent; V4
$\delta_\infty^D=0$ for i.i.d.\ so the dispersion is the leading second-order term; V5
$D\to0$ recovers Theorem 7.29).

\emph{Remark 7.34a (The missing ingredient exists for a memory source: the Gauss--Markov
rate-distortion dispersion realises the lossy second-order interplay).} Remark 7.34 left the
lossy analogue of the Theorem 7.29 interplay open because its missing ingredient --- a
\emph{rate-distortion dispersion for a source with memory} --- was unsettled. That ingredient
\emph{is} available for one canonical memory source: the stationary Gauss--Markov (AR$(1)$)
source $X_t=\rho X_{t-1}+W_t$, $|\rho|<1$, $W_t\sim\mathcal N(0,\sigma_W^2)$, under
mean-squared distortion, for which Tian and Kostina (\emph{The Dispersion of the Gauss--Markov
Source}, IEEE Trans.\ Inf.\ Theory 65(10):6355--6384, 2019) established the operational
second-order rate $R(n,D,\varepsilon)=R(D)+\sqrt{V(D)/n}\,Q^{-1}(\varepsilon)+o(n^{-1/2})$
with $R(D),V(D)$ given by reverse-water-filling on the PSD
$S(\omega)=\sigma_W^2/(1+\rho^2-2\rho\cos\omega)$:
\[
  R(D)=\frac1{2\pi}\!\int_{-\pi}^{\pi}\!\tfrac12\log^{+}\!\tfrac{S(\omega)}{\theta}\,d\omega,
  \quad D=\frac1{2\pi}\!\int_{-\pi}^{\pi}\!\min\{\theta,S(\omega)\}\,d\omega,
  \quad
  V(D)=\frac1{4\pi}\!\int_{-\pi}^{\pi}\!\min\Bigl\{1,\bigl(\tfrac{S(\omega)}{\theta}\bigr)^2\Bigr\}d\omega
\]
(nats$^2$; multiply by $(\log_2 e)^2$ for bits), and crucially identified this with the
\emph{informational} dispersion --- the limiting variance of the block $d$-tilted information
$\jmath_n(X^n;D)$ --- so the operational and information-density second orders coincide exactly
as in the lossless Theorem 7.28/7.29. For low distortion $D\le D_c:=\min_\omega S(\omega)=
\sigma_W^2/(1+|\rho|)^2$ one has $\theta=D$ and $V(D)=\tfrac12$ (nats$^2$), the i.i.d.\
innovation value; for $D_c<D<D_{\max}=\sigma_W^2/(1-\rho^2)$, $V(D)<\tfrac12$.

\emph{The lossy interplay holds for this memory source.} The AR$(1)$ source is genuinely
mixing with memory ($I(\text{past};\text{future})>0$), so a sub-block sync reset incurs a
\emph{positive} cold-context lossy excess: at high rate the first reproduction after a reset
must be coded from the marginal (variance $\sigma_W^2/(1-\rho^2)$) rather than the AR
conditional (innovation variance $\sigma_W^2$), giving
$\delta_\infty^D=-\tfrac12\log(1-\rho^2)\,(\log_2 e)>0$ bits, increasing in $|\rho|$ and
$\to0$ as $\rho\to0$ (the i.i.d.\ limit, where Remark 7.34's transfer already applies). Hence
the lossy RNR repair stream obeys, with probability $1-\varepsilon$,
\[
  L^{\mathrm{rep}}_{\mathrm{RNR},D}=N R(D)+\tfrac NK\,\delta_\infty^D
  +\sqrt{N\,V(D)}\,Q^{-1}(\varepsilon)+o(\sqrt N),
\]
and at the balanced spacing $K=c\sqrt N$ the cold-context random-access overhead
$\tfrac NK\delta_\infty^D=\Theta(\sqrt N)$ and the rate-distortion dispersion
$\sqrt{NV(D)}\,Q^{-1}(\varepsilon)=\Theta(\sqrt N)$ are of the \emph{same order}, with
$N$-independent ratio $\delta_\infty^D/(c\sqrt{V(D)}\,Q^{-1}(\varepsilon))$ --- the exact
Theorem 7.29 random-access/dispersion interplay, now \emph{lossy} and for a source \emph{with
memory}. So the interplay Remark 7.34 could only state for the (penalty-free) memoryless case
is non-vacuous: it is realised, with both terms genuinely present and comparable, on the
Gauss--Markov source.

\emph{What stays open (honest localisation).} This supplies Remark 7.34's missing ingredient
only for the \emph{Gaussian} memory class (continuous source, MSE distortion); RNR's native
lossy setting --- a \emph{discrete} finite-state source under \emph{Hamming} distortion ---
still lacks an operational rate-distortion dispersion. The obstruction is structural: the
Markov RDF is not single-letter (it is an $n$-letter limit $R(D)=\lim_n R_n(D)$, known in
closed form only in special low-distortion regimes; Gray 1970s, Jalali--Weissman bounds), so
the operational $d$-tilted information does \emph{not} single-letterise as it does for
memoryless (Kostina--Verd\'u) or Gaussian (Tian--Kostina) sources, and its dispersion is open.
The single-letter marginal (Blahut--Arimoto) $d$-tilted proxy does have a computable Markov
long-run variance (Krishnamachari 2026, $V_{\mathrm{sl}}=\tfrac{ab(2-a-b)}{(a+b)^3}\log_2^2
\tfrac ab$ for the binary chain), but it is provably distinct from the $n$-letter operational
quantity, whose dispersion remains unidentified. (Continuous Gauss--Markov data enters
RNR's scope only through fine quantisation; the abstract second-order machinery of
Theorems 7.29/7.36 is source-agnostic once the dispersion exists, which is exactly what
Tian--Kostina provide here.) Thus Remark 7.34's open problem narrows from ``all sources with
memory'' to the \emph{discrete-Markov, Hamming-distortion} case specifically. The natural
candidate here --- a $D$-independent plateau at the chain's \emph{lossless} varentropy
$p(1-p)\log_2^2\tfrac{1-p}p$ across the Gray region, mirroring the Gauss--Markov $V(D)=\tfrac12$
plateau --- is \emph{proven} in Remark 7.34b below: $V_{\mathrm{op}}=V_{\mathrm{lossless}}$ throughout
the closed Gray region $0<D\le D_c(p)$ (the converse for all $n$; the achievability as
an asymptotic theorem with explicit constants, the endpoint included via a dedicated verification
round), the first discrete memory case resolved --- the general
discrete-Markov/Hamming dispersion beyond the BSMS remains open (Krishnamachari 2026). Only
\emph{outside} the Gray region
($D>D_c(p)$) does even the converse value fall below $V_{\mathrm{lossless}}$, a separate open question.
Verified in
\texttt{scripts/verify/remark\_7\_34a\_gauss\_markov\_rd\_dispersion.py} (V1 reverse-water-filling
self-consistency; V2 $V(D)=\tfrac12$ nats$^2$ for $D\le D_c$; V3 $V(D)<\tfrac12$ decreasing for
$D_c<D<D_{\max}$; V4 the block $d$-tilted-information variance $V_n\to$ the spectral $V(D)$
(AR$(1)$ Toeplitz covariance eigenvalues, cross-checked by Monte-Carlo of $\jmath_n$); V5 the
$\Theta(\sqrt N)$ interplay with $N$-independent ratio; V6 $\delta_\infty^D=-\tfrac12\log(1-\rho^2)>0$,
increasing in $|\rho|$, $\to0$ as $\rho\to0$).

\textbf{Remark 7.34b (The BSMS Gray-region RD-dispersion theorem:
$V_{\mathrm{op}}=V_{\mathrm{lossless}}$ --- the converse for all $n$, and a fully written asymptotic
achievability).} For the binary symmetric Markov
source (BSMS, switch probability $0<p<\tfrac12$) under Hamming distortion, we prove
\[
  V_{\mathrm{op}}(D)= V_{\mathrm{lossless}}=p(1-p)\log_2^2\tfrac{1-p}p
  \qquad\text{throughout the closed Gray region }0<D\le D_c(p),
\]
with Gray's critical distortion $D_c(p)=\tfrac12\big(1-\sqrt{1-2p}/(1-p)\big)$ (Gray 1970;
Zhu--Alajaji) --- exactly the candidate value of Remark 7.34a. The \emph{converse}
$V_{\mathrm{op}}\ge V_{\mathrm{lossless}}$ holds throughout the closed region $0<D\le D_c(p)$ at every
blocklength (the SLB-exact identity below); the matching \emph{achievability} is established by the
chain developed in this remark --- the dual/polymer representation, the proven $4\times4$ replica
spectral inequality and its unconditional quenched tail lemma, the ordered-constants central-window
assembly, and the per-word steepest-descent (GAP-1) lemma, all written out below with closed-form
constants --- as an \emph{asymptotic} statement (explicit but large $n_0(p,D)$; the $nD\ge5$ regime;
see the scope notes and the concluding theorem below; the endpoint $D=D_c$ is included via the
M\"obius-orbit floor lemma and a dedicated endpoint verification round). This resolves the
discrete-Markov (BSMS,
Hamming) dispersion throughout Gray's region --- the first discrete memory case, the general
discrete-Markov/Hamming dispersion beyond the BSMS remaining open (Krishnamachari 2026;
Tasci--Kostina 2026 list Gauss--Markov as the only previously solved memory case).

\emph{The key identity and the converse.} Solving the $n$-letter rate--distortion problem exactly
(Blahut--Arimoto over the $2^n$-word alphabet) and evaluating the operational $n$-letter d-tilted
information $\jmath_n(x^n,D)=-\lambda D-\log\mathbb{E}_{Y^{*n}}[e^{-\lambda d_H(x^n,Y^{*n})}]$
($\mathbb{E}[\jmath_n]=R_n(D)$ to machine precision): on the \emph{whole} Gray region $0<D\le D_c(p)$
the i.i.d.\ $\mathrm{BSC}(D)$ deconvolution of the source law is valid at \emph{every} blocklength $n$
(mechanism below), so the Shannon lower bound is exactly tight for all $n$ and the finite-block SLB
identity holds,
\[
  \jmath_n(x^n,D)=\imath_n(x^n)-n\,h(D),\qquad \imath_n:=-\log_2 P(X^n).
\]
Since $\imath_n$ is an affine function of the i.i.d.\ switch-count $N_{\mathrm{sw}}=\sum_i\mathbf 1[X_i\ne X_{i-1}]$
(because $P(x^n)=\tfrac12(1-p)^{\mathrm{stays}}p^{\mathrm{switches}}$), $\jmath_n$ is an \emph{exact
i.i.d.-additive functional}, with $\mathrm{Var}(\jmath_n)/n\to V_{\mathrm{lossless}}$. By the
one-shot d-tilted converse of Kostina--Verd\'u 2012 (Thm.~7, excess-distortion criterion; a plain CLT
for $\jmath_n$ suffices for the fixed-$\varepsilon$ liminf), $V_{\mathrm{op}}(D)\ge\lim_n\mathrm{Var}(\jmath_n)/n=V_{\mathrm{lossless}}$
on $(0,D_c(p)]$. (For $D>D_c(p)$, outside the Gray region, the identity breaks and the converse ratio
$\rho(D):=\mathrm{Var}(\jmath_n)/\mathrm{Var}(\imath_n)$ drops $n$-stably into $(0,1)$ --- e.g.\
$\rho\approx0.61$ at $p=0.1,\,D=0.05$ --- the separate regime noted at the end.)

\emph{The mechanism: $D_c$ is the all-$n$ SLB-deconvolution boundary.} The SLB-exact identity
$\jmath_n=\imath_n-n\,h(D)$ holds iff the source law deconvolves by an i.i.d.\ $\mathrm{BSC}(D)$ noise
into a valid output law. Because Hamming distortion is additive, the d-tilted backward kernel
factorises, $e^{-\lambda d_H(x^n,y^n)}=\prod_i e^{-\lambda\,\mathbf 1[x_i\ne y_i]}$, into single-letter
$\mathrm{BSC}$ kernels, so the SLB-achieving backward channel is \emph{forced memoryless}; the
deconvolution in the Walsh (Hadamard) domain is $\hat P_Y(w)=\hat P_X(w)\,(1-2D)^{-|w|}$. From the
transfer recursion ($(1+\alpha)^2-4\alpha(1-2D)^{-2}\ge0$, $\alpha=1-2p$) its all-$n$ positivity
boundary is exactly $\tfrac12\big(1-\sqrt{1-2p}/(1-p)\big)$, which \emph{coincides with Gray's critical
distortion} $D_c(p)$ (verified across $p$). Thus on the \emph{entire} Gray region $0<D\le D_c(p)$ the
deconvolution is valid for every $n$ and the SLB identity holds for all $n$, yielding the converse
above. (For $D>D_c(p)$ the inverse transform acquires negative entries for large $n$ --- e.g.\ at
$p=0.1,\,D=0.05$ valid only at $n=2$, failing for every $n\ge3$ --- so the identity fails and
$V_{\mathrm{conv}}<V_{\mathrm{lossless}}$; but that is outside the Gray region.)

\emph{The achievability residual, correctly located at the dispersion level.} The
matching upper bound $V_{\mathrm{op}}\le V_{\mathrm{lossless}}$ remained open prior to the chain
developed below (the general discrete-Markov/Hamming dispersion beyond the BSMS is still open ---
Krishnamachari 2026 leaves it
explicitly open; before this remark only the Gauss--Markov memory class was resolved --- Tian--Kostina
2019, Tasci--Kostina 2026 --- via $n$-letter informational quantities). Two facts locate the gap. \emph{(a)
The leading dispersion is a CDF Berry--Esseen object; the achievability residual is a (near-mean)
local-limit object.} The $\sqrt{V/n}$ coefficient is carried by the additive d-tilted information
$\jmath_n$ through a CDF Berry--Esseen (lattice-agnostic, ellipticity-free --- this is what the converse
uses). The Kostina--Verd\'u constant-width window $\{nD-\tau<\sum_i d(X_i,\hat Z_i)\le nD\}$, however, is
a \emph{windowed local mass} $F(nD)-F(nD-\tau)\sim\tau\,f(nD)=O(n^{-1/2})$ --- a difference of CDF values
over an $O(1)$ window, i.e.\ an Edgeworth/local-limit quantity, not a plain CDF value --- so matching the
\emph{memory} block-optimal ball to $\jmath_n$ (forcing $V_{\mathrm{op}}=V_{\mathrm{lossless}}$) needs
that local mass: a near-mean lattice local-limit theorem for the optimal-coding distortion sum, uniform
over the $\sqrt n$-shell. (Only the sharp $\tfrac12\log n$ \emph{coefficient} below is genuinely
third-order and dispensable --- being deterministic, it drops from the variance --- not the local-limit
\emph{machinery} that produces it, which the dispersion does need.) \emph{(b) The optimal kernel is temporally correlated:} for a Markov source
the $n$-letter block-optimal test channel is \emph{not} a single-letter tilt of the marginal
(Krishnamachari 2026), so the single-letter $d$-tilted information does not carry the operational
dispersion and the relevant object is the genuinely $n$-letter forward-filter cocycle --- non-elliptic
(atomic filter, no density floor) and schedule-inhomogeneous. Its CDF Berry--Esseen, uniform over the
shell, holds \emph{homogeneously and density-floor-free even on singular laws}: Hennion--Herv\'e, and
sharply Kloeckner (\emph{Ann.\ Appl.\ Probab.}\ 29(3), 2019; arXiv:1703.09623, Thm.~C), give an
$O(1/\sqrt n)$ CDF Berry--Esseen for an additive functional under a spectral gap on a Banach
\emph{algebra} of regular (BV/Lipschitz) functions ($\sigma^2>0$; the regular observable's moments are
automatic) --- with \emph{no} density floor or ellipticity --- the flagship example being the \emph{singular} Bernoulli convolution,
the exact atomic-IFS analogue of our forward filter. Filter stability moreover \emph{supplies} that gap:
a contracting place-dependent IFS upgrades its Wasserstein-$1$ contraction to the full
Lipschitz-algebra-norm spectral gap (Kloeckner--Lopes--Stadlbauer, arXiv:1412.0848, Cor.~1.2; so the
measure-side ``$W_1$ gap $\Rightarrow$ CLT only'' barrier of Komorowski--Walczuk, arXiv:1102.1842, does
\emph{not} bite), surviving the analytic tilt (Herv\'e--P\`ene). \emph{So the binding obstruction is not
a density floor.} It sits one level up, in the \emph{object}: the achievability quantity is the
lower-tail \emph{ball log-probability}
\[
  G_n(x^n):=-\log_2 P_{Y^{*n}}\!\big(B_{nD}(x^n)\big)=\jmath_n(x^n,D)-\log_2 S_n(x^n),
\]
whose leading part $\jmath_n=\imath_n-n\,h(D)$ is the clean homogeneous additive functional (annealed
over the source, $\mathrm{Var}(\jmath_n)/n\to V_{\mathrm{lossless}}$, covered by Kloeckner's theorem
verbatim), while the ball residual $-\log_2 S_n$ is \emph{(i) non-additive} --- not a Birkhoff sum of
any fixed Banach-algebra observable, so the additive-functional Berry--Esseen theorems do not reach it
(within a level set of $\jmath_n$, where $\jmath_n$ is constant, $S_n$ still varies with the
\emph{arrangement} of $x^n$ --- a decoupling probe shows within-class spread $0.02$--$0.09$ bits) ---
and \emph{(ii) schedule-inhomogeneous} (the per-step indicator $\mathbf 1[Y_i\ne x_i]$ is selected by
the schedule symbol $x_i$, and for the correlated $Y^*$ the per-schedule flip-bits are not i.i.d.). The
one inhomogeneous (sequential) CDF Berry--Esseen on the market --- Hafouta (\emph{Nonlinearity} 33,
2020; arXiv:1903.04018, Thm.~2.10 on the Thm.~2.6 RPF triplet) --- needs the covering hypothesis
$T_j^{n_0}B_j(x,\xi)=\mathcal E_{j+n_0}$ that \emph{forces} the eigenfunction floor
$\min_x|h_j^{(z)}|\ge c>0$, a density-floor surrogate the atomic filter lacks; but, as Kloeckner's
singular-case theorem shows, that is a limitation of \emph{that route}, not of the CDF Berry--Esseen
grade (the sharper sequential Edgeworth of Hafouta 2025, arXiv:2511.06414, inherits it a fortiori).
\emph{The dispersion residual is therefore the ball-probability$\,\to\,$additive-functional reduction.}
That $\mathrm{Var}(\log_2 S_n)/n\to0$ uniformly over the $\sqrt n$-shell (equivalently, that $-\log_2
S_n$ is $o_P(\sqrt n)$-dominated by an algebra-additive functional of the enlarged homogeneous chain
$(X_t,b_t,\mathrm{phase})$). Sharpening the \emph{grade} of that reduction: a CDF Berry--Esseen is
\emph{provably insufficient} for it. By Abel summation $S_n=\sum_{k\le t}q_k\,e^{\theta_0(nD-k)}$ is a
$\theta_0$-geometric average of the \emph{local} lattice masses $q_k=P_{Y^{*n}}(d_H(\cdot,x^n)=k)$, each
of order $n^{-1/2}$; a CDF bound of error $O(n^{-1/2})$ cannot resolve a mass of its own order, so
pinning $-\log_2 S_n$ to $o(1)$ needs the local masses --- a \emph{lattice local-limit theorem}, not a
CDF Berry--Esseen. Precisely, the residual is a \emph{near-mean} (Edgeworth window, saddle tilt
$s^*\to0$), density-floor-free, inhomogeneous \emph{lattice} LLT for the tilted forward-filter cocycle
$L_\alpha^{(x_n)}\cdots L_\alpha^{(x_1)}$ at the integer Hamming lattice, uniform over the $\sqrt
n$-shell --- the \emph{near-mean sibling} of the third-order $\tfrac12\log n$-prefactor lemma below (the
same density-floor-free \emph{lattice} object --- non-elliptic \emph{in the filter-side
representation}; the dual source-side representation below is elliptic --- relaxed only from
fixed-fraction to near-mean: a narrowing of the regime, \emph{not} an easier object), not an
ellipticity question.

\emph{The annealed form reduces to a citable theorem; the residual is its quenched upgrade.} Coding the
forward filter to a Gibbs--Markov subshift (the belief a H\"older function of the past), the integer
Hamming cocycle $K_n=\sum_i\mathbf 1[Y^*_i\ne x_i]$ becomes a H\"older additive functional of the
homogeneous joint chain that is \emph{aperiodic} (lattice span $1$, twisted spectral radius $<1$ on
$(0,2\pi)$ --- verified numerically to $n=520$); the Aaronson--Denker 2001 / Gou\"ezel 2010
finite-variance \emph{lattice} local-limit theorem then applies \emph{verbatim} --- Gibbs--Markov $+$
aperiodicity, with \emph{no} ellipticity, density floor, or minorization. So, modulo the (standard)
symbolic coding and the verified aperiodicity, the \emph{annealed} $\mathrm{Var}_{x\sim\mu}(\log_2
S_n)/n\to0$ holds by citation (not merely by the numerics below), which shows the ellipticity ``wall''
is a transfer-operator-route artifact, not intrinsic to the LLT grade. The \emph{sole} residual is the
\emph{quenched}, shell-uniform upgrade.

\emph{Sharpening the residual to a single bounded local phase.} An Abel factoring of the geometric
lower-tail sum gives, with $\phi=e^{\theta_0}\in(0,1)$,
$-\log_2 S_n=-\log_2 q_t-\{nD\}\log_2\phi-\log_2\!\sum_{j\ge0}(q_{t-j}/q_t)\,\phi^{\,j}$; the local
ratios are \emph{not} $\approx1$ at finite $n$ --- correctly, $q_{t-j}/q_t=\exp(\beta j/\sqrt{v_x}
-j^2/(2v_x))(1+o(1))$ with $\beta=(t-\mu_x)/\sqrt{v_x}=O(1)$ on the shell, i.e.\ geometric in $j$ at the
$O(n^{-1/2})$ rate $\beta/\sqrt{v_x}$ --- but the $\phi$-weighted sum converges to a shell-uniform
$O(1)$ factor ($\to\frac1{1-\phi}$) with vanishing shell-variance (verified: ratio model to
median rel.\ err.\ $<15\%$, $W$ within $1\%$ of $\frac1{1-\phi}$, $\mathrm{Var}_{\mathrm{shell}}(\log_2 W)$
decreasing; \texttt{probe\_7\_34\_R2\_Le\_abel\_mechanism.py}), so
$\mathrm{Var}(\log_2 S_n)=\mathrm{Var}(\log_2 q_t)+o(1)$ (verified: gap $\le0.0023$, ratio $0.97$, $n$
to $520$). Writing $q_t=r_t/\sqrt{2\pi v_q}$
with $v_q=\mathrm{Var}_q(K)=\Theta(n)$, the smooth scale part $-\tfrac12\log_2(2\pi v_q)$ has
averaged-Dirichlet energy $O(1/n)$ (single-coordinate sensitivity bounded via filter stability), so the
dispersion residual sharpens to: \emph{the bounded lattice-phase residual $r_t=q_t\sqrt{2\pi v_q}$
(numerically $\approx0.94$, shell-uniform) has $\mathrm{Var}(\log_2 r_t)=O(1)$} --- precisely the
quenched near-mean lattice-LLT regularity (shell-uniform boundedness of the ratios $q_{t-j}/q_t$ and
single-flip stability of $r_t$). On the concentration side the \emph{worst-case} geometric-mixing
bounded-difference inequality (Kontorovich--Ramanan 2008) gives only $\mathrm{Var}(\log_2 S_n)=O(n)$ ---
so $V_{\mathrm{op}}\le V_{\mathrm{lossless}}+O(1)$ is \emph{finite}, never $+\infty$ --- because the
single-coordinate influence is $\Theta(1)$ at \emph{atypical} frozen-radius ($t$ small) configurations;
but in the genuine near-mean regime ($t\asymp nD$ large) the \emph{averaged}
(conditional-variance/Efron--Stein) coordinate-influence energy is numerically \emph{bounded}, $O(1)$ and
$n$-stable, whence the heat-bath Glauber--Poincar\'e inequality on the $1$D-Gibbs source law
(Diaconis--Saloff-Coste 1996; spectral gap $\Theta(1)$ by Dobrushin/Martinelli) \emph{would} give
$\mathrm{Var}(\log_2 r_t)=O(1)$ outright --- the lone gap being the \emph{rigorous} bound on that
averaged energy, which \emph{is} the quenched near-mean lattice-LLT regularity above (this concentration
framing is an equivalent restatement, not a reduction in difficulty). (Were $\mathrm{Var}(\log_2 S_n)/n$
instead to tend to a positive constant, the converse would be loose and
$V_{\mathrm{op}}>V_{\mathrm{lossless}}$ strictly; the numerics below place it at $0$.)

\emph{Why the residual is irreducible to filter stability: it is a cancellation, not a forgetting.} The
exact algebraic identity $D_i\log_2 q_t=D_i\log_2\pi_t-D_i\log_2 M$ (writing $q_t\propto\pi_t e^{\theta_0
t}$ with $\pi_t$ the untilted shell-enumerator mass and $M(x)=\mathbb E_{Y^{*}}[e^{\theta_0 K}]$ the
tilted free energy, $D_i$ the single-flip difference) holds to machine precision, and \emph{both}
single-flip energies $\mathbb E[\sum_i(D_i\log_2\pi_t)^2]$ and $\mathbb E[\sum_i(D_i\log_2 M)^2]$ are
$\Theta(n)$ with correlation $\to1$, while their difference $\mathbb E[\sum_i(D_i\log_2 q_t)^2]=O(1)$
(verified, energy decomposition exact to $10^{-17}$). So the $O(n^{-1/2})$ single-flip smallness is a
\emph{leading-order cancellation} between two $\Theta(n)$-energy objects, of which only the additive
free-energy part $\log_2 M$ carries the filter-stability-controlled (and \emph{$\Theta(1)$, non-small})
influence; \emph{filter stability cannot supply the smallness}, which lives entirely in the cancellation
$D_i\log_2\pi_t=D_i\log_2 M+O(n^{-1/2})$ --- a rank-1 restriction of the same quenched lattice-LLT-ratio
regularity. Two standard routes provably fail to close it: triangle-inequality chaining over the shell
gives only $O(\sqrt n)$ (two shell schedules differ in $\Theta(n)$ coordinates with non-localized,
participation-ratio-$\tfrac12$ single-flip influence; verified $\mathrm{chain}/\sqrt n\to0.20$ constant,
i.e.\ $\mathrm{Var}=O(n)$), and every applicable LLT is either homogeneous (Aaronson--Denker,
Gou\"ezel) or expanding/elliptic (Dolgopyat--Sarig; Dragi\v cevi\'c--Froyland--Gonz\'alez-Tokman--Vaienti
arXiv:1705.02130). The achievability thus reduces, with everything around it proven, to the
\emph{single} new theorem isolated above --- not a conjunction of citable facts, but one quenched,
inhomogeneous, near-mean \emph{lattice} local-limit regularity (non-elliptic in the filter
representation; elliptic in the dual source-side representation below), fully separated from the
rest of the argument. The closest existing tool is the sequential lattice LLT of Hafouta
(arXiv:2510.24244, 2025), which replaces uniform ellipticity by a weaker \emph{path-ball mass}
(``physicality'') floor $\inf_{j}\inf_x P_{j}(B(b,r)\mid x)>0$; this still fails for the
\emph{atomic} forward-filter (a $2$-Dirac kernel, belief-path mass on a measure-zero tree), so the
filter-cocycle framing is genuinely walled. The natural escape --- re-encoding $K_n$ via the finite binary source chain (which trivially has the
floor) --- does \emph{not} by itself close it, for a reason \emph{below} the floor: Hafouta's lattice LLT
\emph{consumes} a bounded additive cocycle $\sum_j f_j$ ($\sup_j\|f_j\|<\infty$) and \emph{produces} a
local mass, whereas the residual $-\log_2 r_t$ \emph{is} that local mass, generated by a source-driven
weight-enumerator transfer matrix of \emph{growing} ($\Theta(i)$) dimension --- not a bounded additive
functional --- with $K_n$ not even $\sigma(X)$-measurable ($Y^*$ is auxiliary). Moreover that theorem is
deterministic-single-chain (annealed over the chain's own randomness): granting the rest, it re-derives
only the \emph{annealed} lattice LLT already cited (Aaronson--Denker/Gou\"ezel), not the \emph{quenched},
$\sqrt n$-shell-uniform upgrade that is the residual. So the wall is the quenched, shell-uniform
\emph{grade} of an additive-cocycle LLT, confirmed \emph{not} ellipticity (the source-side aperiodicity ---
Hafouta's irreducibility condition --- is verified to hold). Whether a cleverer state-space reformulation
recasts $-\log_2 r_t$ as a bounded additive cocycle of a floored chain \emph{with} quenched family-uniformity
is the residual's remaining degree of freedom. \emph{Validated} in \texttt{scripts/verify/probe\_7\_34\_r2\_jumpset.py} (single-flip
influence uniformly $O(n^{-1/2})$, no $\Theta(1)$ jumps, energy $O(1)$),
\texttt{probe\_7\_34\_r2\_filter\_decomp.py} (the $\Theta(n)$-cancellation, $\mathrm{corr}\to0.99997$),
and \texttt{probe\_7\_34\_R1\_chaining\_growth.py} (chaining gives $\Theta(\sqrt n)$, true gap $O(1)$).

\emph{A dual elliptic representation: the filter is eliminable from the quenched object.} Two exact
facts re-locate where the residual lives. \emph{(Dual identity.)} For any word $x^n$, with
$\zeta(\beta)=\frac{1-e^{\beta}}{(1+e^{\beta})(1-2D)}$, $\eta=\frac{1-\zeta}{1+\zeta}$,
$C(\beta)=\frac{(1+e^{\beta})(1+\zeta)}2$,
\[
  \mathbb E_{Y^{*}}\!\big[e^{\beta\,d_H(x^n,Y^{*n})}\big]
  \;=\;C(\beta)^{n}\,\mathbb E_{X'}\!\big[\eta(\beta)^{\,d_H(x^n,X'^n)}\big],
\]
where $X'$ is an \emph{independent copy of the source} --- three lines of Walsh algebra (the
per-coordinate identity $\sum_{y_i}(-1)^{w_iy_i}e^{\beta\mathbf 1[x_i\ne y_i]}=(1+e^\beta)u^{w_i}
(-1)^{w_ix_i}$, $u=\frac{1-e^\beta}{1+e^\beta}$, applied forward to $Y^*$ and in reverse to $X'$);
algebraic in $e^\beta$ (complex $\beta$ included), verified exact to $10^{-15}$. At the saddle
$\beta=\theta_0$: $\zeta=1$, $\eta=0$, $C=\frac1{1-D}$, and $\mathbb E_{Y^*}[e^{\theta_0K}]
=(1-D)^{-n}P_X(x^n)$ exactly --- the SLB identity re-derived source-side. \emph{(Polymer gas.)} The
quenched shell polynomial $S_x(\eta):=\mathbb E_{X'}[\eta^{d_H(x,X')}]/P_X(x)$ is the grand partition
function of a one-dimensional \emph{hard-core domain-wall polymer gas} on $x$: maximal flip-blocks $I$
with activity $\eta^{|I|}\rho^{\pm1}$ per edge bond ($\rho=\frac p{1-p}$; $+1$ stay, $-1$ switch),
blocks separated by $\ge1$ unflipped site (verified exact to $10^{-13}$), and the posterior
characteristic function is $\Phi(s;x)=(C_s/C_0)^n S_x(\eta_s)$ with $\eta_s=\eta(\theta_0+is)$,
$\eta_0=0$. \emph{Consequence for the localization:} the law on the dual side is the
\emph{strictly-positive-kernel (elliptic) source} --- the non-ellipticity of the forward filter is a
property of the filter-side \emph{representation}, not of the residual object, which accordingly should
be read as: lattice $+$ quenched shell-uniformity $+$ the $\Theta(n)$-cancellation, \emph{with
representation-dependent ellipticity}. \emph{What the dual route does not give (refuted closure):} a
contour-wide cluster expansion of $\log S_x(\eta_s)$ fails --- the Kotecky--Preiss radius of the gas is
$4$--$10\times$ \emph{smaller} than $\max_s|\eta_s|$ at every $p$ (the single-flip activity threshold
$(p/(1-p))^2$ is not the convergence radius); $S_x$ carries Fisher/Lee--Yang zeros \emph{on} the contour
near $s=\pm\pi$ for a typical fraction of words (median $\min_s|S_x(\eta_s)|\approx4\times10^{-4}$ at
$p=0.4$, $n=24$), so $\log S_x$ is singular there (the zeros sit at the benign anti-periodic endpoint
where $|\Phi(\pi;x)|\to0$ --- a span-$1$ signature --- so they break the \emph{method}, not the LLT);
and $|C_s/C_0|>1$ on $(0,\pi]$, so the deterministic factor \emph{grows} and all decay of $|\Phi|$ must
come from $S_x$ --- the same $\Theta(n)$-vs-$\Theta(n)$ cancellation, now as $C^n$-growth vs.\
partition-function zeros. The residual is unchanged in content and sharper in form: a quenched,
shell-uniform tail bound $\sup_{|s|\ge\delta_{\mathrm{KP}}}|\Phi(s;x)|\le e^{-c(p)n}$ for typical $x$,
with the KP expansion supplying only the central window $|s|\le\delta_{\mathrm{KP}}$ --- a \emph{fixed}
$\Theta(1)$ window (the first crossing of $|\eta_s|=r_{\mathrm{KP}}$; measured $0.35$--$1.01$ across the
Gray region, stabilizing $\approx0.50$ as $p\to0$), since the committed KP criterion uses the worst-case
bond weight and so holds uniformly over all words. \emph{Validated} in
\texttt{scripts/verify/lemma\_7\_34b\_dual\_identity.py} (the identity, complex $\beta$, saddle facts),
\texttt{lemma\_7\_34b\_polymer\_gas.py} (the gas representation, incl.\ pathological words), and
\texttt{remark\_7\_34b\_kp\_radius\_and\_zeros.py} (the three refutation anchors: KP-vs-contour gap,
on-contour zeros, $|C_s/C_0|>1$).

\emph{The sign is numerically settled in the achievability-closing direction.} An exact
$O(n^2)$-per-sample evaluation of $G_n$ --- via the Walsh/Krawtchouk transfer-matrix identity
$P_{Y^{*n}}(B_{nD}(x^n))=2^{-n}\sum_{j}(1-2D)^{-j}\,\mathcal K_{\le t}(j)\,H_j(x^n)$ ($\mathcal K_{\le
t}$ a partial-Krawtchouk sum, $H_j$ the weight-$j$ Walsh mass from a $2\times2$ transfer matrix in a
formal variable), validated against the brute-force $2^n$ Walsh deconvolution to $10^{-18}$ at $n\le16$
--- gives, annealed over the source in-Gray ($p\in\{0.25,0.4\}$, $D=0.9\,D_c$, blocklengths to
$n=320$, well past the integer-ball-floor $nD\ge5$), a residual variance $\mathrm{Var}(\log_2 S_n)/n$
that \emph{decays to $0$} (weighted extrapolated intercept $0\pm2\times10^{-5}$; $\mathrm{Cov}(\jmath_n,
\log_2 S_n)/n\to0$ likewise), so $\mathrm{Var}(G_n)/n\to V_{\mathrm{lossless}}$. Thus the candidate
$V_{\mathrm{op}}=V_{\mathrm{lossless}}$ is \emph{numerically certified} and the residual is
achievability-consistent (not a loose converse); the \emph{sole} remaining gap was the \emph{analytic}
proof of the ball-probability$\,\to\,$additive-functional reduction $\mathrm{Var}(\log_2 S_n)/n\to0$,
uniform over the shell --- supplied by the chain below. (This numerical evidence of the \emph{sign}
preceded and guided the analytic chain that now completes the theorem.)

\emph{The second-moment (replica) route: the analytic ingredient reduces to one explicit
finite-dimensional spectral inequality plus a routine assembly.} Averaging over the schedule
\emph{before} bounding --- a mechanism orthogonal to every word-by-word negative above --- gives the
annealed two-replica identity
\[
  \mathbb E_x\big|\Phi(s;x)\big|^2=\mathbb E\big[e^{is(K-K')}\big]
  =\Big|\tfrac{C_s}{C_0}\Big|^{2n}\,u^{T}W(s)^{\,n-1}v,
\]
with $K,K'$ two conditionally independent posterior draws given $x$: via the dual identity the triple
$(x,x',x'')$ is a \emph{source}-side chain, and the global-flip reduction leaves an explicit $4\times4$
transfer matrix $W(s)$ (verified against full $2^n$ enumeration to $9\times10^{-15}$;
$\rho(W_8)=\rho(W_4)$ throughout). Set $g(s):=|C_s/C_0|^2\rho(W(s))$. Both
committed walls are \emph{bypassed}, not contradicted: no logarithm of $S_x$ is ever taken, so the
on-contour zeros are harmless ($g$ is in fact \emph{smallest} at $s=\pi$), and the
$|C_s/C_0|^{2n}$ growth is absorbed into $g$.

\textbf{Theorem (the replica spectral inequality, sharp constant, proven).} \emph{For every $p\in(0,\tfrac12)$, every
$0<D\le D_c(p)$, and every $s\in(0,\pi]$,}
\[
  g(s)\;\le\;1-\tfrac{3}{2}\,D(1-D)\,(1-\cos s)\;<\;1 ,
\]
\emph{and the constant $\tfrac32$ is sharp: the floor ratio
$R_2(s):=[1-g(s)]/[D(1-D)(1-\cos s)]$ satisfies $R_2(s)\ge\tfrac32$ throughout
the Gray region, with infimum $\tfrac32$ approached only in the corner $s=\pi$,
$D\to D_c$, $p\to0$ (the endpoint margin at $D=D_c$ is exactly
$\tfrac18p^2+\tfrac14p^3+O(p^4)$).}
\emph{Proof.} (i) \emph{Structure.} Writing $W_4[(d_1,d_2),(e_1,e_2)]=B(d_1\!\oplus\!e_1,d_2\!\oplus\!
e_2)\,\eta^{e_1}\bar\eta^{e_2}$ with $B(f_1,f_2)=\sum_a T_{a\oplus f_1}T_{a\oplus f_2}/T_a$, direct
evaluation gives $B(0,0)=B(0,1)=B(1,0)=1$ and $B(1,1)=1+q$, $q=(1-2p)^2/(p(1-p))>0$; hence
$W_4=\mathbf 1\,m^T+q\,\Pi M$ with $m=(1,\bar\eta,\eta,|\eta|^2)^T$, $M=\mathrm{diag}(m)$, $\Pi$ the
$(d_1,d_2)\mapsto(1-d_1,1-d_2)$ reversal. (ii) \emph{Exact spectrum.} The matrix determinant lemma with
the $\Pi$-pair antidiagonal structure (both pair products $=|\eta|^2$) factorises the characteristic
polynomial \emph{exactly}:
$\det(\lambda I-W_4)=(\lambda^2-q^2|\eta|^2)\,(\lambda^2-|1+\eta|^2\lambda-\kappa^2|\eta|^2)$, with
$\kappa^2=q(q+4)=(1-2p)^2/(p^2(1-p)^2)$ (since $(1-2p)^2+4p(1-p)=1$); the Perron root is always the
second factor's $\lambda_+=\tfrac12\bigl[T+\sqrt{T^2+4\kappa^2|\eta|^2}\bigr]$, $T=|1+\eta|^2$
($\lambda_+\ge\kappa|\eta|>q|\eta|$ and $|\lambda_-|=\kappa^2|\eta|^2/\lambda_+\le\lambda_+$).
(iii) \emph{Dual closed forms.} With $w=e^{is}$ the committed definitions simplify to
$\eta_s=D(1-D)(w-1)/((1-D)^2-D^2w)$, $C_s/C_0=((1-D)^2-D^2w)/(1-2D)$, and --- the key identity ---
$(C_s/C_0)(1+\eta_s)=1-D+Dw$ \emph{exactly}. Hence, with $t=1-\cos s\in(0,2]$,
$F:=|C_s/C_0|^2|1+\eta_s|^2=|1-D+De^{is}|^2=1-2D(1-D)t$ \emph{is the Jensen floor identically} (so the
bound is automatically tight to $O(s^2)$ as $s\to0$, as that floor forces), and
$G:=4\kappa^2|\eta_s|^2|C_s/C_0|^4=8\kappa^2D^2(1-D)^2t\,[(1-2D)^2+2D^2(1-D)^2t]/(1-2D)^4\ge0$, giving
$g(s)=\tfrac12\bigl[F+\sqrt{F^2+G}\bigr]$. (iv) \emph{Reduction and the Gray bound.} Since $0<F<1$,
squaring is an equivalence: $g<1\iff G<4(1-F)=8D(1-D)t\iff\Psi_t(D):=\kappa^2D(1-D)[(1-2D)^2
+2D^2(1-D)^2t]/(1-2D)^4<1$; $\Psi_t$ is affine increasing in $t$, so the worst case is $t=2$. The map
$D\mapsto\Psi_2(D)$ is strictly increasing on $(0,\tfrac12)$ (using $\tfrac{d}{dD}\tfrac{D(1-D)}
{(1-2D)^2}=\tfrac1{(1-2D)^3}>0$), and at the Gray boundary, where $1-2D_c=\sqrt{1-2p}/(1-p)$ and
$D_c(1-D_c)=p^2/(4(1-p)^2)$,
\[
  \Psi_2(D_c)=\frac{1-2p}{4(1-p)^2}+\frac{p^4}{16(1-p)^4}<\frac14+\frac1{16}=\frac5{16}
\]
(the first term $<\tfrac14$ since $1-2p<(1-p)^2$). Hence $\Psi_t(D)<\tfrac5{16}$ on the whole closed
Gray region, and, since $\sqrt{F^2+G}<2-F\le2$ bounds the denominator by $2[(2-F)+\sqrt{F^2+G}]\le8$,
$1-g=\tfrac{4(1-F)-G}{2[(2-F)+\sqrt{F^2+G}]}\ge\tfrac{8D(1-D)t\,(1-5/16)}{8}
=\tfrac{11}{16}D(1-D)t$ --- the soft uniform floor, retained because its
$\Psi$-chain constants ($\Psi_t<\tfrac5{16}$, $\Psi_0<\tfrac14$) also feed the
curvature corollary and the central-window variance floor below. The sharp
constant takes two further steps. (v) \emph{Monotonicity in $s$ (the
discriminant certificate).} Write $a:=D(1-D)\in(0,\tfrac14)$ (so $1-4a=(1-2D)^2>0$)
and $H:=F^2+G=1+Bt+Ct^2$, with, from (iii), $B=-4a+8\kappa^2a^2/(1-4a)$ and
$C=4a^2+16\kappa^2a^4/(1-4a)^2$. From $g=\tfrac12[F+\sqrt H]$ the floor ratio is
$R_2=[(1+2at)-\sqrt H]/(2at)$, and --- using the exact identity $2H-tH'=2+Bt$ ---
\[
  \frac{dR_2}{dt}=\frac{(2+Bt)-2\sqrt H}{4at^2\sqrt H},
\]
so $R_2$ is non-increasing in $t$ iff $2\sqrt H\ge2+Bt$. The squaring is safe
one-sidedly: if $4C\ge B^2$ then $4H=4+4Bt+4Ct^2\ge(2+Bt)^2$ pointwise, i.e.\
$2\sqrt H\ge|2+Bt|\ge2+Bt$, regardless of the sign of $2+Bt$. And the
discriminant is exact: with $m:=8\kappa^2a^2/(1-4a)^2>0$,
\[
  4C-B^2=8am\bigl(1-(3+\kappa^2)\,a\bigr),\qquad\text{so}\qquad
  \frac{dR_2}{dt}\le0\ \text{ for all }t\iff a\le\frac1{3+\kappa^2}.
\]
On the Gray region $a\le a_c:=D_c(1-D_c)=p^2/(4(1-p)^2)$, and
$(3+\kappa^2)\,a_c\le1$ is, cleared of denominators, the quartic inequality
$4(1-p)^4-3p^2(1-p)^2-(1-2p)^2=p^4-10p^3+17p^2-12p+3\ge0$, which in
$u:=\tfrac12-p\in(0,\tfrac12)$ reads $(16u^4+128u^3+56u^2+32u+1)/16>0$ ---
\emph{all} coefficients positive. Hence $4C-B^2>0$ strictly and $R_2$ is
strictly decreasing in $t$ on the entire Gray region: the binding point is
$t=2$, i.e.\ $s=\pi$. (vi) \emph{The sharp endpoint.} At $t=2$, $F(\pi)=1-4a$
and $G(\pi)=16\kappa^2a^2(1-2a)^2/(1-4a)^2$; since $1-2a>0$, the single
squaring below is an equivalence:
\[
  R_2(\pi)\ge\tfrac32\iff(1-2a)\ge\sqrt{(1-4a)^2+G(\pi)}
  \iff P(a,p):=(1-3a)(1-4a)^2-4\kappa^2a(1-2a)^2\ \ge\ 0
\]
(the middle step uses $(1-2a)^2-(1-4a)^2=4a(1-3a)$). Equivalently
$\kappa^2\le f(a):=(1-3a)(1-4a)^2/[4a(1-2a)^2]$, and
\[
  \frac{d\log f}{da}=-\frac{3}{1-3a}-\frac1a-\frac{4}{(1-2a)(1-4a)}<0
  \quad\text{on }(0,\tfrac14)
\]
(each term strictly negative there), so $f$ is strictly decreasing on
$(0,\tfrac14)\supseteq(0,a_c]$ and the interior inequality reduces to the
boundary $a=a_c$, where \emph{exactly}
\[
  P(a_c,p)=\frac{p^2(1-2p)^3}{4(1-p)^8}\;>\;0\qquad(p\in(0,\tfrac12)).
\]
Hence $R_2(s)\ge R_2(\pi)\ge\tfrac32$ on the whole Gray region; sharpness is
the corner ladder $R_2(\pi)\big|_{D_c}-\tfrac32=\tfrac18p^2+\tfrac14p^3+O(p^4)
\to0$ as $p\to0$ (Remark 7.34m$''''$). \hfill$\square$

\emph{Corollary (curvature $=$ posterior variance).} $g(s)=1-\bar v\,s^2+O(s^4)$ with
$\bar v=\bar v(p,D)=D(1-D)\bigl(1-\Psi_0(D)\bigr)>0$ in closed form, which reproduces the measured
curvature range $[0.0014,0.0894]$ exactly and \emph{is} the per-site posterior variance ($O(1/n)$ gaps, halving
with $n$). The theorem was found numerically first (hostile scans: $77$ combos, corner drills to
$p=0.005$, margins positive throughout including $D=D_c$; measured prefactor
$\mathbb E_x|\Phi|^2/g(s)^n\le1$ from $n\ge50$) and then proven as above; two independent adversarial
reviews re-derived every step from the raw definitions (symbolic identity checks, a $64{,}192$-point
hostile sweep, exact-rational verification at adversarially chosen corners).
The sharp-constant steps (v)--(vi) are machine-certified separately --- the
discriminant identity $4C-B^2=8am(1-(3+\kappa^2)a)$, the positive-coefficient
boundary quartic, and the assembled $dR_2/dt<0$ in
\texttt{scripts/verify/probe\_7\_34m\_monotone\_a2\_proof.py} (T1--T5); the
endpoint reduction, the exact boundary margin $p^2(1-2p)^3/(4(1-p)^8)$, and the
$f$-monotonicity closure in \texttt{probe\_7\_34m\_endpoint\_sharp\_a2.py}
(S1--S7) --- all PASS.
With the inequality proven,
Markov plus an $e^{cn}$-net (the word-uniform bound
$|\partial_s\Phi|\le\mathbb E_Q[K\mid x]\le n$) yields the quenched uniform tail lemma
\emph{unconditionally}:
\[
  \Pr_x\Big(\sup_{\varepsilon\le|s|\le\pi}\big|\Phi(s;x)\big|>2e^{-cn}\Big)\le \pi C_1\,n\,e^{-cn},
  \qquad c=\tfrac14\ln(1/\bar g),\quad \bar g:=\max_{[\varepsilon,\pi]}g<1
\]
--- exactly the restated tail residual, for all but an exponentially small fraction of words (an
exceptional set that is \emph{intrinsic}: atypical schedules such as the all-stay word genuinely
violate the bound at large $n$, so no word-uniform version exists, and the $e^{-cn}$ bad-set Var
contribution is absorbed by the pointwise $|\log_2 S_n|=O(n)$ floor). \emph{And the central window
provably does not hide the $\Theta(n)$-cancellation:} $\Phi(0;x)=1$ \emph{identically} ($\eta_0=0$, the
empty gas), so the tilted free energy $\log_2 M(x)$ of the cancellation identity above is divided out
\emph{algebraically} in forming the posterior characteristic function --- the cancellation is
structural, not asymptotic, and never has to be re-proven. On the fixed window
$|s|\le\delta_{\mathrm{KP}}$ the $x$-dependence reduces to the cumulants $\mu_x,v_x$ --- sums of
exponentially local, $O(1)$-influence cluster terms in the i.i.d.\ bonds (Azuma) --- and quenched
steepest descent gives $q_t=(2\pi v_x)^{-1/2}e^{-\beta^2/2}\bigl(1+O(\log^3 n/\sqrt n)\bigr)
+O(e^{-cn})$ with $\beta=(t-\mu_x)/\sqrt{v_x}$ sub-Gaussian $O(1)$ over the shell
($\mathbb E_x[\mu_x]=nD$ \emph{exactly}, $K\sim\mathrm{Bin}(n,D)$ annealed; $nD(1-D)=\mathbb E[v_x]
+\mathrm{Var}(\mu_x)$ splits the total variance): the fixed-$t$ centering costs only the budgeted factor
$e^{-\beta^2/2}$ --- precisely the committed $\mathrm{Var}(\log_2 r_t)=O(1)$ allowance, now
\emph{explaining} the $O(n^{-1/2})$ single-flip influence ($D_i\log_2 q_t\approx-\beta\,D_i\mu_x/
\sqrt{v_x}$) --- whence $\mathrm{Var}(\log_2 S_n\mid\text{shell})=O(1)$ and the reduction assembles.
There is no middle zone: the KP window is fixed while $g(s)\sim1-\bar v s^2$ ($\bar v>0$) lets the replica bound
reach $s\sim n^{-1/2+\kappa}$, so the handoff overlaps with $\Theta(1)$ slack. Machine-checked
end-to-end at $n$ up to $1024$ ($150$ shell words per $n$): $\mathrm{Var}(\log_2 q_t)$ $n$-stable and
matched by the Gaussian model, steepest-descent residual decaying as $n^{-1/2}$, tail suprema within the
lemma bound (\texttt{scripts/verify/probe\_7\_34\_replica\_secondmoment\_tail.py},
\texttt{probe\_7\_34\_replica\_assembly\_central.py}; the theorem's algebra and inequality chain in
\texttt{probe\_7\_34\_G\_analytic\_bounds.py}, V0--V11b); independently, the \emph{per-word quenched}
rates $g_n(s;x)=-\tfrac1n\log|\Phi(s;x)|$ evaluated directly by an $O(n)$ two-term polymer recursion
(validated against $2^n$ brute force; $|\Phi|\le1$ checked throughout) are $n$-stable positive and sit
above the replica-implied floor $\tfrac{11}{32}D(1-D)(1-\cos s)$ at every $p$, on $nD$-matched
schedules to $nD=10$ ($n=3589$ at $p=0.1$), with the small-$p$ flat-$|\Phi|$ appearance resolved as the
frozen-radius ($nD<1$) artifact plus the $n$-independent Gaussian shoulder
(\texttt{probe\_7\_34\_quenched\_tail\_rate.py}). The quenched tail bound --- the form in which the
residual sat after the dual/polymer
analysis --- is \emph{proven}, with explicit $p$-uniform constants
($c=\tfrac14\ln(1/\bar g)$, $\bar g\le1-\tfrac{11}{16}D(1-D)(1-\cos\varepsilon)$). The
ordered-constants central-window chain \emph{(R2)} is moreover now executed and adversarially verified
step by step: the window is \emph{analytic}, $\delta_{\mathrm{KP}}\ge r_{\mathrm{KP}}(p)/A(D)>0$ with
$A(D)=\frac{D(1-D)}{1-2D}$ the exact contour slope ($|\eta_s|\le A(D)|s|$, the denominator's closest
approach to $0$ being exactly $1-2D$); the one genuinely at-risk constant is \emph{closed in closed
form}: $\mathrm{Var}_x(\mu_x)=\Theta(n)$ (so the total-variance split $nD(1-D)=\mathbb
E[v_x]+\mathrm{Var}(\mu_x)$ does \emph{not} give positivity for free), yet the per-site posterior
variance is $\bar v=\lim\mathbb E[v_x]/n=D(1-D)(1-\Psi_0(D))$ with
$\Psi_0(D)=\kappa^2D(1-D)/(1-2D)^2<\tfrac14$ on the \emph{closed} Gray region (the $t=0$ case of the
proven $\Psi_t<\tfrac5{16}$ chain), whence $\bar v\ge\tfrac34D(1-D)>0$ \emph{including at} $D=D_c$
(verified in closed form against measurement to $n=800$ across $p$, incl.\ corner $p=0.05$;
\texttt{probe\_7\_34\_R2\_Lb\_variance\_split.py}); the cumulants $(\mu_x,v_x)$ concentrate by
McDiarmid from the $O(1)$-influence cluster representation; the Abel step holds via the corrected
weighted-sum mechanism above; and the tail/window handoff tiles all of $(0,\pi]$ with a \emph{fixed}
$\varepsilon=\min\bigl(\delta_{\mathrm{KP}},\,c'\bar v/C_3\bigr)$ --- still a fixed $\Theta(1)$
threshold (no shrinking-$\varepsilon$ is ever needed, so $c>0$ is fixed; the second entry is where the
zone-2 quadratic domination $-\log|\Phi|\ge v_xs^2/4$ requires $C_3ns\le v_x/4$).
\emph{The last estimate to be written out} (GAP-1) \emph{is}: the \emph{per-word, shell-uniform}
steepest-descent remainder $R_n(x)=O(\log^3n/\sqrt n)$ in $q_t=(2\pi v_x)^{-1/2}e^{-\beta^2/2}(1+R_n)
+O(e^{-cn})$. This is of the standard saddle-point/Edgeworth \emph{type}, but we flag honestly that the
classical sources (Petrov; Bhattacharya--Rao) treat \emph{independent} summands and do not literally
contain the per-word statement --- the closest weak-dependence framework is G\"otze--Hipp 1983
(Edgeworth expansions under exponential strong mixing plus a \emph{conditional Cram\'er} condition,
\emph{Z.\ Wahrsch.\ verw.\ Geb.}\ 64:211--239; Lahiri 1993 refinements), whose conditional-Cram\'er
role is here played by the \emph{proven} tail lemma; under the proven KP analyticity (word-uniform Cauchy bound
$|\partial_s^3\log\Phi|\le C_3n$ on the half-window) plus the unconditional tail lemma it is a
technique-level adaptation --- a complex-saddle contour shift $z_0=i(\mu_x-t)/v_x$ making the cubic
error \emph{multiplicative}, with the vertical-end linear term of favorable sign --- rather than a
quotable theorem; it is now \emph{written out below} with explicit constants (and is numerically
certified \emph{tight} in the $(1+|\beta|)^3/\sqrt n$ form on adversarial fixed-bias ensembles:
$\mathrm{RMS}(R_n)=0.063\to0.035\to0.028$ at $n=256/512/1024$, consistent with $\log^3n/\sqrt n$). Scope notes (honest): the assembly holds for
$0<D\le D_c$ and $nD\ge5$ --- every constant is finite and positive on the closed region ($\bar
v(D_c)\ge\tfrac34D(1-D)$), the endpoint $D=D_c$ being covered by its dedicated verification round
(the M\"obius-orbit floor lemma below for the marginal discriminant, and the full probe suites re-run
at $D_c$; \texttt{probe\_7\_34\_dc\_endpoint.py}); the exceptional-set
absorption $O(n^3e^{-cn})\to0$ holds for each fixed $(p,D)$ but its crossover is large where the
spectral margin $c$ is small (small $D$, or $p$ near $\tfrac12$). The residual thus \emph{transforms} twice over: from a quenched, inhomogeneous,
near-mean lattice local-limit regularity --- the form in which it sat in the empty intersection of the
LLT literature --- to an explicit $4\times4$ spectral-radius inequality (now a \emph{theorem}) plus
\emph{one} per-word-uniform Edgeworth remainder, now written out below with every constant
closed and machine-checked. With both in hand the chain is complete; the closure is stated as the
concluding theorem after the write-out.

\emph{The GAP-1 write-out (two-scale architecture; adversarially verified; every numbered inequality machine-checked in \texttt{scripts/verify/probe\_7\_34\_gap1\_writeout\_v2.py} at $(p,D)\in\{(0.25,0.9D_c),(0.4,0.5D_c)\}$, $n$ to $2048$, two seeds, plus referee corner checks at $(0.1,0.5D_c)$ and $(0.45,0.9D_c)$).} \emph{Standing objects (all committed; cited, not re-derived).}
Fix $p\in(0,\tfrac12)$ and $0<D\le D_c(p)$ in the closed Gray region, and
write $\rho=\tfrac{p}{1-p}$, $\theta_0=\log\tfrac{D}{1-D}$,
$t=\lfloor nD\rfloor$ (we assume the committed floor $nD\ge5$).  For a
source word $x=x^n$ let $\Phi(s;x)=\mathbb E_Q[e^{isK}\mid x]
=(C_s/C_0)^n S_x(\eta_s)$ be the posterior characteristic function of the
Hamming distance $K$ of a posterior draw ($K$ integer, $\Phi$
$2\pi$-periodic, $|\Phi|\le1$, $\Phi(0;x)=1$), $q_t=q_t(x)=\Pr_Q(K=t\mid
x)$, and let $\mu_x=\mathbb E_Q[K\mid x]$, $v_x=\mathrm{Var}_Q(K\mid x)$,
\[
  \beta=\frac{t-\mu_x}{\sqrt{v_x}},\qquad
  y_0=\frac{\mu_x-t}{v_x}=-\frac{\beta}{\sqrt{v_x}},\qquad
  z_0=i\,y_0 .
\]
The committed inputs are:

\begin{itemize}
\item[(I1)] \emph{(KP analyticity; committed L-a/R1.)}  With
$r_{\mathrm{lb}}:=\tfrac{p^2}{12e(1-p)^2}\le r_{\mathrm{KP}}$,
$A(D):=\tfrac{D(1-D)}{1-2D}$, and
$\delta_{\mathrm c}:=\min\!\bigl(\tfrac{r_{\mathrm{lb}}}{4A(D)},\,
\tau_{\max},\,\tfrac12\bigr)$ the committed complex window
($\tau_{\max}=\min(\ln2,\ln(1+\tfrac{1-2D}{2D^2}))$; at all
verification points used here and in the referee corner checks,
$\delta_{\mathrm c}=\tfrac{r_{\mathrm{lb}}}{4A}$),
the map $z\mapsto\log\Phi(z;x)$ is analytic on the closed disk
$D_2:=\{|z|\le\delta_{\mathrm c}\}$ (complex contour slope
$|\eta_z|\le4A(D)|z|\le r_{\mathrm{lb}}$ there; $S_x$ zero-free on
$|\eta|\le r_{\mathrm{lb}}$), with the word-uniform extensivity bound
$\sup_{D_2}|\log\Phi(z;x)|\le K_\Phi\,n$,
$K_\Phi:=K_C+K_S$ ($K_C=\sup_{D_2}|\log(C_z/C_0)|$ closed-form,
$K_S=\sup_{|\eta|\le r_{\mathrm{lb}}}\sup_x|\log S_x(\eta)|/n$ from the
KP cluster expansion).
\item[(I2)] \emph{(Cauchy third-derivative bound; committed.)}  Define
the \textbf{strip scale}
\[
  \boxed{\;\varepsilon_{\mathrm{strip}}:=\tfrac12\,\delta_{\mathrm c}
  \;\Bigl(=\tfrac{r_{\mathrm{lb}}}{8A(D)}\text{ at the physical
  points}\Bigr)\;}
\]
Then for every $|z|\le\tfrac43\varepsilon_{\mathrm{strip}}$ the disk of
radius $\tfrac23\varepsilon_{\mathrm{strip}}$ about $z$ lies in $D_2$,
and Cauchy gives, word-uniformly,
\[
  \tag{I2}
  \sup_{|z|\le\frac43\varepsilon_{\mathrm{strip}}}
  \bigl|\partial_z^3\log\Phi(z;x)\bigr|\;\le\;C_3\,n,
  \qquad
  C_3:=\frac{81\,K_\Phi}{4\,\varepsilon_{\mathrm{strip}}^{3}} ,
\]
\emph{or any smaller constant for which (I2) holds} (the chain below is
monotone in $C_3$; the numerical instantiation uses the measured
admissible value $C_3^{\mathrm{meas}}$, reported by the probe).
\item[(I3)] \emph{(Taylor with exact cumulants.)}  By (I2) and Taylor
with integral remainder along $[0,z]\subset D_2$,
\[
  \tag{I3}
  \log\Phi(z;x)=i\mu_x z-\tfrac{v_x}{2}z^2+r_3(z;x),\qquad
  |r_3(z;x)|\le\tfrac{C_3 n}{6}|z|^3
  \quad\text{for }|z|\le\tfrac43\varepsilon_{\mathrm{strip}} .
\]
\item[(I4)] \emph{(Quenched uniform tail lemma; the proven theorem.)}
For any fixed $\varepsilon>0$, off a bad set of probability $\le\pi
C_1ne^{-cn}$,
$\sup_{\varepsilon\le|s|\le\pi}|\Phi(s;x)|\le2e^{-cn}$ with
$c=c(\varepsilon)=\tfrac14\ln(1/\bar g(\varepsilon))$,
$\bar g(\varepsilon)\le1-\tfrac{11}{16}D(1-D)(1-\cos\varepsilon)$.
\item[(I5)] \emph{(Replica second moment; the proven spectral
inequality, quantitative form of the tail-lemma mechanism.)}  With
$c_g:=\tfrac{11}{16}D(1-D)$,
\[
  \tag{I5}
  \mathbb E_x\bigl|\Phi(s;x)\bigr|^2\;\le\;C_1\,g(s)^{\,n-1}
  \;\le\;C_1\,e^{-c_g(1-\cos s)(n-1)}
  \;\le\;C_1\,e^{-\frac{c_g}{\pi^2}\,n s^2}\quad(n\ge2),
\]
the last step by Jordan's inequality $1-\cos s\ge\tfrac{2}{\pi^2}s^2$ on
$[0,\pi]$ and $n-1\ge n/2$.
\item[(I6)] \emph{(Word-uniform derivative.)}  $|\partial_s\Phi(s;x)|
\le\mathbb E_Q[K\mid x]\le n$.
\item[(I7)] \emph{(McDiarmid; committed L-c/R3.)}  Single-bond
influences $c_\mu=2A(D)\rho^{-1}(\rho^{-1}-\rho)$ on $\mu_x$ and
$c_v=O(1)$ on $v_x$ (cluster representation), giving
$\Pr(|\mu_x-nD|\ge\lambda)\le2e^{-2\lambda^2/((n-1)c_\mu^2)}$ and the
same for $v_x$ with $c_v$; $\mathbb E[\mu_x]=nD$ exactly.
\item[(I8)] \emph{(Variance floor; committed L-b/R2.)}
$\mathbb E[v_x]=\bar v\,n+b_0$ exactly (affine), with
$\bar v=D(1-D)(1-\Psi_0(D))\ge\tfrac34D(1-D)>0$ on the closed Gray
region.
\end{itemize}

\emph{The two scales and the derived constants.}
The architecture keeps \emph{two distinct, never-conflated scales}:
\[
  \varepsilon_2:=\min\Bigl(\varepsilon_{\mathrm{strip}},\,
  \frac{3\bar v}{32\,C_3}\Bigr)
  \qquad\text{(contour scale; $=\varepsilon_{\mathrm{strip}}$ at the
  physical points with $C_3^{\mathrm{meas}}$)},
\]
\[
  \varepsilon_{\mathrm{split}}:=\min\Bigl(\delta_{\mathrm{KP}},\,
  \frac{\bar v}{4C_3}\Bigr)
  \qquad\text{(the \emph{committed} tail handoff;
  $\delta_{\mathrm{KP}}\ge r_{\mathrm{lb}}/A(D)$)} .
\]
Note $\varepsilon_2\le\varepsilon_{\mathrm{split}}$ always
($\varepsilon_{\mathrm{strip}}\le\delta_{\mathrm{KP}}/8$ and
$\tfrac{3}{32}<\tfrac14$).  Further set
\[
  \sigma_n:=\Bigl(\frac{3}{4C_3 n}\Bigr)^{1/3},\qquad
  \tilde\sigma:=\min(\sigma_n,\varepsilon_2),\qquad
  c_2:=\frac{c_g}{4\pi^2},\qquad
  c_4:=\frac{c_g}{16\pi^2},\qquad
  c_5:=\frac{c_g}{4\pi^2},
\]
\[
  B_3:=\frac{3c_\mu}{\sqrt{\bar v}},\qquad
  c:=\tfrac14\ln\bigl(1/\bar g(\varepsilon_{\mathrm{split}})\bigr),\qquad
  J:=\bigl\lceil\log_2(\varepsilon_{\mathrm{split}}/\varepsilon_2)
  \bigr\rceil ,
\]
\[
  C_a:=\frac{8e\,C_3}{3\sqrt{2\pi}}\Bigl(\frac{2}{\bar v}\Bigr)^{3/2},
  \qquad
  C_b:=\frac{2e\,C_3}{3}\Bigl(\frac{2}{\bar v}\Bigr)^{3/2},\qquad
  \boxed{\;C:=C_a+C_b+5\;}.
\]

\emph{Good sets.}  Let
\[
  \widetilde E_1:=\underbrace{\Bigl\{\sup_{\varepsilon_{\mathrm{split}}
  \le|s|\le\pi}|\Phi(s;x)|>2e^{-cn}\Bigr\}}_{E_1\text{ (committed tail
  bad set)}}
  \;\cup\;
  \underbrace{\Bigl\{\exists\,s\in[\varepsilon_2,
  \varepsilon_{\mathrm{split}}]:\ |\Phi(s;x)|>e^{-c_4 n s^2}
  \Bigr\}}_{E_{2'}\text{ (zone-$2'$ bad set, Step 5)}},
\]
\[
  E_2:=\Bigl\{v_x<\tfrac{\bar v\,n}{2}\Bigr\},\qquad
  E_3:=\Bigl\{|t-\mu_x|>B_3\sqrt{v_x\log n}\Bigr\}
  \quad(\text{i.e.\ }|\beta|>B_3\sqrt{\log n}) .
\]
Their probabilities (Step 6):
\[
  \tag{P}
  \Pr(\widetilde E_1)\le\pi C_1ne^{-cn}
  +12\,C_1 J\,n\,\varepsilon_{\mathrm{split}}\,
  e^{-c_5 n\varepsilon_2^2},
  \qquad
  \Pr(E_2)\le e^{-\frac{\bar v^2}{8c_v^2}n}\ (n\ge n_a),
  \qquad
  \Pr(E_3\setminus E_2)\le2n^{-3}\ (n\ge n_b),
\]
with $n_a:=\lceil4|b_0|/\bar v\rceil$ and $n_b$ the explicit threshold
in Step 6.  All three are summable against the pointwise
$|\log_2 S_n|=O(n)$ floor, which the following lemma supplies.

\emph{The pointwise floor (M\"obius-orbit lemma; proven at every $0<D\le D_c$, including the marginal
discriminant).} The deconvolved output law has the exact transfer form
$P_{Y^*}(y)=2^{-n}\,u(s_1)^TM(s_2)\cdots M(s_n)\,v$ with
$M(s)=\bigl(\begin{smallmatrix}1&s\\ c\alpha s&\alpha\end{smallmatrix}\bigr)$, $\alpha=1-2p$,
$c=(1-2D)^{-2}$; the adversarial component-ratio envelope obeys $x_{i+1}=\psi(x_i)$,
$\psi(x)=\tfrac{1-\alpha x}{1-c\alpha x}$, $x_1=1$, with
$\psi(x)-x=\tfrac{c\alpha(x-x_-)(x-x_+)}{1-c\alpha x}$, the roots being the transfer-recursion fixed
points. At $D=D_c$ --- the marginal discriminant, double root $x_\pm=x^*=\tfrac1{1-p}$ and
$c\alpha=(1-p)^2$ \emph{exactly} (rational; the odd Walsh masks vanish identically, so the irrational
$1-2D_c$ never enters) --- the pole margin is \emph{strict}: $c\alpha x^*=1-p<1$. The orbit therefore
rises monotonically to $x^*$ strictly below the pole, every conditional satisfies
$P(y_{i+1}\mid y^i)\ge p/2$, and
\[
  \min_y P_{Y^*}(y)\;\ge\;2^{-n}\,p^{\,n-1}
  \qquad\text{at }D=D_c,\ \text{and a fortiori (larger fixed-point margin) for all }0<D\le D_c,
\]
asymptotically tight (the greedy envelope word is the exact arg-min; per-site factor decreasing to
$p/2$ --- the historical $60$-digit value $\approx0.0547$ at $p=0.1$ is this orbit at $n\approx70$).
Consequently $0\le G_n\le n(1+\log_2\tfrac1p)$, while $|\jmath_n|\le n\log_2\tfrac1p+1$ per word from
the displayed product law, so $|\log_2 S_n|=|G_n-\jmath_n|=O(n)$ for \emph{every}
word: the floor used above and below --- previously asserted, and on the open interior only implicitly
backed by the strict deconvolution margin --- is hereby proven on the whole \emph{closed} Gray region,
and the all-$n$ endpoint deconvolution positivity is upgraded from numerical to exact. Verified in
exact integer/rational arithmetic (exhaustive minimum over all words to $n=18$ at three $p$, in float
to $n=24$ and greedy to $n=64$; greedy arg-min match; orbit identities in rationals; the transfer form
cross-checked against the raw Hadamard deconvolution at $n=10,12$, whose all-words minimum concurs) in
\texttt{scripts/verify/probe\_7\_34\_dc\_endpoint.py} (DC2).

\bigskip
\noindent\textbf{Lemma (GAP-1: per-word, shell-uniform steepest-descent
remainder).}
\emph{Fix $p\in(0,\tfrac12)$, $0<D\le D_c(p)$.  There is an explicit
$n_0=n_0(p,D)$ --- the largest crossing point of the finitely many
explicit conditions (N0)--(N5) below together with $n_a,n_b$ and
$nD\ge5$ --- such that for every $n\ge n_0$ and every word
$x\notin\widetilde E_1\cup E_2\cup E_3$,}
\[
  q_t(x)\;=\;\frac{1}{\sqrt{2\pi v_x}}\;e^{-\beta^2/2}\,
  \bigl(1+R_n(x)\bigr),
  \qquad
  |R_n(x)|\;\le\;\frac{C\,(1+|\beta|)^3}{\sqrt n},
\]
\emph{with $C=C_a+C_b+5$ as above.  ($n_0$ is large --- the chain is
asymptotic --- but every constant is closed-form; the probe prints
$n_0$ at the physical points.  $n_0(p,D)$ is finite at every fixed
point of the closed Gray region including $D=D_c$, but is not claimed
uniform as $(p,D)$ approaches the boundary.)}

\bigskip
\noindent\emph{Proof.}

\textbf{Step 0 (inversion and contour).}  Since $K$ is integer-valued,
\[
  \tag{E0}
  2\pi\,q_t=\int_{-\pi}^{\pi}\Phi(s;x)\,e^{-ist}\,ds .
\]
On the good set, by (N0) below,
\[
  |y_0|=\frac{|\beta|}{\sqrt{v_x}}
  \le\frac{B_3\sqrt{\log n}}{\sqrt{\bar v n/2}}
  \le\min\Bigl(\tilde\sigma,\;
  \frac{\varepsilon_{\mathrm{strip}}}{2}\Bigr).
  \tag{N0}
\]
Replace the central segment $[-\varepsilon_2,\varepsilon_2]$ by the
three sides of the rectangle with vertices $\pm\varepsilon_2$,
$\pm\varepsilon_2+iy_0$.  Every point $z$ of this rectangle satisfies
$|z|\le\sqrt{\varepsilon_2^2+y_0^2}\le
\sqrt{1+\tfrac14}\,\varepsilon_{\mathrm{strip}}<
\tfrac43\varepsilon_{\mathrm{strip}}$, so the rectangle lies in the
analyticity domain of (I1)/(I3) and Cauchy's theorem gives
\[
  \tag{E1}
  2\pi q_t=\underbrace{\int_{-\varepsilon_2}^{\varepsilon_2}
  \Phi(s+iy_0)e^{-i(s+iy_0)t}ds}_{\text{central (Steps 1--3)}}
  \;+\;\underbrace{V_++V_-}_{\text{vertical ends (Step 4)}}
  \;+\;\underbrace{\int_{\varepsilon_2\le|s|\le
  \varepsilon_{\mathrm{split}}}\Phi e^{-ist}ds}_{Z_{2'}\text{ (Step 5)}}
  \;+\;\underbrace{\int_{\varepsilon_{\mathrm{split}}\le|s|\le\pi}
  \Phi e^{-ist}ds}_{Z_3\text{ (Step 5)}} .
\]

\textbf{Step 1 (exact saddle cancellation).}  For $z=s+iy_0$, by (I3)
and the \emph{algebraic identity}
\[
  i(\mu_x-t)(s+iy_0)-\tfrac{v_x}{2}(s+iy_0)^2
  =-\tfrac{\beta^2}{2}-\tfrac{v_x}{2}s^2
  \qquad\Bigl(\text{since } y_0=\tfrac{\mu_x-t}{v_x},\
  v_xy_0^2=\beta^2\Bigr),
\]
we get, \emph{exactly},
\[
  \tag{E2}
  \psi(s):=\log\Phi(s+iy_0;x)-i(s+iy_0)t
  =-\frac{\beta^2}{2}-\frac{v_x s^2}{2}+r_3(s+iy_0;x).
\]
The linear term cancels identically --- no $e^{t|y_0|}=e^{\Theta(\sqrt
n)}$ phase factor ever appears (the refuted draft's pitfall (ii)): the
modulus $|e^{-izt}|=e^{ty}$ is combined with
$\mathrm{Re}\,[i\mu_x z]=-\mu_x y$ \emph{before} any bound is taken,
leaving $(t-\mu_x)y=\Theta(\beta^2)$, not $\Theta(t\,y)$.

\textbf{Step 2 (Gaussian completion).}
\[
  \tag{E3}
  \int_{-\varepsilon_2}^{\varepsilon_2}e^{-v_xs^2/2}ds
  =\sqrt{\frac{2\pi}{v_x}}-T_G,
  \qquad
  0\le T_G\le\frac{2}{v_x\varepsilon_2}\,
  e^{-v_x\varepsilon_2^2/2}.
\]

\textbf{Step 3 (cubic error, split at $\tilde\sigma$).}  Write the
central integral as
$e^{-\beta^2/2}\bigl[\int_{-\varepsilon_2}^{\varepsilon_2}
e^{-v_xs^2/2}ds+E_{\mathrm{in}}+E_{\mathrm{out}}\bigr]$ with
$E_{\bullet}=\int_{\bullet}e^{-v_xs^2/2}(e^{r_3(s+iy_0)}-1)\,ds$ over
$|s|\le\tilde\sigma$ resp.\ $\tilde\sigma\le|s|\le\varepsilon_2$.

\emph{(3a) Inner.}  For $|s|\le\tilde\sigma$, by (N0)
$|y_0|\le\tilde\sigma\le\sigma_n$, so
$|r_3|\le\tfrac{C_3n}{6}(|s|+|y_0|)^3\le\tfrac{C_3n}{6}(2\sigma_n)^3=1$;
hence $|e^{r_3}-1|\le e\,|r_3|$ and, extending the integral to
$\mathbb R$ and using $(a+b)^3\le4(a^3+b^3)$,
$\int_{\mathbb R}|s|^3e^{-v_xs^2/2}ds=4/v_x^2$:
\[
  \tag{E5}
  |E_{\mathrm{in}}|
  \le e\,\frac{C_3n}{6}\cdot4\Bigl[\frac{4}{v_x^{2}}
  +|y_0|^3\sqrt{\frac{2\pi}{v_x}}\Bigr]
  =\frac{2e\,C_3n}{3}\Bigl[\frac{4}{v_x^{2}}
  +|y_0|^3\sqrt{\frac{2\pi}{v_x}}\Bigr].
\]
Dividing by $\sqrt{2\pi/v_x}$ and using $v_x\ge\bar vn/2$ (off $E_2$),
$|y_0|^3=|\beta|^3v_x^{-3/2}$:
\[
  \tag{E5$'$}
  \frac{|E_{\mathrm{in}}|}{\sqrt{2\pi/v_x}}
  \le\frac{C_a}{\sqrt n}+\frac{C_b|\beta|^3}{\sqrt n}.
\]

\emph{(3b) Outer (empty unless $\sigma_n<\varepsilon_2$, i.e.
$n>\tfrac{3}{4C_3\varepsilon_2^3}$).}  For
$\tilde\sigma\le|s|\le\varepsilon_2$ we have $|y_0|\le\tilde\sigma\le
|s|$, so $(|s|+|y_0|)^3\le8|s|^3$, and since
$\varepsilon_2\le\tfrac{3\bar v}{32C_3}\le\tfrac{3v_x}{16C_3n}$ (off
$E_2$),
\[
  \tag{E6}
  \mathrm{Re}\,r_3(s+iy_0)\le\frac{4C_3n}{3}|s|^3
  \le\frac{v_x s^2}{4},
  \qquad\text{hence}\quad
  e^{-v_xs^2/2}\bigl|e^{r_3}-1\bigr|\le2e^{-v_xs^2/4},
\]
\[
  \tag{E7}
  |E_{\mathrm{out}}|\le4\int_{\sigma_n}^{\infty}e^{-v_xs^2/4}ds
  \le\frac{8}{v_x\sigma_n}e^{-v_x\sigma_n^2/4},
  \qquad
  \frac{|E_{\mathrm{out}}|}{\sqrt{2\pi/v_x}}
  \le\frac{8\,e^{-\frac{\bar v}{8}n\sigma_n^2}}
  {\sigma_n\sqrt{2\pi\bar vn/2}}
  \;\overset{\text{(N2)}}{\le}\;\frac{1}{\sqrt n}.
\]
(The exponent is $\tfrac{\bar v}{8}(\tfrac{3}{4C_3})^{2/3}n^{1/3}$ ---
genuinely $e^{-\Theta(n^{1/3})}$.)

\textbf{Step 4 (vertical ends).}  For $z=\pm\varepsilon_2+iy$, $y$
between $0$ and $y_0$ ($y=\tau y_0$, $\tau\in[0,1]$), by (I3),
\[
  \mathrm{Re}\,\bigl[\log\Phi(z)-izt\bigr]
  =\underbrace{-v_x y_0 y+\tfrac{v_x}{2}y^2}_{=-\beta^2(\tau-\tau^2/2)
  \,\le\,0}
  \;-\;\frac{v_x\varepsilon_2^2}{2}
  \;+\;\mathrm{Re}\,r_3(z),
\]
the first brace being the \emph{favorable-sign} exact saddle residual
(again no $t|y_0|$ appears).  Since $|y_0|\le\varepsilon_2$ by (N0),
$\mathrm{Re}\,r_3\le\tfrac{C_3n}{6}(2\varepsilon_2)^3\le
\tfrac{v_x\varepsilon_2^2}{4}$ as in (E6), so
\[
  \tag{E8}
  \sup_{y}\bigl|\Phi(\pm\varepsilon_2+iy)\,e^{-izt}\bigr|
  \le e^{-v_x\varepsilon_2^2/4},
  \qquad
  |V_++V_-|\le2|y_0|\,e^{-v_x\varepsilon_2^2/4}.
\]

\textbf{Step 5 (the two real-line zones --- the two scales kept
distinct).}

\emph{(5a) Zone-$2'$, $[\varepsilon_2,\varepsilon_{\mathrm{split}}]$:
quenched quadratic domination by Markov + dyadic net.}  Partition
$[\varepsilon_2,\varepsilon_{\mathrm{split}}]$ into $J$ dyadic blocks
$I_k=[a_k,\min(2a_k,\varepsilon_{\mathrm{split}})]$, $a_k=2^k
\varepsilon_2$.  On $I_k$ put a net of spacing
$\delta_k=\theta_k/n$ with threshold
$\theta_k=\tfrac12e^{-c_2na_k^2}$, so $M_k\le2na_ke^{c_2na_k^2}+1$
points.  Per net point $u\ge a_k$, by Markov and (I5),
\[
  \tag{E9a}
  \Pr\bigl(|\Phi(u;x)|>\theta_k\bigr)
  \le\frac{\mathbb E|\Phi(u)|^2}{\theta_k^2}
  \le4C_1\,e^{2c_2na_k^2}e^{-\frac{c_g}{\pi^2}nu^2}
  \le4C_1\,e^{-\frac{c_g}{2\pi^2}na_k^2}
  \quad\Bigl(c_2=\frac{c_g}{4\pi^2}\Bigr),
\]
\[
  \tag{E9b}
  \Pr\Bigl(\bigcup_k\bigcup_{u\in N_k}\{|\Phi(u)|>\theta_k\}\Bigr)
  \le\sum_k M_k\cdot4C_1e^{-\frac{c_g}{2\pi^2}na_k^2}
  \le12\,C_1J\,n\,\varepsilon_{\mathrm{split}}\,
  e^{-c_5n\varepsilon_2^2}
\]
(the $k$-sum is dominated by $k=0$ since $a_k^2=4^k\varepsilon_2^2$;
the net-count exponent $c_2na_k^2$ is beaten \emph{blockwise} by the
margin $\tfrac{c_g}{\pi^2}-3c_2=\tfrac{c_g}{4\pi^2}>0$ --- this is why
the net must be dyadic-local, not global; the constant $12$ assumes
$n\varepsilon_2\ge1$ --- the binding block being $k=0$ with
$a_0=\varepsilon_2$ --- subsumed in $n_0$).  Off this event, for any
$s\in I_k$ with nearest net point $u$, by (I6)
$|\Phi(s)|\le\theta_k+n\delta_k=2\theta_k=e^{-c_2na_k^2}\le
e^{-c_2ns^2/4}$ (as $s\le2a_k$), i.e.\ the \emph{per-word} band bound
\[
  \tag{E9}
  |\Phi(s;x)|\le e^{-c_4ns^2}\qquad
  \forall\,s\in[\varepsilon_2,\varepsilon_{\mathrm{split}}],\
  x\notin E_{2'},\qquad c_4=\tfrac{c_g}{16\pi^2} .
\]
(The measured per-site rate on this band is $\approx\bar v s^2/2$ ---
two orders above this provable floor; the probe checks both.)  Hence
\[
  \tag{E10}
  |Z_{2'}|\le2\int_{\varepsilon_2}^{\infty}e^{-c_4ns^2}ds
  \le\frac{1}{c_4n\varepsilon_2}\,e^{-c_4n\varepsilon_2^2}.
\]

\emph{(5b) Zone 3, $[\varepsilon_{\mathrm{split}},\pi]$: the committed
tail lemma verbatim at $\varepsilon=\varepsilon_{\mathrm{split}}$.}
Off $E_1$,
\[
  \tag{E11}
  |Z_3|\le2(\pi-\varepsilon_{\mathrm{split}})\cdot2e^{-cn}
  \le4\pi e^{-cn}.
\]
\emph{Two-scale discipline (the refuted draft's pitfall (i)):} the tail
bound (E11) with $c=c(\varepsilon_{\mathrm{split}})$ is invoked
\emph{only} for $|s|\ge\varepsilon_{\mathrm{split}}\approx0.1$; at
$|s|=\varepsilon_2\approx0.013$--$0.025$ one has $|\Phi|\approx1$
(per-site rate $O(\varepsilon_2^2\bar v)\sim10^{-5}$) and \emph{only}
the quadratic bounds (E8)--(E9) are used there.

\textbf{Step 6 (good-set probabilities).}  $E_1$, $E_{2'}$: (I4) and
(E9b), giving (P).  $E_2$: by (I7)--(I8), for $n\ge
n_a=\lceil4|b_0|/\bar v\rceil$, $\mathbb E[v_x]-\tfrac{\bar vn}{2}\ge
\tfrac{\bar vn}{4}$, so $\Pr(E_2)\le e^{-2(\bar vn/4)^2/((n-1)c_v^2)}
\le e^{-\bar v^2n/(8c_v^2)}$.  $E_3$: off $E_2$,
$|\beta|>B_3\sqrt{\log n}$ forces $|\mu_x-nD|\ge
B_3\sqrt{\bar v/2}\sqrt{n\log n}-1$, so by (I7)
$\Pr(E_3\setminus E_2)\le2\exp\bigl(-\tfrac{2(B_3\sqrt{\bar v/2}
\sqrt{n\log n}-1)^2}{(n-1)c_\mu^2}\bigr)\le2n^{-3}$ for $n\ge n_b$:
with $n_b$ the explicit point past which
$1\le\tfrac16B_3\sqrt{\bar v/2}\sqrt{n\log n}$, the exponent is
$\ge\tfrac{25}{36}\cdot\tfrac{B_3^2\bar v}{c_\mu^2}\log n
=\tfrac{25}{36}\cdot9\log n=6.25\log n\ge3\log n$
(using $B_3=3c_\mu/\sqrt{\bar v}$ and $\tfrac{n}{n-1}\ge1$); the probe
evaluates $n_b$ numerically ($n_b\approx16$ at both points).

\textbf{Step 7 (assembly and the explicit $n_0$).}  Collecting
(E1)--(E11),
\[
  q_t=\frac{e^{-\beta^2/2}}{\sqrt{2\pi v_x}}
  \Bigl(1+R_n\Bigr),\qquad
  R_n=\frac{E_{\mathrm{in}}+E_{\mathrm{out}}-T_G}{\sqrt{2\pi/v_x}}
  +\frac{(V_++V_-)+Z_{2'}+Z_3}{\sqrt{2\pi/v_x}}\,e^{\beta^2/2},
\]
and with $v_x\ge\bar vn/2$, $v_x\le n^2/4$,
$|\beta|\le B_3\sqrt{\log n}$ (so $e^{\beta^2/2}\le n^{B_3^2/2}$), the
five exponentially small terms are each $\le n^{-1/2}$ once
\[
  \tag{N1}
  \frac{2}{\varepsilon_2\sqrt{\pi\bar v n}}\,
  e^{-\frac{\bar v}{4}n\varepsilon_2^2}\le\frac1{\sqrt n}
  \qquad(T_G),
\]
\[
  \tag{N2}
  \frac{8}{\sigma_n\sqrt{\pi\bar v n}}\,
  e^{-\frac{\bar v}{8}n\sigma_n^2}\le\frac1{\sqrt n}
  \qquad(E_{\mathrm{out}};\ \text{vacuous while }
  \sigma_n\ge\varepsilon_2),
\]
\[
  \tag{N3}
  \frac{2B_3\sqrt{\log n}}{\sqrt{2\pi}}\,n^{B_3^2/2}\,
  e^{-\frac{\bar v}{8}n\varepsilon_2^2}\le\frac1{\sqrt n}
  \qquad(V_\pm),
\]
\[
  \tag{N4}
  \frac{1}{c_4n\varepsilon_2}\cdot\frac{n}{2\sqrt{2\pi}}\,
  n^{B_3^2/2}\,e^{-c_4n\varepsilon_2^2}\le\frac1{\sqrt n}
  \qquad(Z_{2'}),
\]
\[
  \tag{N5}
  4\pi e^{-cn}\cdot\frac{n}{2\sqrt{2\pi}}\,
  n^{B_3^2/2}\le\frac1{\sqrt n}
  \qquad(Z_3).
\]
Together with (E5$'$),
\[
  |R_n|\le\frac{C_a+C_b|\beta|^3+5}{\sqrt n}
  \le\frac{(C_a+C_b+5)(1+|\beta|)^3}{\sqrt n}
  =\frac{C(1+|\beta|)^3}{\sqrt n}.
\]
$n_0(p,D):=$ the largest crossing point of (N0)--(N5), $n_a$, $n_b$,
$nD\ge5$, $n\ge2$ --- each an explicit inequality in $n$ with
closed-form constants, hence computable (the probe prints it).
\hfill$\square$

\bigskip
\emph{Remarks (deviations from the prescribed route, justified;
honest scope).}
\begin{itemize}
\item[(a)] \emph{Two scales.}  $\varepsilon_2$ (contour/strip,
$\approx0.013$--$0.025$) and $\varepsilon_{\mathrm{split}}$ (tail
handoff, $\approx0.1$) carry different symbols throughout; the proof
never evaluates the tail bound at the strip scale.  This is the exact
repair of the refuted draft's eps-conflation
(\texttt{probe\_7\_34\_gap1\_ADV\_eps\_conflation.py}).
\item[(b)] \emph{The zone-1/zone-2 split at $\sigma_n=\Theta(n^{-1/3})$
instead of the prescribed $\delta_1=A'\log n/\sqrt n$.}  With the
$\delta_1$-cut the inner error must be bounded by its worst case on the
window, $e^{O(n(\delta_1+|y_0|)^3)}$, which yields only
$|R_n|=O((\log n+|\beta|)^3/\sqrt n)$ --- the committed text's
$\log^3n/\sqrt n$ form.  The stated target $C(1+|\beta|)^3/\sqrt n$
(no polylog) \emph{requires} weighting the cubic against the Gaussian
(E5), and then the natural cut is where the cubic reaches $O(1)$,
i.e.\ $\sigma_n$.  Geometry (shift, strip, vertical ends, zones
$2'$/3) is unchanged from the prescription.  Note $\sigma_n>
\varepsilon_2$ for all $n\lesssim3/(4C_3\varepsilon_2^3)$ (measured:
$1.1\times10^7$ resp.\ $9.7\times10^5$ at the two physical points): at
accessible $n$ the outer zone is empty and (E5) covers the whole
shifted segment --- the probe verifies (E5) directly and (E7)/(N2) by
explicit arithmetic in the engaged regime.
\item[(c)] \emph{$C_3$ robustness.}  The chain is stated for any
admissible $C_3$ in (I2); with the worst-case Cauchy value
$C_3=\tfrac{81K_\Phi}{4\varepsilon_{\mathrm{strip}}^3}$ all constants
remain explicit ($\varepsilon_2$ shrinks to $\tfrac{3\bar v}{32C_3}$,
$n_0$ grows); with the measured admissible $C_3^{\mathrm{meas}}
\approx0.015$ resp.\ $0.048$ at the two physical points,
$\varepsilon_2=\varepsilon_{\mathrm{strip}}$.
\item[(d)] \emph{Honest size of $n_0$.}  Measured crossings: $n_0
\approx5\times10^8$ at $(0.4,\,0.5D_c)$ (dominated by (N4)) and
$n_0\approx9\times10^{12}$ at $(0.25,\,0.9D_c)$ (dominated by (N0),
i.e.\ by the lossy McDiarmid cap $B_3=3c_\mu/\sqrt{\bar v}\approx8.9$
there).  The lemma is asymptotic with explicit constants, exactly the
committed scope (``$n_0$ large'').
\item[(e)] \emph{What is new vs the committed text.}  Nothing is
assumed beyond (I1)--(I8), all committed; the new content is the
write-out itself: (E2)'s exact cancellation, the (E5) weighted-cubic
bound, the (E8) favorable-sign vertical bound, and the (E9) dyadic-net
quenched quadratic domination with explicit Jordan constants.
\item[(f)] \emph{Uniformity in the target count (the Abel handoff).}
The proof is verbatim uniform in $t'$ over $|t'-\mu_x|\le
B'\sqrt{v_x\log n}$ --- $t$ enters only through $\beta$ and $y_0$ ---
so the lemma applies unchanged at the shifted counts $t-j$,
$j\le K\log^2n$, of the committed Abel step with the committed
$B'=B_3+1$ bookkeeping ($|\beta_j|\le(B_3+1)\sqrt{\log n}$ there);
only the crossings (N0), (N3)--(N5) shift $B_3\to B_3+1$, changing
$n_0$ but not $C$.
\end{itemize}

\textbf{Theorem (BSMS Gray-region rate--distortion dispersion).} \emph{Fix $p\in(0,\tfrac12)$ and
$0<D\le D_c(p)$ --- the closed Gray region --- with $V_{\mathrm{op}}(D)$ the operational dispersion
under the fixed-$\varepsilon$
excess-distortion criterion (the second-order coefficient of $R(n,D,\varepsilon)$,
$\varepsilon\in(0,1)$ fixed; criterion restated in Remark 7.34d below). Then}
\[
  V_{\mathrm{op}}(D)\;=\;V_{\mathrm{lossless}}\;=\;p(1-p)\log_2^2\tfrac{1-p}p .
\]
\emph{Proof.} The converse $V_{\mathrm{op}}\ge V_{\mathrm{lossless}}$ follows from the SLB-exact
identity together with the one-shot Kostina--Verd\'u $d$-tilted converse and the varentropy CLT (the
opening of this remark). For the
achievability: by the GAP-1 lemma just proven, for $n\ge n_0(p,D)$ (which enforces $nD\ge5$) and every
word outside
$\widetilde E_1\cup E_2\cup E_3$, $-\log_2 q_t=\tfrac12\log_2(2\pi v_x)+\tfrac{\beta^2}{2}\log_2e
-\log_2(1+R_n)$ with $|R_n|\le C(1+|\beta|)^3/\sqrt n$; the Abel factoring (with remark (f)'s
uniformity in the shifted counts) converts this to $-\log_2 S_n=\tfrac12\log_2(2\pi v_x)
+\tfrac{\beta^2}{2}\log_2e-\{nD\}\log_2\phi+\log_2(1-\phi)+o(1)$ on the good set (the last two terms
deterministic per $n$); $\beta$ is sub-Gaussian $O(1)$ and
$\mathrm{Var}(\log_2 v_x)=O(1/n)$ (the (R2) chain). One threshold bookkeeping point: the
$\log_2(1+R_n)$ term is bounded \emph{proof-grade} by the worst-case majorant only once
$C(1+B_3\sqrt{\log n})^3/\sqrt n<1$, i.e.\ for $n\ge n_1(p,D)$, which exceeds the GAP-1 threshold
($n_1\approx1.2\times10^{14}$ vs $n_0\approx8.8\times10^{12}$ at $(0.25,0.9D_c)$); we therefore take
the theorem's threshold to be $n_0':=\max(n_0,n_1)$ --- immaterial to the asymptotic statement, while
the \emph{measured} $R_n$ decay (RMS $0.060\to0.017$ over $n=256$--$2048$, $n$-stable form constant)
corroborates the bound far below $n_1$. So $\mathrm{Var}(\log_2 S_n\mid\text{shell})=O(1)$,
the bad sets being absorbed by the pointwise $|\log_2 S_n|=O(n)$ floor against their exponential
($\widetilde E_1$, $E_2$) and polynomial ($E_3$: $2n^{-3}$) probabilities. Hence
$\mathrm{Var}(\log_2 S_n)/n\to0$ uniformly over the
$\sqrt n$-typical shell, $\mathrm{Var}(G_n)/n\to V_{\mathrm{lossless}}$ with
$\mathrm{Cov}(\jmath_n,\log_2 S_n)/n\to0$ (Cauchy--Schwarz from the $O(1)$ variance), and the committed
ball-probability reduction yields $V_{\mathrm{op}}\le V_{\mathrm{lossless}}$. \hfill$\square$

\emph{Scope of the theorem (honest restatement).} Asymptotic with explicit but large $n_0(p,D)$
(remark (d); at the small-$D_c$ corner $p=0.1$, $D=D_c$, the crossing reaches $\approx2.9\times
10^{17}$). The endpoint $D=D_c$ is included by a dedicated verification round: every chain statement
already covers the closed region (the replica theorem, the variance floor, the GAP-1 lemma's own
scope), the marginal-discriminant floor is supplied by the M\"obius-orbit lemma above, and the full
GAP-1 (V0--V9) and central-window suites re-run \emph{at} $D=D_c$ for $p\in\{0.1,0.25,0.4\}$ pass
identically to the interior profile (\texttt{probe\_7\_34\_dc\_endpoint.py}, DC0--DC4). The
exceptional word set is intrinsic (atypical schedules genuinely violate the
tail bound) and enters only through its measure. Beyond-BSMS (non-symmetric chains, larger alphabets)
and beyond-Gray ($D>D_c$) remain open, as does the third-order prefactor below.

\emph{The sharp third-order prefactor (a separate, also-open refinement).} Pinning the exact
$\tfrac12\log n$ term would additionally require, for the integer Hamming ball, the lower-tail
prefactor $-\log P_{Y^{*n}}(B_{nD}(x^n))=\jmath_n(x^n,D)+\tfrac12\log n+O(1)$ --- a \emph{lattice}
sharp-LD local-limit theorem for the same cocycle. A spectral (Nagaev--Guivarc'h)
analysis reduces this to a Bahadur--Rao for a \emph{filter functional} and clears two of its three
sub-obstructions: the exact form $e^{-nI}/(\lambda\sqrt{2\pi n\sigma^2})$ with $nI=\jmath_n$ holds
(density-floor-free, via the dynamical complex-RPF route detailed below), and the \emph{full}
large-deviation regime applies (the Dolgopyat--Sarig partial-regime degeneracy needs a
vanishing-probability symbol, absent here as $Y^*$ has both symbols at probability $\tfrac12$). The
residual is the combination no single theorem spans: a \emph{fixed-fraction} deviation ($D=\Theta(1)$,
beyond Dolgopyat--Sarig's small-$\varepsilon$ window) with a \emph{time-inhomogeneous} test function
(the schedule $x_i$, beyond Kontoyiannis--Meyn's single fixed $F$), \emph{uniform} over the
$\sqrt n$-typical shell. In the Gray region $Y^*$ has \emph{Hankel rank exactly $2$} (a rank-$2$
observable-operator model; verified by the in-Gray Walsh deconvolution) yet is genuinely
infinite-order with \emph{no finite positive (HMM) realization} (its observable operators carry
irremovable negative entries), so its natural Markov state is the continuous \emph{forward filter},
geometrically ergodic under filter stability (the [RNR-II, §7.15]/7.30a tool). Steps (i)--(iii), (iv-a), and (iv-c) are
constructed: the filter realization with a spectral gap (the gap \emph{substituting} for ellipticity,
the atomic filter kernel lacking a density floor); the tilted operator $L_\alpha^{(c)}$ with analytic
spectral gap (Keller--Liverani/Herv\'e--P\`ene); the single-schedule lattice Bahadur--Rao via the
\emph{dynamical} complex Ruelle--Perron--Frobenius route (Lalley 1989; Pollicott--Sharp;
Kesseb\"ohmer--Kombrink), density-floor-free on a H\"older space --- \emph{not} Kontoyiannis--Meyn
2005, whose exact lattice expansion requires their domination condition (DV3+)(ii), which they prove
cannot be removed and which the \emph{continuum} filter kernel cannot satisfy --- its moving Dirac
atoms admit no finite dominating measure (the same obstruction that rules out the elliptic routes;
a finite-state atomic chain would satisfy it via counting measure, but the filter lives on a
continuum). The contracting IFS is coded to a one-sided symbolic system (the belief a H\"older
function of the past) so the RPF hypotheses hold directly, with the lattice step a checkable
cohomological non-arithmeticity condition on the filter; the schedule-driven slope/top-Lyapunov exponent via Birkhoff--Hopf cone
contraction (Hennion 1997, stationary driving), giving the Legendre rate $nI=\jmath_n$; and the
$\sqrt n$-shell uniformity via analytic perturbation of the leading eigenvalue. Within this
third-order refinement the remaining gap is the $\tfrac12\log n+O(1)$ \emph{prefactor} of the lattice
sharp-large-deviation local-limit theorem for the schedule-driven product
$L_\alpha^{(x_n)}\cdots L_\alpha^{(x_1)}$, uniform over the shell. The natural route --- transferring the Dolgopyat--Sarig 2023 inhomogeneous
local-limit theorem with Hilbert-cone contraction in place of ellipticity --- is \emph{ruled out}:
their decorrelation (Prop.~2.13) is indeed contraction-coefficient-driven and transfers, but their
change-of-measure, bridge, and hexagon constructions use ellipticity \emph{essentially}, through
\emph{pointwise} density lower bounds $p_n(x\to z)\ge\varepsilon_0$ that a contraction coefficient does
not supply (e.g.\ $P=\big(\begin{smallmatrix}1&0\\ \tfrac12&\tfrac12\end{smallmatrix}\big)$ has
$\delta(P^2)<1$ yet $P^2(1,2)=0$), and the atomic filter kernel has no density floor, which a bounded
tilt cannot create. The \emph{nearest} existing tool is Dragi\v cevi\'c--Hafouta 2020 (Thm.~4.18),
which already proves a quenched \emph{sharp} large-deviation prefactor of order $1/\sqrt n$ for tilted
transfer-operator \emph{cocycles} (an Edgeworth/sharp-LD expansion) --- so the sharp inhomogeneous
prefactor is \emph{not} unprecedented. It does not cover our exact object, and this third-order lattice
LLT is precisely its adaptation across five axes: (1) a \emph{deterministic} source schedule (vs.\ their quenched-a.e.\
random driving); (2) a \emph{contracting atomic} IFS operator (vs.\ expanding/Anosov --- bridged by
coding the contracting filter to a one-sided symbolic system, the belief being a H\"older function of
the past, so the RPF hypotheses hold directly); (3) a \emph{lattice} (integer) tail, needing the
arithmetic/peripheral-spectrum (non-arithmeticity) condition rather than their non-lattice one;
(4) a \emph{fixed-fraction} deviation $D<D_c<\tfrac12$ in the interior of the Legendre domain (vs.\
their near-mean/small-LD regime); and (5) \emph{uniformity} over the $\sqrt n$-typical shell of
schedules. Of these, (1) adapts (Hafouta 2025's deterministic \emph{sequential} transfer-operator
theory, density-floor-free), (2) and (5) adapt (symbolic coding; analytic perturbation), and (4)
reduces \emph{onto} (3) via the Doob tilt. The hard blocker is axis (3) --- the \emph{lattice} case
for the non-elliptic inhomogeneous product --- which exposes a sharp \emph{structural dichotomy} in
the literature: every inhomogeneous lattice-or-large-deviation result (Dolgopyat--Sarig 2023;
Dolgopyat--Hafouta 2022) requires \emph{uniform ellipticity}, while every density-floor-free
spectral-gap result (Hafouta 2025) is \emph{non-lattice and near-mean} only; the target
$-\log P(A_n\le nD)=\jmath_n+\tfrac12\log n+O(1)$ for an integer-valued tail sits in their
\emph{empty intersection}, and the atomic filter kernel has no ellipticity (in the \emph{filter-side}
representation; the dual source-side representation of the dispersion-grade object above is elliptic
but trades the missing floor for the partition-function-zeros obstruction, so the intersection
statement is unchanged in substance). So the sharp third-order prefactor needs a
genuinely \emph{new} theorem --- a lattice sharp-LD local-limit expansion for inhomogeneous
non-elliptic transfer-operator cocycles --- absent from the literature; the dispersion residual
\emph{was} its
\emph{near-mean sibling} --- the \emph{same} density-floor-free \emph{lattice}
local-limit object of (a)--(b) above (representation-dependent ellipticity), relaxed only from
fixed-fraction to \emph{near-mean} ($s^*\to0$,
the Edgeworth window), not down to CDF Berry--Esseen grade (a CDF bound's $O(n^{-1/2})$ error is exactly
the order of the local mass $q_t\sim(2\pi v_q)^{-1/2}$ it must resolve; verified
$r_t=q_t\sqrt{2\pi v_q}\approx0.94$, $\mathrm{Var}(\log_2 S_n)=O(1)$) --- and that near-mean sibling is
exactly what the GAP-1 lemma above now supplies, closing the dispersion while the sharp-LD
\emph{prefactor} here stays open. So with the dispersion resolved (the concluding theorem above), what
this paragraph leaves open is the third-order $\tfrac12\log n$ refinement \emph{only}; the
discrete-Markov dispersion beyond the BSMS also remains open in general (Krishnamachari 2026). (The
resolved BSMS case obeys the
memoryless/Gauss--Markov ``lossy dispersion $=$ lossless varentropy in the SLB-tight region'' pattern,
with the plateau bounded by $D_c(p)$.)

\emph{Beyond the Gray region ($D>D_c(p)$): a separate open question.} Outside the Gray region the
SLB is no longer tight, $\jmath_n$ ceases to be additive, and the converse dispersion drops
$n$-stably to $\rho(D)V_{\mathrm{lossless}}$ with $\rho(D)<1$ (e.g.\ $\rho\approx0.61$ at $p=0.1$,
$D=0.05$). The ratio $\rho(D)$ has \emph{no elementary closed form} --- the best one-parameter fit
$((1-2D)/(1-2D_c))^a$ has a $p$-dependent exponent ($a\approx5.2,3.5,2.8$ at $p=0.1,0.25,0.4$) and
max-error $\approx0.02$--$0.08$, far above the $10^{-6}$ a closed form would need, and $\rho$ is not a
separable function of $(D,p,D_c)$. Its clean characterization is the Green--Kubo variance rate
$V_{\mathrm{conv}}(D)=\lim_n\mathrm{Var}(\jmath_n)/n=\chi''(0)$ (second cumulant-rate) of the
$\lambda^*$-tilted joint (source, optimal-output) transfer operator: \emph{elementary in Gray}
(order-$1$ output, $\rho\equiv1$, the $V_{\mathrm{lossless}}$ plateau) but, off Gray, the variance rate
of an \emph{infinite-dimensional} operator, because the $n$-letter optimal output requires \emph{growing}
Markov order (Eswaran--Gastpar 2022) --- \emph{the same growing-order obstruction that walls the in-Gray
achievability}, now seen to govern the off-Gray converse-ratio's non-closed-form as well. There the
operational dispersion is unresolved: the converse gives
$V_{\mathrm{op}}\ge\lim_n\mathrm{Var}(\jmath_n)/n$, the first order remains achievable (order-$k$
Markov codebooks, Elshafiy--Rose 2022), and $\jmath_n$ is numerically $\approx91$--$99\%$ the i.i.d.\
switch-count, but the matching second-order achievability needs a \emph{quenched, time-inhomogeneous}
Bahadur--Rao / lossy-AEP lemma --- the discrete-Hamming analogue of the shell-probability step
Tian--Kostina close for Gauss--Markov only via Karhunen--Lo\`eve decorrelation, with no discrete
counterpart (the optimal-output ball measure is degenerate, and a fixed-order codebook is
first-order rate-suboptimal). We leave the $D>D_c$ dispersion open; it does not affect the Gray-region
resolution above. \emph{Validated} in
\texttt{scripts/verify/remark\_7\_34b\_bsms\_rd\_dispersion\_numeric.py} (PART A: reproduces the
proven memoryless Kostina--Verd\'u plateau exactly; PART C: the all-$n$ deconvolution is valid iff
$D\le D_c(p)$, verified across $p=0.1,0.25,0.4$, establishing $V_{\mathrm{op}}=V_{\mathrm{lossless}}$ on
the Gray region; PART B: the $D>D_c$ converse ratio $\rho(D)<1$, $n$-stable),
\texttt{remark\_7\_34c\_bsms\_achievability\_transfermatrix.py} (the $D>D_c$ regime), and
\texttt{probe\_7\_34\_sign\_var\_logS.py} (the exact $O(n^2)$ Walsh/Krawtchouk evaluation of $G_n$,
brute-force-validated to $10^{-18}$ at $n\le16$, certifying $\mathrm{Var}(\log_2 S_n)/n\to0$ in-Gray ---
the achievability-consistent sign) and \texttt{probe\_7\_34\_saddlepoint\_logS\_llt.py} (the saddlepoint
reduction $-\log_2 S_n=\tfrac12\log_2(2\pi n\sigma_x^2)+\Psi(\{nD\})+o(1)$: near-mean saddle $s^*\approx0$,
boundary LLT ratio $r_t\approx0.94$, $\mathrm{Var}(\log_2 S_n)=O(1)$ --- locating the residual as a
near-mean lattice LLT, not a CDF Berry--Esseen), and \texttt{probe\_7\_34\_offgray\_rho.py} (the off-Gray
$\rho(D)$: $n$-stable, $\rho=1$ at $D_c$ and $\to0$ as $D\to\tfrac12$; no elementary closed form; the
Green--Kubo $\chi''(0)$ characterization with the growing-Markov-order obstruction off Gray).

\textbf{Remark 7.34d (The general converse: $V_{\mathrm{op}}\ge V_{\mathrm{lossless}}$ on the all-$n$
SLB-tight region of \emph{any} finite-alphabet source --- unifying the memoryless, symmetric-BSMS, and
non-symmetric-Markov cases).} The converse of Remark 7.34b is not special to the \emph{symmetric}
chain: its mechanism is \emph{source-agnostic}, and isolating that yields one principle covering a
whole type lattice of sources --- with the symmetric BSMS and the i.i.d.\ plateau as two instances and
the non-symmetric binary Markov chain as a new one.

\emph{The principle.} Let $X$ be a stationary finite-alphabet source under Hamming distortion such that
(i) its information density $\imath_n:=-\log_2 P(X^n)$ obeys a central limit theorem with
$\mathrm{Var}(\imath_n)/n\to V_{\mathrm{lossless}}$ (the source entropy-rate varentropy), and (ii) the
finite-blocklength redundancy is negligible at the dispersion scale,
$\mathbb E[\jmath_n]-nR(D)=o(\sqrt n)$; assume $V_{\mathrm{lossless}}\in(0,\infty)$ (otherwise the bound
is vacuous). Every irreducible aperiodic finite-state --- in particular finite-order Markov --- source
meets (i)--(ii): the CLT by the Markov varentropy theorem for additive functionals (Nagaev;
Meyn--Tweedie), and (ii) because $\mathbb E[\jmath_n]=H(X^n)-n\,h(D)=nR(D)+O(1)$ via $H(X^n)=n\bar h+O(1)$
($V_{\mathrm{lossless}}>0$ there unless the log-kernel is a coboundary --- the degenerate
uniform-entropy chains, excluded by the standing hypothesis). Call $D$ \emph{all-$n$
SLB-tight} if the source law deconvolves by the maximizing i.i.d.\ noise --- in the worked binary case
the $\mathrm{BSC}(D)$ kernel, $\widehat P_Y(w)=\widehat P_X(w)(1-2D)^{-|w|}$ in the Walsh domain; for a
$q$-ary alphabet the $q$-ary symmetric kernel with noise entropy $h(D)+D\log(q-1)$ in place of $h(D)$,
the displayed Walsh formula being the $2$-element-character-group specialisation --- into a valid output
law at \emph{every} blocklength, the region $0<D\le D_c$ of Remark 7.34b's mechanism. Let
$V_{\mathrm{op}}(D)$ denote the \emph{fixed-$\varepsilon$ excess-distortion dispersion}, i.e.\ the
second-order coefficient in $R(n,D,\varepsilon)=R(D)+\sqrt{V_{\mathrm{op}}(D)/n}\,Q^{-1}(\varepsilon)+o(1/\sqrt n)$
at a fixed $\varepsilon\in(0,\tfrac12)$ under the excess-distortion criterion
$\Pr[d_H(X^n,\hat X^n)>nD]\le\varepsilon$ (Kostina--Verd\'u 2012). \emph{Claim:} on the all-$n$
SLB-tight region this dispersion obeys
\[
  V_{\mathrm{op}}(D)\;\ge\;V_{\mathrm{lossless}}.
\]

\emph{Proof (converse only; identical machinery to Remark 7.34b).} On the all-$n$ SLB-tight region the
Shannon lower bound is exactly tight at every $n$, so with $\lambda_n^*=n\log_2\frac{1-D}D$ the additive
Hamming kernel forces the binary SLB identity
\[
  \jmath_n(x^n,D)=\imath_n(x^n)-n\,h(D)+c_n,
\]
where $c_n$ is a \emph{deterministic} ($x^n$-independent) finite-blocklength boundary constant with
$c_n/n\to0$ (numerically $c_n=0$ to machine precision against the $\imath_n$ baseline --- per-word
standard deviation $\lesssim10^{-8}$, falling to $\sim10^{-15}$ deep in-Gray, a
Blahut--Arimoto/near-singular-deconvolution numeric, the analytic shift being exact; a nonzero $O(1)$
value appears only if one instead normalises against the asymptotic rate $n(\bar h-h(D))$, $c_n$ then
being the entropy boundary term $H(X_1)-\bar h$).
For a $q$-ary alphabet the displayed constant becomes $n[h(D)+D\log(q-1)]$. The point is that \emph{only
the determinism of the shift is load-bearing}, so its exact value and alphabet dependence are immaterial:
$\mathrm{Var}(\jmath_n)=\mathrm{Var}(\imath_n)$ \emph{exactly}, whence
$\mathrm{Var}(\jmath_n)/n\to V_{\mathrm{lossless}}$ and $\jmath_n$ inherits the source's varentropy CLT.
Now invoke the one-shot d-tilted converse of Kostina--Verd\'u 2012 (Thm.~7,
$\varepsilon\ge\sup_{\gamma\ge0}\{\Pr[\jmath_n\ge\log M+\gamma]-e^{-\gamma}\}$), valid for an
\emph{arbitrary} source under the excess-distortion criterion. Fix $\varepsilon\in(0,\tfrac12)$ and a
slack $\delta>0$, and suppose some code sequence achieved rate
$\log M_n\le\mathbb E[\jmath_n]+\sqrt{nV_{\mathrm{lossless}}}\,(Q^{-1}(\varepsilon)-\delta)-\gamma_n$
with $\gamma_n\to\infty$, $\gamma_n=o(\sqrt n)$. Then
$\log M_n+\gamma_n\le t_n:=\mathbb E[\jmath_n]+\sqrt{nV_{\mathrm{lossless}}}\,(Q^{-1}(\varepsilon)-\delta)$,
so $\Pr[\jmath_n\ge\log M_n+\gamma_n]\ge\Pr[\jmath_n\ge t_n]\to Q(Q^{-1}(\varepsilon)-\delta)>\varepsilon$
(the CLT for $\jmath_n$ at the fixed interior point $Q^{-1}(\varepsilon)-\delta$) while $e^{-\gamma_n}\to0$,
so the converse forces $\varepsilon\ge Q(Q^{-1}(\varepsilon)-\delta)-o(1)>\varepsilon$
for large $n$ --- a contradiction. Hence every $\varepsilon$-code obeys
$\log M_n\ge\mathbb E[\jmath_n]+\sqrt{nV_{\mathrm{lossless}}}\,(Q^{-1}(\varepsilon)-\delta)-o(\sqrt n)$;
by hypothesis (ii) $\mathbb E[\jmath_n]=nR(D)+o(\sqrt n)$, so this is
$R(n,D,\varepsilon)\ge R(D)+\sqrt{V_{\mathrm{lossless}}/n}\,(Q^{-1}(\varepsilon)-\delta)+o(1/\sqrt n)$,
and letting $\delta\downarrow0$ yields
$V_{\mathrm{op}}(D)\ge\lim_n\mathrm{Var}(\jmath_n)/n=V_{\mathrm{lossless}}$ (a \emph{plain} CLT suffices
for this fixed-$\varepsilon$ liminf converse via the $\gamma_n\!\to\!\infty$ plus $\delta$-slack
argument; the $\varepsilon\to0$ rate-dispersion limit of Kostina--Verd\'u Def.~7 would additionally need
uniform/moderate-deviation tail control, not invoked here). \emph{The crucial point} is that nothing used
the symmetry of the chain or even that it is binary: only (i) an all-$n$ deconvolution region (via the
appropriate balanced kernel), making $\jmath_n$ a deterministic affine image of $\imath_n$, and (ii) a
varentropy CLT for $\imath_n$, which every finite-state source has. $\square$

\emph{The instances (a type lattice).} \emph{(0) Memoryless} ($X$ i.i.d., \emph{any} finite alphabet
under Hamming, via the $q$-ary balanced kernel): the all-$n$ region is the entire single-letter
SLB-tight range, $V_{\mathrm{lossless}}$ is the single-letter varentropy, and here the matching
\emph{achievability} also holds (Kostina 2017), so $V_{\mathrm{op}}=V_{\mathrm{lossless}}$ --- both
bounds (this is the instance that exercises the non-binary alphabet). \emph{(1) Symmetric BSMS}: Remark 7.34b, region $(0,D_c(p)]$,
$V_{\mathrm{lossless}}=p(1-p)\log_2^2\tfrac{1-p}p$. \emph{(2) Non-symmetric binary Markov} (new; switch
probabilities $P(0\!\to\!1)=a\ne b=P(1\!\to\!0)$, stationary $\pi_0=b/(a+b)$): the same Walsh
deconvolution has a \emph{positive} all-$n$ threshold $D_c(a,b)>0$ --- e.g.\ $D_c(0.1,0.3)\approx0.0129$
and $D_c(0.2,0.4)\approx0.0455$ --- on which $\jmath_n=\imath_n-n\,h(D)$ holds (the converse ratio
$\rho(D)=\mathrm{Var}(\jmath_n)/\mathrm{Var}(\imath_n)=1$ to $12$ digits at $n=12$), so
$V_{\mathrm{op}}\ge V_{\mathrm{lossless}}(a,b)$ there; outside it $\rho<1$ ($\approx0.73$--$0.88$),
exactly as for the symmetric chain. Here $\imath_n$ is no longer affine in a single i.i.d.\ switch-count
(two distinct transition logs), but it remains an additive Markov functional, so the varentropy CLT ---
hence the converse --- still applies; this is precisely the generalisation the abstract principle buys.

\emph{Scope (honest).} This is a \emph{converse} statement; the matching achievability
$V_{\mathrm{op}}\le V_{\mathrm{lossless}}$ is \emph{closed for the symmetric BSMS} --- instance (1) ---
throughout the closed Gray
region (Remark 7.34b's concluding theorem) but remains open for the non-symmetric memory instance
(2), with the residual now \emph{sharply mapped} by a dedicated re-instantiation probe. What
generalizes at \emph{proof} grade (both sides of the dual identity are degree-$\le n$ polynomials in
$e^\beta$ over $\mathbb Q(D)$, so exact-rational equality at $>n+1$ points proves it identically ---
all complex $\beta$, no floating point): the dual identity itself is \emph{chain-agnostic} (its
derivation uses only the binary character group and the validity of the BSC$(D)$ deconvolution, which
defines the region, not the identity); the $\mathrm{BSC}(D)$ backward channel survives asymmetry
\emph{structurally} (Hamming is a difference distortion on $(\mathbb Z/2)^n$, so the tilt
$e^{-\lambda d}$ is symbol-independent; KKT marginal consistency collapses to the saddle identity
$M(\theta_0;x)=(1-D)^{-n}P_X(x)$, certified exactly, whence $\jmath_n=\imath_n-n\,h(D)$ and
$\lambda^*=\ln\tfrac{1-D}{D}$ unchanged --- independently confirmed by random-start Blahut--Arimoto);
the annealed facts $K\sim\mathrm{Bin}(n,D)$, $\mathbb E[\mu_x]=nD$, and the structural cancellation
$\Phi(0;x)=1$ are chain-agnostic. The polymer gas survives with \emph{richer} activities (block
interiors no longer cancel: site-value-dependent ratios join the bond weights; the worst-case-bond KP
criterion still applies with a smaller radius). What breaks is the \emph{closed-form} replica
inequality: for $a\ne b$ the global-flip reduction fails, the $\mathbf 1m^T+q\Pi M$ structure and the
$\tfrac{11}{16}$ chain are lost, and the replica matrix stays $8\times8$ --- yet numerically
$g(s)=|C_s/C_0|^2\rho(W_8(s))<1$ holds uniformly on $[10^{-3},\pi]$ at every tested corner (slow/fast
mixing, extreme asymmetry, near-symmetric; $g\sim1-\bar vs^2$ with the curvature still equal to the
per-site posterior variance), so the instance-(2) achievability residual is exactly: (i) a proof (or
per-$(a,b,D)$ interval certification --- mechanical, though the slow-mixing small-$D_c$ corner has
curvature $\sim10^{-4}$ and needs care) of $\sup_{[\varepsilon,\pi]}g<1$ for the asymmetric
$8\times8$, now discharged per-$(a,b,D)$ --- on all of $(0,\pi]$, floored, with an explicit margin
constant --- by Lemma 7.34e below; (ii) the routine Markov bounded-differences replacement of the i.i.d.\
switch-bond McDiarmid step --- now supplied: by Paulin (2015, \emph{EJP} 20(79), Thm.~2.1; Marton
coupling) with one-step Dobrushin coefficient $\delta=|1-a-b|=|\lambda_2|$, McDiarmid holds with the
factor $1/(1-\delta)^2=1/(a+b)^2$, verified sub-Gaussian for $\mu_x$ at the test corners; and (iii) the
all-$n$ characterization of the asymmetric deconvolution threshold $D_c(a,b)$ --- now \emph{solved in
closed form} by Lemma 7.34f below: $D_c(a,b)=\tfrac12\bigl(1-\sqrt{1-4ab/(2-a-b)^2}\bigr)$ (the earlier
finite-$n$ measurements $0.0075$--$0.044$, ``still creeping at $n=14$,'' are explained exactly: above
$D_c$ the binding \emph{alternating} word first goes negative at $n\approx3/\sqrt{D/D_c-1}$, the
complex-rotation angle of the period-2 transfer product).
Everything else in the chain --- the tail-lemma Markov$+$net assembly, the KP window analyticity and
Cauchy bounds, the GAP-1 saddle machinery --- consumes only these inputs (re-instantiation checks,
exact-rational or machine-precision throughout, in
\texttt{scripts/verify/probe\_7\_34\_nonsym\_dual\_identity.py}). The general
discrete-Markov/Hamming dispersion beyond the binary chain remains open (Krishnamachari 2026). What is new is the \emph{unification}: a single source-agnostic lower bound that
collapses to Kostina 2017 in the memoryless case and extends the 7.34b converse to every finite-state
source possessing an all-$n$ SLB-tight region, with the non-symmetric chain a worked second memory
instance. \emph{Validated} in
\texttt{scripts/verify/remark\_7\_34d\_general\_converse\_discrete\_markov.py} (C1 the non-symmetric
Gray threshold $D_c(a,b)>0$; C2 in-Gray $\rho=\mathrm{Var}(\jmath_n)/\mathrm{Var}(\imath_n)=1\Rightarrow V_{\mathrm{conv}}=V_{\mathrm{lossless}}$;
C3 out-of-Gray $\rho<1$; C4 the shift $c_n=\jmath_n-(\imath_n-n\,h(D))$ is deterministic in-Gray
(per-word std $\lesssim10^{-8}$, so $\mathrm{Var}(\jmath_n)=\mathrm{Var}(\imath_n)$ exactly); C5 the
centering $\kappa_n=\mathbb E[\jmath_n]-nR(D)=H(X^n)-n\bar h$ is $O(1)$ and $n$-stable (hypothesis (ii):
$\kappa=0.2392,0.1134$ at $n=10,12,14$, exactly constant); across $(a,b)\in\{(0.1,0.3),(0.2,0.4)\}$).

\textbf{Lemma 7.34e (instance (2): the asymmetric replica spectral
inequality --- exact block-CP structure, the generalised curvature
identity, and a one-shot exact certification on all of $(0,\pi]$).}
\emph{Let $T=\bigl(\begin{smallmatrix}1-a&a\\ b&1-b\end{smallmatrix}\bigr)$
with $a,b\in(0,1)$, $a+b<1$, $a\ne b$, let $0<D<\tfrac12$, and let
$W_8(s)$, $\eta_s$, $C_s$, $g(s)=|C_s/C_0|^2\rho(W_8(s))$ be the committed
asymmetric replica objects (the dual identity being chain-agnostic at proof
grade). Write $x\in\{0,1\}$ for the outer (source) index and $(x',x'')$ for
the replica pair, $\bar d:=D(1-D)$, $c:=1-2D$.}

\emph{\textbf{(i) (Block-Kraus form and the variational certificate.)} With
$M_y:=T\,\mathrm{diag}_y(\eta_s)$ --- $\mathrm{diag}_0=\mathrm{diag}(1,\eta_s)$,
$\mathrm{diag}_1=\mathrm{diag}(\eta_s,1)$, i.e.\
$M_y[d,e]=T(d,e)\,\eta_s^{\,y\oplus e}$ --- the replica matrix has the exact
block form}
\[
  W_8\bigl[(x,\cdot),(y,\cdot)\bigr]\;=\;\frac1{T(x,y)}\;M_y\otimes\overline{M_y},
\]
\emph{so its column action is the completely positive map on pairs of
$2\times2$ matrices $\Psi_s(Z)_x=\sum_y M_y Z_y M_y^\dagger/T(x,y)$.
Consequently, for \emph{every} Hermitian positive-definite pair
$Z=(Z_0,Z_1)$ and every $\nu>0$:}
\[
  \Psi_s(Z)_x\prec\nu\,Z_x\ (x=0,1)
  \;\Longrightarrow\;\rho\bigl(W_8(s)\bigr)<\nu .
\]
\emph{\textbf{(ii) (The curvature identity generalises in closed form.)}
$W_8$ is multilinear in $(\eta,\bar\eta)$, and $W_8(0)$ has rank two with
nonzero spectrum $\mathrm{spec}(T)=\{1,\,1-a-b\}$; the Perron eigenvalue $1$
is simple, with right eigenvector $r[(x,x',x'')]=\sum_yT(x',y)T(x'',y)/T(x,y)$
(as a pair: $Z_x(0)=\sum_y v_yv_y^{T}/T(x,y)$, $v_y:=T[\,\cdot\,,y]$, carrying
\emph{unit} --- not stationary --- weights) and left functional
$l=\pi_0e_{(000)}+\pi_1e_{(111)}$, $lr=1$. Analytic perturbation at this
simple eigenvalue, with the first/second-order coefficients evaluated
\emph{exactly} symbolically over $\mathbb Q(a,b)$ ($\alpha_1=1$, $\alpha_2=0$,
$\alpha_{11}$ in closed form), gives}
\[
  g(s)=1-\bar v\,s^2+O(s^4),\qquad
  \bar v(a,b,D)=D(1-D)\Bigl[\,1-\kappa_{\mathrm{eff}}^2(a,b)\,
  \frac{D(1-D)}{(1-2D)^2}\Bigr],\qquad
  \kappa_{\mathrm{eff}}^2:=\alpha_{11}-1,
\]
\emph{where $\kappa_{\mathrm{eff}}^2(a,b)$ is an explicit $D$-independent
rational function (denominator $ab(1-a)^2(1-b)^2(a+b)$, degree-$7$ numerator
symmetric in $(a,b)$; e.g.\ $\kappa_{\mathrm{eff}}^2(0.1,0.3)=1080/49$,
$\kappa_{\mathrm{eff}}^2(0.2,0.4)=130/27$) with
$\kappa_{\mathrm{eff}}^2(p,p)=(1-2p)^2/(p^2(1-p)^2)=\kappa^2$ --- the
committed symmetric curvature recovered exactly --- and the curvature again
equals the per-site posterior variance (Richardson-extrapolated brute force
agrees to $10^{-14}$). In particular $g<1$ on a punctured neighbourhood of
$s=0$ whenever $\bar\Psi_0(a,b,D):=\kappa_{\mathrm{eff}}^2\,D(1-D)/(1-2D)^2<1$,
an exact rational test; $\bar\Psi_0\le0.385$ across the scanned asymmetric
Gray grid (maximum at $(a,b)=(0.05,0.45)$, $D=D_c$).}

\emph{\textbf{(iii) (One-shot exact certification on all of $(0,\pi]$.)} Let
$(a,b,D)$ and $\tilde c\in[0,2)$ be rational, and define the \emph{s-dressed
pair} $Z(s):=r+\eta_s\,c_{10}+\bar\eta_s\,c_{01}$ (reshaped to a pair of
$2\times2$ matrices), where $c_{10}=SW_{10}r$, $c_{01}=SW_{01}r$ are the
exact first-order perturbation corrections at the rank-two object ($S$ the
reduced resolvent, $W_{10},W_{01}$ the multilinear parts) --- all data
rational over $\mathbb Q(a,b)$, and $Z(s)$ Hermitian identically. Under
$t=\tan(s/2)$, $e^{is}=\frac{(1-t^2)+2it}{1+t^2}$, every entry of $Z$,
$\Psi_s(Z)$ and $|C_s/C_0|^2$ is rational in $t$ over $\mathbb Q(i)$ with
explicit positive denominators; clearing them turns the family of matrix
inequalities}
\[
  Z(s)\succ0\quad\text{and}\quad
  \Psi_s\bigl(Z(s)\bigr)_x\;\prec\;
  \frac{1-\tilde c\,D(1-D)(1-\cos s)}{|C_s/C_0|^2}\;Z(s)_x\qquad(x=0,1)
\]
\emph{into the positivity on $t\in(0,\infty)$ of \emph{eight explicit
polynomials} over $\mathbb Q$ (the traces and determinants of the cleared
$2\times2$ Hermitian forms; degrees $\le24$), each decided exactly by a
Sturm root count after factoring out the $t=0$ root of orders $t^2$
(traces) and $t^4$ (determinants) --- precisely the curvature order, since
$Z(0)$ is the exact Perron pair and the dressing is first-order exact ---
the $s=\pi$ endpoint being checked directly in exact rational arithmetic.
By (i), success of this single finite computation \emph{proves}}
\[
  g(s)\;<\;1-\tilde c\,D(1-D)(1-\cos s)\qquad\text{for all }s\in(0,\pi]
\]
\emph{at that $(a,b,D)$ --- the exact per-point analogue of the symmetric
$\tfrac{11}{16}$ theorem. Executed (zero floating point): at the four test
corners $\times\ D\in\{0.9D_c,D_c\}$ the certificates succeed at
$\tilde c=\lfloor64(1-\bar\Psi_0)\rfloor/64$ on the first attempt every
time --- $\tilde c\in[\tfrac{43}{64},\tfrac{25}{32}]\approx[0.67,0.78]$,
seven of the eight values above the symmetric $\tfrac{11}{16}=0.6875$, the
one below being the slow-mixing corner $(0.05,0.3)$ at $D=D_c$ exactly ---
while the off-Gray control $(0.1,0.2,D=0.03)$, where $g>1$, is correctly
refused.}

\emph{Proof.} (i) The block identity is term-by-term algebra: the committed
entry $W_8[(x,x',x''),(y,y',y'')]=T(x',y')T(x'',y'')\,
\eta^{y\oplus y'}\bar\eta^{\,y\oplus y''}/T(x,y)$ equals
$M_y[x',y']\,\overline{M_y[x'',y'']}/T(x,y)$ by the definition of $M_y$.
$\Psi_s$ is completely positive (Kraus form, operators
$M_y/\sqrt{T(x,y)}$ per component). For the certificate, let $Z$ be a
Hermitian PD pair with $\Psi_s(Z)_x\prec\nu Z_x$ and set
$\Gamma_Z(X)_x:=Z_x^{1/2}X_xZ_x^{1/2}$, an invertible congruence of
$M_2\oplus M_2$; then $\Phi:=\Gamma_Z^{-1}\circ\Psi_s\circ\Gamma_Z$ is
completely positive and similar to $\Psi_s$ as a linear operator, so
$\rho(W_8)=\rho(\Psi_s)=\rho(\Phi)\le\|\Phi\|$, and a positive map on a
unital $C^*$-algebra attains its norm at the identity (Russo--Dye; see
Paulsen, \emph{Completely Bounded Maps and Operator Algebras}, 2002,
Cor.~2.9): $\|\Phi\|=\|\Phi(I)\|=\max_x\lambda_{\max}\bigl(Z_x^{-1/2}
\Psi_s(Z)_xZ_x^{-1/2}\bigr)<\nu$. (The bound is tight: the Perron
eigenvector of the positive map lies in the PSD cone by Kre\u{\i}n--Rutman,
is strictly PD at every tested parameter, and the power-iterated pair
tightens geometrically at the subdominant ratio.)

(ii) Multilinearity is read off the entries ($D$ enters only through
$\eta_s$). At $\eta=0$ only the columns $(y,y,y)$ survive, the reduced
$2\times2$ weight-evolution matrix is $T$ itself (row-stochastic, hence the
\emph{unit} weights of $r$; the stationary law belongs to the left
functional $l$), giving rank two and the displayed spectrum; $1-a-b\in(0,1)$
makes the eigenvalue $1$ simple, so Kato analytic perturbation applies to
the multilinear family. Rayleigh--Schr\"odinger at $(e_1,e_2):=(\eta,\bar\eta)$
gives $\lambda_{\mathrm{top}}=1+\alpha_1(e_1+e_2)+\alpha_2(e_1^2+e_2^2)
+\alpha_{11}e_1e_2+O(|e|^3)$ with $\alpha_1=lW_{10}r=lW_{01}r$,
$\alpha_2=lW_{10}SW_{10}r$, $\alpha_{11}=lW_{11}r+lW_{10}SW_{01}r
+lW_{01}SW_{10}r$; the exact symbolic evaluation over $\mathbb Q(a,b)$
yields $\alpha_1=1$ and $\alpha_2=0$, so on the contour $e_2=\bar e_1$:
$\lambda_{\mathrm{top}}=|1+\eta_s|^2+\kappa_{\mathrm{eff}}^2|\eta_s|^2
+O(|\eta_s|^3)$. The two chain-agnostic exact identities
$(C_s/C_0)(1+\eta_s)=1-D+De^{is}$ and $(C_s/C_0)\,\eta_s=\bar d\,(e^{is}-1)/c$
give $|C_s/C_0|^2|1+\eta_s|^2=1-2\bar d\,(1-\cos s)$ and
$|C_s/C_0|^2|\eta_s|^2=2\bar d^{\,2}(1-\cos s)/c^2$, whence
$g(s)=1-2\bar d(1-\cos s)\bigl[1-\kappa_{\mathrm{eff}}^2\bar d/c^2\bigr]
+O(s^3)$; $g$ is even in $s$ (complex conjugation of the contour), so the
remainder is $O(s^4)$ and the displayed $\bar v$ follows. The symmetric
specialisation is an exact symbolic identity.

(iii) On $t\in(0,\infty)$ one has $\eta_s=p(t)/q(t)$ with
$p=2\bar d\,t(i-t)$, $q=(A-B)+(A+B)t^2-2iBt$ ($A=(1-D)^2$, $B=D^2$);
$|q|^2>0$ on $\mathbb R$ since the root $e^{is}=A/B>1$ of $A-Be^{is}$ is off
the unit circle; $|C_s/C_0|^{-2}=c^2(1+t^2)/\bigl(c^2(1+t^2)+4\bar
d^{\,2}t^2\bigr)$, and $1-\cos s=2t^2/(1+t^2)$. Multiplying the two matrix
inequalities by the positive quantities $|q|^2$, $|q|^4$, the
$|C|^{-2}$-denominator and $(1+t^2)$ clears all denominators; Hermiticity
of the cleared pair $\widehat Z$ is an exact symbolic identity (the swap
symmetry $c_{01}=\mathrm{swap}(c_{10})$), so positive-definiteness of each
cleared $2\times2$ Hermitian form is equivalent to positivity of its trace
and determinant --- real polynomials in $t$ over $\mathbb Q$. A polynomial
positive at $0$ (after factoring out an exact $t=0$ root) with Sturm count
zero on $[0,\infty)$ is positive there by the intermediate value theorem;
the multiplicities $2$ and $4$ are forced because
$\Psi_0(Z(0))=Z(0)$ exactly and the dressing removes the $O(s)$ error, so
the certified margin opens at the curvature order $\bar vs^2$ --- and
indeed the certified $\tilde c\approx1-\bar\Psi_0=\bar v/\bar d$ is
one half of the small-$s$ optimum $2\bar v/\bar d$, leaving room. The
$s=\pi$ point ($t=\infty$) is a single exact rational $2\times2$ PSD check
($\eta_\pi=-2\bar d/(A+B)$ real). Success at $(a,b,D,\tilde c)$ therefore
proves, via (i) with $\nu=\nu_{\tilde c}(s)/|C_s/C_0|^2$, that
$g(s)<\nu_{\tilde c}(s)=1-\tilde c\,\bar d\,(1-\cos s)$ for every
$s\in(0,\pi]$. $\square$

\emph{Two negative structural facts make the per-point certificate the
right grade here, not a stopgap: the characteristic polynomial of $W_8$ ---
multilinear in $(\eta,\bar\eta)$ over $\mathbb Q(a,b)$ --- is
\emph{irreducible} over $\mathbb Q(a,b)$ (a single degree-$8$ factor,
fully symbolic; at $a=b$ it factors and the committed flip-even quartic
divides it), so the symmetric quadratic factorisation is genuinely lost for
$a\ne b$; and the naive transplant of the symmetric Perron root with
$\kappa\mapsto\kappa_{\mathrm{eff}}$, although it matches
$\lambda_{\mathrm{top}}$ through $O(|\eta|^2)$, is \emph{refuted} as a
global bound (violated near $s=\pi$ at every tested asymmetric corner).
Combined with the quenched tail lemma --- whose Markov$+$net assembly is
chain-agnostic and consumes only $\bar g:=\max_{[\varepsilon,\pi]}g<1$ ---
each certified $(a,b,D)$ inherits the explicit constants
$\bar g\le1-\tilde c\,D(1-D)(1-\cos\varepsilon)$,
$c=\tfrac14\ln(1/\bar g)$, exactly as in the symmetric pipeline. Residual
item (i) of instance (2) therefore now reads: \emph{proven near $s=0$ in
closed form (the curvature identity), and proven on all of $(0,\pi]$ per
rational $(a,b,D)$ modulo one finite exact computation --- scheme specified,
margins explicit, executed at the test corners}. What remains open is only
the \emph{uniform-in-$(a,b,D)$} statement, and with the exact threshold of
Lemma 7.34f the situation is now \emph{decided}: the Route-A uniform Gray
bound \emph{fails} --- $\bar\Psi_0(a,b,D_c(a,b))\to1$ as $a+b\to1$ (the
earlier ``scan maximum $0.385$'' was a grid artifact that never approached
the $a+b\to1$ boundary; verified $0.794$ at $a+b=0.95$, $0.993$ at $0.999$,
fixed $a=0.05$) --- so the uniform statement holds only on \emph{bounded-mixing}
subfamilies $\{a+b\le1-\eta\}$ (e.g.\ $\sup_{a+b\le0.97}\bar\Psi_0=0.4375<1$),
or per-$(a,b)$ in $D$ via the CAD route (two-variable $(t,D)$ positivity of
the same eight polynomials). Item (ii) (Paulin bounded differences, factor
$1/(a+b)^2$) and item (iii) ($D_c(a,b)$, closed form) are discharged above
and in Lemma 7.34f. Validated in
\texttt{scripts/verify/probe\_7\_34\_asym\_replica\_structure.py}
(A0--A11: the anchored block-Kraus identity to $9\times10^{-16}$; CP/cone
preservation and Collatz--Wielandt validity on random PD pairs; the
symbolic irreducibility, $\alpha_1=1$, $\alpha_2=0$, and the symmetric
reduction of $\kappa_{\mathrm{eff}}^2$; curvature against brute-force
posterior variance to $10^{-14}$; the dressed-pair coverage map with floor
ratio $\ge1.09$ at every corner; the eight Sturm certificates and the
off-Gray refusal).}

\textbf{Remark 7.34e$'$ (the binding endpoint $s=\pi$ and the residual
convexity conjecture).} The per-point certificate of (iii) has a sharp
structural reduction. Numerically (50~dps, $4000+$ parameter points, no
exception) the ratio $R(a,b,D,s):=[1-g(s)]/[D(1-D)(1-\cos s)]$ is
\emph{minimised at the endpoint $s=\pi$}; equivalently, in $u=\cos s$, the
map $G(u):=g(\arccos u)$ is \emph{convex} on $[-1,1]$ with
$G(1)=\lim_{s\to0}g(s)=1$, so the secant slope
$[G(1)-G(u)]/(1-u)=R$ is nondecreasing and attains its minimum
$R(a,b,D,\pi)$ at $u=-1$. The endpoint is exactly solvable: at $s=\pi$ the
contour data is real and rational,
$\eta_\pi=-2D(1-D)/\bigl((1-D)^2+D^2\bigr)$ and
$|C_\pi/C_0|^2=(c^2+4D^2(1-D)^2)/c^2$ ($c=1-2D$), and $g(\pi)<1$ holds on the
open Gray interior by a closed-form-in-$D$ certificate: since $\eta_\pi<0$
makes $W_8$ \emph{signed}, the Kre\u{\i}n--Rutman/Collatz--Wielandt PSD-cone
bound (\emph{not} the M-matrix minor test, which is invalid for a signed
matrix) applies, and the once-iterated dressed pair $Z_1=\Psi_\pi(Z_{\mathrm
{dr}})$ satisfies $\Psi_\pi(Z_1)_x\preceq\mu\,Z_{1,x}$ with
$\mu=c^2/(c^2+4D^2(1-D)^2)$ throughout $D\in(0,D_c]$ --- one exact
Sturm-in-$D$ per rational $(a,b)$ (executed at representative corners; a
single CP-map iteration is load-bearing, the bare pair over-estimating
$\rho$ in a thin sliver below $D_c$). Hence
$\tilde c(a,b,D):=R(a,b,D,\pi)=[1-g(\pi)]/(2D(1-D))>0$ is explicit, though an
algebraic number of degree $6$ over $\mathbb Q(a,b)$ (the $s=\pi$ Perron
value sits in the irreducible degree-$6$ symmetric-pair block of the split
$8=2{+}6$ characteristic polynomial --- no radical form for $a\ne b$).
\emph{Were the convexity of $G$ proven}, item (i) would collapse to this
single endpoint constant, $g(s)\le1-\tilde c(a,b,D)\,D(1-D)(1-\cos s)$ on
$(0,\pi]$, eliminating the $s$-Sturm; the convexity is then the sole
residual. The obstruction is \emph{not} merely that $W_8(s)$ is complex: even
after the real, single-valued reduction below ($G$ is an algebraic, branch-free
function of $u$), the classical Perron-root convexity theorems still do not
apply, because $u$ enters through \emph{off-diagonal}, \emph{non-log-convex}
M\"obius entries. Cohen (1981), Kato (1982) and Nussbaum (1986) give $\rho$
convex (resp.\ the spectral bound convex) only when the parameter enters as a
\emph{diagonal} (multiplication) perturbation --- the off-diagonal case can
fail (Cohen's counterexample) --- and Kingman (1961) gives log-convexity only
under entrywise \emph{log-convex} dependence, which the M\"obius entries are
not. The generic non-self-adjoint variational handles point the wrong way: the
Collatz--Wielandt formula $\rho=\inf\{c:\Psi\varphi\le c\varphi,\varphi>0\}$
(Donsker--Varadhan; Anantharam--Borkar 2019) is an \emph{infimum} of
affine-in-$\Psi$ functionals, hence \emph{concave} in an affine parameter, and
the Rayleigh sup-form yields convexity only for self-adjoint families (\S below:
$M(u)$ is self-adjoint only w.r.t.\ an \emph{indefinite} fixed metric). A
single fixed pair likewise cannot stand in for the $s$-dressed one (the
frozen-$Z(\pi)$ Collatz--Wielandt slack is $\approx$constant in $s$ and
overshoots away from $s=\pi$). The only interior-or-endpoint point with real
$\eta$ besides $s=0$ is $s=\pi$.

\emph{Where the obstruction really lies (a real-representation
sharpening).} The complexity of $W_8(s)$ is, however, \emph{not} the genuine
barrier. Writing $u:=\cos s$ and acting $\Psi_s$ on a real Pauli basis
$\{I,\sigma_x,\sigma_y,\sigma_z\}$ of each $2\times2$ block gives an explicit
\emph{real} $8\times8$ representation $M(u)$ with $\rho(M(u))=\rho(W_8(s))$.
On the contour one has the exact closed forms (proven symbolically; common
linear denominator $\mathrm{DEN}(u):=(A^2+B^2)-2AB\,u>0$, $A=(1-D)^2$,
$B=D^2$): $\operatorname{Re}\eta_s$, $|\eta_s|^2$ and $|C_s/C_0|^2$ are
\emph{rational} in $u$, while
$(\operatorname{Im}\eta_s)^2=\bar d^{\,2}c^2(1-u^2)/\mathrm{DEN}(u)^2$ ---
the lone $\sqrt{1-u^2}$ enters $M(u)=P(u)+\sqrt{1-u^2}\,Q(u)$ \emph{linearly},
and $Q$ couples \emph{only} the $\sigma_y$ channel to the rest. Eliminating
that channel by a Schur complement replaces $\sqrt{1-u^2}$ by $1-u^2$, so the
characteristic polynomial of $M(u)$ is a \emph{polynomial} in $(\lambda,u)$
over $\mathbb Q(a,b,D)$ (equivalently: the char poly of $W_8$ is symmetric in
$\eta_s\leftrightarrow\bar\eta_s$, hence a polynomial in
$(\eta_s+\bar\eta_s,\,\eta_s\bar\eta_s)=(2\operatorname{Re}\eta_s,|\eta_s|^2)$,
both rational in $u$). Consequently the Perron branch
$G(u)=|C_s/C_0|^2\,\lambda(u)$ is an \emph{algebraic} function of $u$,
real-analytic on $[-1,1]$ with \emph{no} branch point at $u=\pm1$ (verified by
single-valued analytic continuation). Thus ``$G$ convex on $[-1,1]$'' is a
\emph{decidable} polynomial-positivity ($\mathrm{d}^2/\mathrm{d}u^2$ of an
algebraic function $\ge0$, one CAD/Sturm question per rational $(a,b,D)$), not
a transcendental statement --- consistent with, and explaining, the per-point
certificate of (iii). The residual barrier to a \emph{closed-form} proof is
purely \emph{quantitative}: by Rellich perturbation
$\lambda''(u)=\langle l|M''|r\rangle+2\langle l|M'\,S\,M'|r\rangle$ is a small
positive residual of a near-cancellation between the (concave) $M''$ curvature
$\langle l|M''|r\rangle<0$ (the $w''=-(1-u^2)^{-3/2}$ from the $\sqrt{}$) and
the (positive) second-order level-repulsion $2\langle l|M'SM'|r\rangle>0$,
with $(|t_1|+|t_2|)/|t_1+t_2|$ up to $\sim40$ on the test grid --- no
term-wise sign certificate. The repulsion term is itself \emph{not} a sum of
squares: although $M(u)$ is self-adjoint with respect to a single
$u$-\emph{independent} bilinear form $H$ (a one-dimensional nullspace of the
stacked commutator constraint, residual $\sim10^{-57}$ at 60~dps), that $H$ is
\emph{indefinite}, so the standard Rayleigh--Schr\"odinger expansion
$\sum_{k}|\langle l|M'|r_k\rangle|^2/(\rho-\lambda_k)$ has mixed-sign per-mode
contributions (4 of 7 negative at a representative point) and is positive only
on \emph{net}. The naive variational route is likewise refuted: for a
\emph{fixed} pair $Z$ the Collatz--Wielandt value
$\max_x\lambda_{\max}(Z_x^{-1/2}\Psi_s(Z)_xZ_x^{-1/2})$ is itself
\emph{non}-convex in $u$, and the cone Perron root is the \emph{inf} envelope
of these, not a sup of convex functions. The convexity margin is moreover
genuinely tight, degenerating as $D^2$ when $D\to0$ and vanishing at $u\to1$
(the curvature point), so any closed-form certificate must capture an
$O(D^2)$-small positive quantity at the endpoint. Validated in
\texttt{scripts/verify/lemma\_7\_34e\_spi\_closed\_form.py} (the closed-form
$\eta_\pi,|C_\pi/C_0|^2$; the exact $D$-Sturm for $g(\pi)<1$ on $(0,D_c]$;
the signed-cone correction and the once-iterated fix),
\texttt{probe\_7\_34e\_endpoint\_reduction.py} (the $s=\pi$ minimisation, the
$G$-convexity to 50~dps, and the refuted one-pair closer), and
\texttt{probe\_7\_34e\_trackA\_real\_rep\_convexity.py} (the real $8\times8$
representation; the rational reductions; the $\sqrt{}$-elimination /
algebraicity of $G(u)$; the refuted variational route; the
curvature-vs-repulsion near-cancellation), and
\texttt{probe\_7\_34\_trackB\_perturbation.py} (the exact
Rayleigh--Schr\"odinger $\rho''(u)$ decomposition; the $u$-independent but
indefinite self-adjointness metric; the mixed-sign per-mode contributions).

\textbf{Lemma 7.34f (instance (2): the asymmetric all-$n$ deconvolution
threshold in closed form).} \emph{Let $X$ be the stationary binary Markov
chain with $P(0\!\to\!1)=a$, $P(1\!\to\!0)=b$, $a,b\in(0,1)$, $a+b<1$. The
all-$n$ $\mathrm{BSC}(D)$-deconvolution region of $X$ (the asymmetric Gray
region) is $0<D\le D_c(a,b)$ with}
\[
  D_c(a,b)\;=\;\frac12\Bigl(1-\sqrt{1-\frac{4ab}{(2-a-b)^2}}\Bigr),
\]
\emph{which reduces to Gray's $D_c(p)=\tfrac12(1-\sqrt{1-2p}/(1-p))$ at
$a=b=p$ and to the i.i.d.\ Bernoulli threshold $\min(\pi_1,\pi_0)$ at
$a+b=1$.}

\emph{Proof.} Combining the two Walsh steps coordinate-wise, the deconvolved
law is the exact $2\times2$ transfer contraction
$P_{Y^*}(y)=\pi^T\bigl[\prod_{t<n}G_{y_t}T\bigr]G_{y_n}\mathbf 1$ with
$G_0=\mathrm{diag}\bigl(\tfrac{1+\zeta}2,\tfrac{1-\zeta}2\bigr)$,
$G_1$ its swap, $\zeta=(1-2D)^{-1}$ (verified exactly against the $2^n$
Walsh sum). \emph{(Binding word.)} The minimizing word is the
\emph{alternating} word $0101\cdots$ at every tested $(a,b)$ (exhaustive over
all $2^n$ words, $n\le16$): below threshold all words are nonnegative, and
the first word to go negative above it is alternating. \emph{(Discriminant.)}
The alternating word's value is governed by the period-2 product
$M=B_1B_0$, $B_y=G_yT$; its positivity for \emph{all} $n$ holds iff the
eigenvalues of $M$ are real, i.e.\ iff
$\mathrm{disc}(M)=\mathrm{tr}(M)^2-4\det(M)\ge0$, which factors exactly
(symbolic, sympy) as
\[
  16\,\mathrm{disc}(M)=\underbrace{\bigl[(a+b)^2-(a-b)^2\zeta^2\bigr]}_{F_1}
  \cdot\underbrace{\bigl[(2-a-b)^2-\bigl((2-a-b)^2-4ab\bigr)\zeta^2\bigr]}_{F_2},
\]
both factors equal to $4ab>0$ at $\zeta=1$ ($D=0$) and decreasing in
$\zeta^2$; the \emph{second} factor binds: its root
$\zeta_2^2=(2-a-b)^2/\bigl((2-a-b)^2-4ab\bigr)$ precedes the first's
$\zeta_1^2=(a+b)^2/(a-b)^2$ by the cross-difference identity
\[
  (a+b)^2\bigl[(2-a-b)^2-4ab\bigr]-(2-a-b)^2(a-b)^2\;=\;16\,ab\,(1-a-b)\;>\;0
\]
on the whole domain ($ab>0$, $a+b<1$; at $a=b$ the first factor is the
constant $(2p)^2>0$ and never vanishes), so the first sign change of
$\mathrm{disc}$ is at $\zeta^2=\zeta_2^2$, i.e.\
$(1-2D_c)^2=1-4ab/(2-a-b)^2$ --- the displayed closed form
(at the root the eigenvalues turn complex and $P_{Y^*}(\mathrm{alt})\sim
r^n\cos(n\theta+\varphi)$ oscillates negative; just above threshold the
first negative index is $n\approx3/\sqrt{D/D_c-1}$, matching the observed
finite-$n$ ``creep''). That no \emph{other} word binds at a smaller $D$:
constant words are excluded \emph{analytically} ---
$\det B_y=\det G_y\det T=\tfrac{1-\zeta^2}4(1-a-b)<0$ for every $D<\tfrac12$
(as $\zeta>1$), so $B_y$ has real eigenvalues of opposite sign at \emph{all}
$D$ and the constant-word value never oscillates --- while the alternating
product $M=B_1B_0$ has $\det M=\det B_0\det B_1>0$, the unique period-$2$
pattern whose eigenvalues \emph{can} turn complex; for the remaining mixed
words the binding role of the alternating pattern is verified exhaustively
(all $2^n$ words, $n\le16$, every tested $(a,b)$) and two-sided at
$n\le51$, $40$~dps (positive at $D_c(1-10^{-6})$ for all $n$; first
negativity at the predicted index for $D>D_c$); the word-uniform
invariant-cone argument for the mixed case is supplied separately in
Lemma 7.34g below (a two-state sign automaton, proven as a certificate
criterion and shown sharp; its existence below $D_c$ holds for all
$(a,b)$ with $a+b<1$ simultaneously, by the closed-form chain
\eqref{eq:G1chain} of Lemma 7.34h --- no per-rational computation is
needed). The radicand's
numerator $(2-a-b)^2-4ab=4-4(a+b)+(a-b)^2$ (the $\zeta^2$-coefficient of
the binding factor $F_2$) recovers Gray's discriminant $4(1-2p)$ at
$a=b$. \hfill$\square$

\emph{Validated} in
\texttt{scripts/verify/probe\_7\_34\_dc\_ab\_characterization.py} (A1 the
transfer identity; A2 the binding-word exhaustive scan; A3 the symbolic
discriminant and closed form; A4 disc-root $=$ the $n\to\infty$ threshold;
A5 the Gray reduction; A6 the Route-A re-evaluation with exact thresholds;
(ii) the Paulin concentration), plus the two-sided $n\le51$ flip check and
the first-negativity scaling $n\approx3/\sqrt{\mathrm{ovr}}$ ($97/31/14/7$
at $\mathrm{ovr}=10^{-3}/10^{-2}/0.05/0.2$, three chains).

\textbf{Lemma 7.34g (the mixed-word leaf: a sign-automaton certificate, and
its sharpness).} \emph{Notation as in Lemma 7.34f. Write the backward
transfer recursion $v_n=B_{y_n}\mathbf1$, $v_k=B_{y_k}v_{k+1}$, so that
$P_{Y^*}(y)=\pi^Tv_1$; track the slope $s_k=v_{k,2}/v_{k,1}$ and sign
$\tau_k=\operatorname{sign}v_{k,1}$. Then $B_0,B_1$ act on the slope by the
M\"obius maps $m_0(s)=\nu\,\psi(s)$, $m_1(s)=\psi(s)/\nu$, where
$\nu=(1-\zeta)/(1+\zeta)\in(-1,0)$ and $\psi(s)=\bigl(b+(1-b)s\bigr)/
\bigl((1-a)+as\bigr)$, while $(B_yv)_1=g_y\cdot\tfrac{1+\zeta}2$-type
factor times $\bigl((1-a)+as\bigr)v_1$ with $g_0=+1$, $g_1=-1$; the pole of
$\psi$ sits at $p_T=-(1-a)/a$ and the final functional $\pi^Tv_1$ changes
sign at $p_\pi=-b/a>p_T$.}

\emph{\textbf{(I) Certificate criterion.} If there are intervals
$A\subset(p_\pi,0)$ and $B\subset(-\infty,p_T)$ with}
\[
  m_0(A)\subseteq A,\;\; m_0(B)\subseteq A,\;\; m_1(A)\subseteq B,\;\;
  m_1(B)\subseteq B,\;\; \nu\in A,\;\; \nu^{-1}\in B,
\]
\emph{then $P_{Y^*}(y)>0$ for every word $y\in\{0,1\}^n$ and every $n$. (Backward
induction: a state with $y_k=0$ stays in $\{+\}\times A$, one with $y_k=1$ in
$\{-\}\times B$; the $1$-symbol's sign flip $g_1=-1$ is compensated by landing
below $p_T$, and $\tau_1=+,s_1>p_\pi$ or $\tau_1=-,s_1<p_\pi$ both give
$\pi^Tv_1>0$.) \textbf{(II) Downward closure.} If $P_{Y^*}\ge0$ for all words at
$D$, the same holds at every $D'\in(0,D]$ (BSC divisibility $K_\delta
K_{D'}=K_D$, $\delta=(D-D')/(1-2D')\in(0,\tfrac12)$, applied to the nonnegative
deconvolved law). \textbf{(III) Sharpness.} No such $A,B$ (open, closed, bounded, or unbounded)
exist for $D>D_c$: $m_1\circ m_0$ would carry $B$, hence its closure
$\overline B$ --- a \emph{proper} closed arc of the projective line
$\mathbb{RP}^1$, since $\overline B\subseteq[-\infty,p_T]$ misses the arc
$(p_T,0]\supseteq A$ --- into itself (the composition is a genuine M\"obius
homeomorphism of $\mathbb{RP}^1$, as $\det B_1B_0>0$, and $m_0(B)\subseteq A$
avoids the pole $p_T$ of $m_1$). A real M\"obius map taking a proper arc of
$\mathbb{RP}^1$ into itself cannot be elliptic --- an elliptic map is a
fixed-point-free rotation, which maps any proper arc to an arc of equal length
and hence never into a sub-arc of itself --- so it is parabolic or hyperbolic,
with a real fixed point, i.e.\ a real eigenvector of $B_1B_0$. Thus
$\operatorname{disc}(B_1B_0)=F_1F_2/16\ge0$, and (as $F_1>0$) $F_2\ge0$, i.e.\
$D\le D_c$. Above $D_c$ the composite is exactly elliptic (complex-conjugate
eigenvalue pair, $\det>0$) and the alternating word itself goes negative
(Lemma 7.34f).
\textbf{(IV) Existence below $D_c$ (uniform in $(a,b)$; closes G1).} For
$1<\zeta<\zeta_c$ the canonical certificate exists for \emph{all} $(a,b)$ with
$a+b<1$ and \emph{all} $\zeta\in(1,\zeta_c)$ at once. Take $a_1$ to be the
\emph{larger} real root of the fixed-point quadratic $\mathcal Q(s):=
A_\zeta s^2+B_\zeta s+C_\zeta$ of the slope map $m_0\circ m_1$ of $B_0B_1$,
where
\[
  A_\zeta=-\frac{a\,W}{\zeta-1},\quad
  C_\zeta=-\frac{b\,W}{\zeta+1},\quad
  B_\zeta=\frac{(a-b)(a+b-2)(\zeta^2-1)-4ab\,\zeta}{\zeta^2-1},\quad
  W:=(a-b)\zeta+(2-a-b),
\]
and set $a_2=m_0(a_1)$, $b_2=m_1(a_1)$, $b_1=m_1(b_2)$, $A=[a_1,a_2]$,
$B=[b_1,b_2]$. The entire certificate is then equivalent to the single strict
chain of real inequalities
\begin{equation}\label{eq:G1chain}
  b_1<\nu^{-1}<b_2<p_T<p_\pi<a_1<\nu<a_2<0,
\tag{$\ast$}
\end{equation}
each link of which is established in closed form, uniformly in $(a,b,\zeta)$, by
Lemma 7.34h below. The all-$n$ deconvolution region is therefore exactly
$(0,D_c(a,b)]$ on the open interior, unconditionally.}

\emph{Provenance and scope.} \emph{The symmetric case $a=b$ is Gray's 1969/1970
theorem} (R.\,M.\ Gray, \emph{Information rates of autoregressive sources}, USC-EE
Report 369, 1969, Ch.\ IV; \emph{Information rates of autoregressive processes},
IEEE Trans.\ IT-16(4):412--421, 1970): his proof is exactly an invariant cone
anchored at the eigenvectors of the alternating composite $GM_0$ (his
Fig.\ 4.3, eqs.\ (4.25)--(4.28)), with sharpness by complex-eigenvalue rotation
when the discriminant goes negative --- the present mechanism at $a=b$, where his
worked example $\alpha=3/8\Rightarrow D_c=\tfrac1{10}$ matches
$D_c(\tfrac38,\tfrac38)=\tfrac1{10}$ and his $e^{-\rho}=D/(1-D)=-\nu$. Gray's
reduction to a single generator plus the swap involution uses $M_1=GM_0G$,
which is precisely the $a=b$ flip symmetry; the asymmetric $a\ne b$ case
requires the genuine two-generator (sign-automaton / path-complete) certificate
above, which is new (cf.\ Forni--Jungers--Sepulchre, \emph{Path-complete
positivity of switching systems}, IFAC 2017, for the cone-family formalism, and
Avila--Bochi--Yoccoz, Comment.\ Math.\ Helv.\ 85(4), 2010, in whose language
$D_c$ is where the $\{B_0,B_1\}$-cocycle leaves the uniformly-hyperbolic
multicone locus through a parabolic period-2 product). No published asymmetric
extension exists (Gray 1971 gives only existence of a positive $D_c$ for
finite-state sources; Avram--Berger 1985 an underbound). \emph{Gap G1 now
closed:} the uniform-in-$(a,b)$ existence is proven in closed form in
Lemma 7.34h --- the asymmetric two-generator replacement for Gray's flip
$M_1=GM_0G$ is to anchor the certificate at the larger root of the
$B_0B_1$ slope-quadratic and collapse all eight closure inclusions to the
single chain \eqref{eq:G1chain}, four links being M\"obius identities and the
remaining four each a one-line sign reduction. No per-rational $(a,b)$
computation is needed.

\emph{Validated} (independent cold reproduction) in
\texttt{scripts/verify/lemma\_7\_34g\_sign\_automaton.py} (V0 transfer form
$=$ direct $(K^{-1})^{\otimes n}$ deconvolution, exact; V1 the slope maps and
sign factor, sympy; V2 the sign-chain parity identity $=$ direct sign,
exhaustive $n\le10$ below \emph{and} above $D_c$; V3 the criterion --- exact
rational certificates constructed and exhaustive words positive, plus a
$50$-point high-precision existence grid with zero failures; V4 sharpness ---
$\operatorname{disc}(B_1B_0)<0$ and the alternating word negative just above
$D_c$; V5 the Gray cross-check; V6 BSC divisibility).

\textbf{Lemma 7.34h (closure of G1: the chain \eqref{eq:G1chain} holds for all
$a,b\in(0,1)$, $a+b<1$, $\zeta\in(1,\zeta_c)$).} \emph{With $a_1$ the larger
root of $\mathcal Q$ and $a_2,b_1,b_2,A,B$ as in part (IV) of Lemma 7.34g,
every link of \eqref{eq:G1chain} holds strictly; consequently $A\subset
(p_\pi,0)$, $B\subset(-\infty,p_T)$, $\nu\in A$, $\nu^{-1}\in B$, and $m_0(A\cup
B)\subseteq A$, $m_1(A\cup B)\subseteq B$. The certificate criterion {\rm(I)} of
Lemma 7.34g is therefore met for every $(a,b,\zeta)$ in range.}

\emph{Proof.} Throughout, $a,b\in(0,1)$, $a+b<1$, $\zeta\in(1,\zeta_c)$, so
$\nu=(1-\zeta)/(1+\zeta)\in(-1,0)$. Write
$f_a:=(a-b)\zeta-(a+b)$, $W:=(a-b)\zeta+(2-a-b)$,
$W':=(a-b)\zeta+(a+b-2)$, and the standing factorisations from Lemma 7.34f,
$16\,\mathrm{disc}=F_1F_2$, with the cross-difference
$(a+b)^2[(2-a-b)^2-4ab]-(2-a-b)^2(a-b)^2=16ab(1-a-b)>0$, which forces
$\zeta_1:=(a+b)/|a-b|>\zeta_c$ (the root of $F_1$), and the elementary
identity $(2-a-b)^2-4ab=4(1-a-b)+(a-b)^2$.

\emph{(0) Foundations.} $W>0$ on $(1,\zeta_c]$: if $a\ge b$, $W$ is
nondecreasing in $\zeta$ with $W(1)=2(1-b)>0$; if $a<b$, $W$ is nonincreasing
and $W(\zeta_c)>0$ because $(2-a-b)^2-(b-a)^2-4ab=-4(a+b-1)>0$. Hence
$A_\zeta=-aW/(\zeta-1)<0$: $\mathcal Q$ is a downward parabola. Its
discriminant is $F_1F_2/(\zeta^2-1)^2>0$ on the open interval (both $F_1,F_2>0$:
$F_1(1)=4ab>0$ and $\zeta_1>\zeta_c$; $F_2>0$ strictly below $\zeta_c$), so
$\mathcal Q$ has two distinct real roots; write $r_-<r_+=:a_1$. Their product
$C_\zeta/A_\zeta=b(\zeta-1)/(a(\zeta+1))=-(b/a)\nu>0$ and sum
$-B_\zeta/A_\zeta$ are both positive-product/negative-sum because the trace
numerator $\mathrm{Bnum}:=(a-b)(a+b-2)(\zeta^2-1)-4ab\zeta$ is negative on
$(1,\zeta_c]$ ($\mathrm{Bnum}(1)=-4ab<0$; for $a<b$ it is upward in $\zeta$ with
$\mathrm{Bnum}(\zeta_c)=-(4ab(2-a-b)/d^2)\bigl((a-b)+d\bigr)<0$,
$d=\sqrt{(2-a-b)^2-4ab}>|a-b|$); thus \emph{both roots are negative}.

\emph{Q-test.} For the downward parabola $\mathcal Q$ with vertex
$v=(r_-+r_+)/2$ and $\mathcal Q'(s)=2A_\zeta(s-v)$, one has the elementary
equivalences $\ s<a_1=r_+\iff \mathcal Q(s)>0\ \text{or}\ \mathcal Q'(s)>0$,
and $\ s>a_1\iff \mathcal Q(s)<0\ \text{and}\ \mathcal Q'(s)<0$.

\emph{(1) $a_1>p_\pi$ and $a_1<\nu$.} A direct computation gives
$\mathcal Q(p_\pi)=-2b\zeta(a+b-1)f_a/\bigl(a(\zeta^2-1)\bigr)$ and
$\mathcal Q(\nu)=-2\zeta(a+b-1)f_a/(\zeta+1)^2$. Now $f_a<0$ on
$(1,\zeta_c]$ ($f_a(1)=-2b<0$ and its root is $\zeta_1>\zeta_c$), so both
values are $<0$ (signs $(-)(-)(-)$ over a positive denominator). Their
derivatives factor as $\mathcal Q'(p_\pi)=f_a\,g_1/(\zeta^2-1)$ with
$g_1:=(a+3b-2)\zeta+(a+b-2)$, and $\mathcal Q'(\nu)=f_a\,g_2/(\zeta^2-1)$ with
$g_2:=(3a+b-2)\zeta+(2-a-b)$. Here $g_1<0$: $g_1(1)=2(a+2b-2)<0$, and if
$a+3b-2>0$ then $g_1(\zeta_c)<0$ because $(2-a-b)^2-4ab-(a+3b-2)^2=8b(1-a-b)>0$;
while $g_2>0$: $g_2(1)=2a>0$, and if $3a+b-2<0$ then $g_2(\zeta_c)>0$ because
$(2-a-b)^2-4ab-(2-3a-b)^2=8a(1-a-b)>0$. Thus $\mathcal Q'(p_\pi)>0$
(so $p_\pi<a_1$, i.e.\ $a_1>p_\pi$) and $\mathcal Q'(\nu)<0$ with
$\mathcal Q(\nu)<0$ (so $\nu>a_1$, i.e.\ $a_1<\nu$). In particular $a_1>p_\pi>
p_T$ since $p_\pi-p_T=(1-a-b)/a>0$, so the denominator $(1-a)+aa_1>0$.

\emph{(2) $\nu<a_2<0$.} The pointwise identity
$\psi(s)-1=(a+b-1)(1-s)/\bigl((1-a)+as\bigr)$ shows $\psi(a_1)<1$ (numerator
$<0$ as $a+b<1$, $a_1<0$; denominator $>0$), hence $a_2=\nu\psi(a_1)>\nu$
($\nu<0$). For $a_2<0$ it suffices that $\psi(a_1)>0$, i.e.\ $a_1>-b/(1-b)$, the
zero of $\psi$ (the denominator being positive). Since $a_1=r_+$ is the
\emph{larger} root, it exceeds the vertex $v:=\tfrac12(r_-+r_+)=-B_\zeta/(2A_\zeta)$
of the downward parabola; and $v>-b/(1-b)$ because
\[
  v+\frac{b}{1-b}=\frac{\mathcal N(a,b,\zeta)}{2a(1-b)(\zeta+1)W},
  \qquad \mathcal N=U(a,b)\,\zeta^2+V(a,b)\ \text{affine in }\zeta^2,
\]
has positive denominator ($W>0$ from (0)) and a numerator positive at both
endpoints of $[1,\zeta_c]$: $\ \mathcal N|_{\zeta=1}=4ab(1-b)>0\ $ and
$\ \mathcal N|_{\zeta=\zeta_c}=8ab\,(2-a-b)(1-a-b)/\bigl((2-a-b)^2-4ab\bigr)>0\ $
(its denominator $=4(1-a-b)+(a-b)^2>0$), so $\mathcal N>0$ on all of
$[1,\zeta_c]$ by affineness in $\zeta^2$. Hence $-b/(1-b)<v<a_1$, whence
$\psi(a_1)>0$, giving $a_2<0$; with $\psi(a_1)<1$ this is $\nu<a_2<0$.
(The direct value $\mathcal Q(-b/(1-b))=-b(a+b-1)W'/\bigl((b-1)^2(\zeta+1)\bigr)$
is in fact $<0$ --- the point lies below \emph{both} roots --- consistent with
$-b/(1-b)<r_-<v<a_1$.)

\emph{(3) $b_2<p_T$ and $b_1<\nu^{-1}<b_2$.} Since $\psi(a_1)\in(0,1)$ from (2),
$b_2=\psi(a_1)/\nu$ satisfies $b_2-\nu^{-1}=(\psi(a_1)-1)/\nu>0$, i.e.\
$\nu^{-1}<b_2$. The inequality $b_2<p_T$ is equivalent (clearing the negative
denominator $\nu$) to $a\psi(a_1)+\nu(1-a)>0$, whose numerator is
\emph{linear} in $a_1$ with leading coefficient $aW>0$; thus it reads
$a_1>s_4$ where $\mathcal Q(s_4)=-(\zeta-1)(a+b-1)^2/(aW)<0$ and
$\mathcal Q'(s_4)=-\mathrm{Bnum}/(\zeta^2-1)>0$ (as $\mathrm{Bnum}<0$ from (0)),
so by the Q-test $s_4<a_1$ and $b_2<p_T$. Finally
$b_1=\psi(b_2)/\nu<\nu^{-1}\iff\psi(b_2)>1$, and
$\psi(b_2)-1=(a+b-1)(1-b_2)/\bigl((1-a)+ab_2\bigr)>0$ because $b_2<p_T$ makes the
denominator negative while the numerator $(a+b-1)(1-b_2)<0$. This proves
$b_1<\nu^{-1}<b_2$, and $p_T<p_\pi$ is $(1-a-b)/a>0$.

\emph{(4) Closure.} The maps $m_0,m_1$ are M\"obius homeomorphisms of
$\mathbb{RP}^1$ sharing the single pole $p_T$ (denominator $(1-a)+as=a(s-p_T)$),
with $m_y(p_T)=\infty$. The four identities $m_0(a_1)=a_2$, $m_1(a_1)=b_2$,
$m_1(b_2)=b_1$, and $m_0(b_2)=a_1$ (the last because $a_1$ is fixed by
$m_0\circ m_1$) are exact. Let $J:=\{s\ge a_1\}\cup\{s\le b_2\}\cup\{\infty\}$,
the closed arc complementary to the open arc $(b_2,a_1)$. By \eqref{eq:G1chain},
$A=[a_1,a_2]\subseteq\{s\ge a_1\}\subseteq J$ and $B=[b_1,b_2]\subseteq\{s\le
b_2\}\subseteq J$, while $p_T\in(b_2,a_1)$, so $p_T\notin J$. Hence $m_0$ is a
pole-free homeomorphism on $J$; by injectivity $\infty=m_0(p_T)\notin m_0(J)$,
so $m_0(J)$ is the \emph{finite} closed arc between $m_0(a_1)=a_2$ and
$m_0(b_2)=a_1$, namely $[a_1,a_2]=A$. Identically $m_1(J)=[b_1,b_2]=B$. Thus
$m_0(A\cup B)\subseteq m_0(J)=A$ and $m_1(A\cup B)\subseteq m_1(J)=B$. Together
with $\nu\in A$, $\nu^{-1}\in B$ (from \eqref{eq:G1chain}) this is exactly the
hypothesis of criterion {\rm(I)}. \hfill$\square$

\emph{Validated} in
\texttt{scripts/verify/lemma\_7\_34g\_G1\_uniform\_existence.py} (Step 0 the
fixed-point quadratic $=(A_\zeta,B_\zeta,C_\zeta)$, $\det=(1-a-b)^2$,
$\mathrm{disc}=F_1F_2/(\zeta^2-1)^2$; Step 1 the foundations $W>0$, $A_\zeta<0$,
$F_1>0$, both roots negative; Steps 2--3 the closed-form sign reductions of all
eight links via the Q-test, every algebraic identity confirmed symbolically;
Step 4 the closure collapse $p_T\notin J\Rightarrow m_0(J)=A$, $m_1(J)=B$,
including $m_0(m_1(a_1))=a_1$ modulo the fixed-point quadratic; Step 5 a
$1352$-point high-precision witness for \eqref{eq:G1chain} with positive margins
and an independent corner/seed check of $m_0(A\cup B)\subseteq A$,
$m_1(A\cup B)\subseteq B$ --- all PASS).

\textbf{Theorem 7.34i (asymmetric BSMS operational rate-distortion dispersion;
the consolidation of instance (2)).} \emph{Let $X$ be the stationary binary
Markov chain of Lemma 7.34f --- $P(0\!\to\!1)=a$, $P(1\!\to\!0)=b$,
$a,b\in(0,1)$, $a+b<1$, $a\ne b$ --- under Hamming distortion and the
fixed-$\varepsilon$ excess-distortion criterion. Then for every
$D\in(0,D_c(a,b))$ at which the one-shot spectral certificate of
Lemma 7.34e\,(iii) succeeds --- every $(a,b,D)$ tested, and unconditionally on
any bounded-mixing subfamily $\{a+b\le1-\eta\}$ below threshold (Remark
7.34e$'$, by compactness) --- the operational second-order dispersion equals
the lossless varentropy rate of $X$ and is $D$-independent across the Gray
interval:}
\[
  V_{\mathrm{op}}(D)\;=\;V_{\mathrm{lossless}}(a,b)
  \;:=\;\lim_{n}\frac1n\operatorname{Var}\bigl(-\log_2 P(X^n)\bigr),
\]
\emph{the Gordin/Green--Kubo varentropy rate, which decomposes as}
\[
  V_{\mathrm{lossless}}(a,b)=\underbrace{\pi_0\,a(1-a)\log_2^2\tfrac{1-a}{a}
  +\pi_1\,b(1-b)\log_2^2\tfrac{1-b}{b}}_{\text{conditional-varentropy mixture}}
  \;+\;V_{\mathrm{mem}}(a,b),\quad \pi_0=\tfrac b{a+b},
\]
\emph{the stationary mixture of the two per-state conditional varentropies plus
a memory term $V_{\mathrm{mem}}(a,b)\ge0$ with $V_{\mathrm{mem}}\equiv0$ iff
$a=b$ (so the asymmetric rate strictly exceeds the conditional mixture). At
$a=b=p$ the statement reduces to the symmetric Remark 7.34b:
$V_{\mathrm{lossless}}=p(1-p)\log_2^2\tfrac{1-p}p$ and
$D_c=\tfrac12(1-\sqrt{1-2p}/(1-p))$. This is the first \emph{asymmetric}
discrete memory source with a proven operational rate-distortion dispersion.}

\emph{Proof (by assembly of the committed pieces).} \emph{Converse}
($V_{\mathrm{op}}\ge V_{\mathrm{lossless}}$): Remark 7.34d. In the Gray region
the SLB is tight and the chain-agnostic dual identity gives the
\emph{deterministic} shift $\jmath_n=\imath_n-nh(D)+c_n$ with $c_n$
word-independent, so $\operatorname{Var}(\jmath_n)=\operatorname{Var}(\imath_n)$
at every $n$, $\rho:=\operatorname{Var}(\jmath_n)/\operatorname{Var}(\imath_n)=1$,
and the Kostina--Verd\'u fixed-$\varepsilon$ liminf converse yields
$V_{\mathrm{op}}\ge\lim_n\operatorname{Var}(\jmath_n)/n=V_{\mathrm{lossless}}(a,b)$
(centering $\kappa_n=H(X^n)-n\bar h=O(1)$, $n$-stable). \emph{Achievability}
($V_{\mathrm{op}}\le V_{\mathrm{lossless}}$): the symmetric closure machinery of
Remark 7.34b --- the quenched uniform tail lemma, the KP window analyticity and
Cauchy bounds, the (R2) ordered-constants central window, and the complex-saddle
GAP-1 lemma --- is \emph{chain-agnostic}, consuming only (i) the replica
spectral gap $\bar g=\max_{[\varepsilon,\pi]}g<1$ with the positive curvature
floor $\bar v=D(1-D)(1-\bar\Psi_0)>0$, supplied per-$(a,b,D)$ by
Lemma 7.34e\,(iii) --- the symmetric constant $\tfrac{11}{16}$ enters the
GAP-1 write-out solely as $c_g=\tfrac{11}{16}D(1-D)$, and every downstream
constant ($c_2=\tfrac{c_g}{4\pi^2}$, $c_4=\tfrac{c_g}{16\pi^2}$,
$c_5=\tfrac{c_g}{4\pi^2}$, the $E_{2'}$ margin $\tfrac{c_g}{4\pi^2}>0$) is a
positive multiple of $c_g$, so the substitution
$c_g\mapsto\tilde c(a,b,D)D(1-D)$, $\tilde c\in[\tfrac{43}{64},\tfrac{25}{32}]$,
is monotone with no sign change anywhere in the chain; (ii) the Markov
bounded-differences concentration of $\mu_x$, supplied by Paulin (2015,
Thm.~2.1; Marton coupling) with Dobrushin coefficient $\delta=|1-a-b|$ and
factor $1/(a+b)^2$; and (iii) the all-$n$ validity of the $\mathrm{BSC}(D)$
deconvolution on $(0,D_c(a,b))$, supplied by Lemmas 7.34f--h; together with the
chain-agnostic facts $K\sim\mathrm{Bin}(n,D)$, $\mathbb E[\mu_x]=nD$,
$\Phi(0;x)=1$, $\lambda^*=\ln\tfrac{1-D}D$. The \emph{only} $a=b$-specific
structure (global-flip $W$-invariance, the $W_4=\mathbf 1m^T+q\Pi M$ form, the
closed-form $\tfrac{11}{16}$ replica inequality) is confined to the symmetric
replica theorem, i.e.\ to item (i), which Lemma 7.34e replaces at proof grade
for $a\ne b$. The open-interval scope avoids the symmetric Möbius-orbit
endpoint floor (a genuinely $a=b$-specific lemma); the closed endpoint
$D=D_c(a,b)$ is reached per-$(a,b)$ by the closed-form-in-$D$ $s=\pi$ certificate
of Remark 7.34e$'$. \hfill$\square$

\emph{Scope (honest).} The statement is per-source and per-$(a,b,D)$: item (i)
is certified by one finite exact computation per point, not a single
closed-form inequality. The uniform-in-$(a,b,D)$ version reduces to the
convexity conjecture of Remark 7.34e$'$ and otherwise holds only on
bounded-mixing subfamilies ($\bar\Psi_0\to1$ as $a+b\to1$ on asymmetric paths,
where the curvature degenerates). The discrete-Markov dispersion beyond the
binary chain, and beyond the Gray region ($D>D_c$), remain open (Remarks
7.34m$'''$ and 7.34o).
\emph{Validated} in
\texttt{scripts/verify/theorem\_7\_34i\_asym\_dispersion.py} (T1 the Gordin
rate, the $V=V_{\mathrm{mix}}+V_{\mathrm{mem}}$ decomposition with
$V_{\mathrm{mem}}=0$ iff $a=b$, the simulation cross-check, and the exact
symmetric reduction; T2 the $c_g$-monotonicity of the GAP-1 constants), on top
of the per-piece verifiers
(\texttt{remark\_7\_34d\_general\_converse\_discrete\_markov.py},
\texttt{probe\_7\_34\_nonsym\_dual\_identity.py},
\texttt{probe\_7\_34\_asym\_replica\_structure.py},
\texttt{lemma\_7\_34e\_spi\_closed\_form.py},
\texttt{lemma\_7\_34g\_sign\_automaton.py},
\texttt{lemma\_7\_34g\_G1\_uniform\_existence.py}).

\textbf{Lemma 7.34j (the first leaf beyond binary: the ternary symmetric Gray
threshold).} \emph{Let $X$ be the stationary symmetric ternary Markov chain on
$\{0,1,2\}$ with $T(i,i)=1-p$, $T(i,j)=p/2$ ($j\ne i$), $\pi$ uniform, under
Hamming distortion, and let the backward test channel be the ternary symmetric
channel $K=(1-\tfrac{3D}2)I+\tfrac D2 J$ ($J$ all-ones), so that
$K^{-1}=\tfrac1{2-3D}\,(2I-D\,J)$ ($D\ne\tfrac23$). Define the all-$n$ Gray
threshold $D_c^{(3)}(p)=\sup\{D:(K^{-1})^{\otimes n}P_X\ge0\ \forall n\}$. Then
for $p\in(0,p^*)$, $p^*\approx0.67$:}
\begin{itemize}
\item[(i)] \emph{$D_c^{(3)}(p)>0$ --- the ternary Gray region is nonempty;}
\item[(ii)] \emph{the binding word is the period-$2$ two-symbol \emph{alternating}
word (using only $2$ of the $3$ states), governed by the $3\times3$ period
product $M=B_1B_0$, $B_y[a,b]=K^{-1}[y,a]\,T(a,b)$; $D_c^{(3)}(p)$ is the
complex-onset of $M$'s dominant eigenpair --- the smallest root in $(0,\tfrac23)$
of the cubic-discriminant binding factor, the quartic}
\[\begin{aligned}
  Q(D,p)=\;&9p^2(3p-4)^2D^4-4p(3p-4)(27p^2-30p-4)D^3\\
  &+(441p^4-864p^3+432p^2+64)D^2\\
  &-8(27p^4-33p^3+36p^2-32p+16)D+16p^2(p^2+1);
\end{aligned}\]
\emph{the period-$3$ word $012\cdots$ is phase-protected (complex eigenpair but
never binds) and constant/run words have real $M$-spectra and never oscillate;}
\item[(iii)] \emph{$Q$ is \emph{irreducible} over $\mathbb Q(p)$ and does not
reduce through a quadratic in any $1/(1-2D)$-type variable (the mechanism behind
the binary clean radical); hence, although $D_c^{(3)}(p)$ as a quartic root does
admit the general nested-radical quartic-formula expression, it has \emph{no
clean single-radical closed form} analogous to the binary
$D_c=\tfrac12(1-\sqrt{1-2p}/(1-p))$ --- a genuine degree-$4$ algebraic number,
the structural break from Lemma 7.34f, traced to the period product being
$3\times3$ rather than $2\times2$;}
\item[(iv)] \emph{small-$p$, $D_c^{(3)}(p)=\tfrac{p^2}8+\tfrac{p^3}4
+\tfrac{45p^4}{128}+O(p^5)$, exactly \emph{half} the binary leading coefficient
($D_c=\tfrac{p^2}4+\tfrac{p^3}2+\cdots$), so the ternary Gray region is
asymptotically half the binary one.}
\end{itemize}
\emph{Proof grade.} (i)--(iv) are PROVEN by exact symbolic computation ($K^{-1}$,
the quartic $Q$ and its irreducibility, the matched small-$p$ series) together
with the period-$2$ eigenvalue-discriminant mechanism (the two dominant
eigenvalues of $M$ collide into a complex-conjugate pair $re^{\pm i\theta}$ at
the smallest root of $Q$, after which $P_{Y^*}(\mathrm{alt})\sim r^n\cos(n\theta
+\varphi)$ oscillates negative; just above threshold the first-negative index
diverges as $\pi/2\theta$, the ternary analog of the binary creep) plus an
exhaustive word-positivity scan --- the same proof grade as Lemma 7.34f.
\emph{Open} (the ternary analog of Lemmas 7.34g/h): the word-uniform
invariant-cone ``all words, all $n$'' certificate; the non-monotone large-$p$
regime ($p\gtrsim0.67$, where the constant word binds first --- a near-i.i.d.\
crossover); and the full A-ary ($A\ge4$) generalisation (now Lemma 7.34k).
\emph{Validated} in
\texttt{scripts/verify/probe\_7\_34j\_ternary\_gray.py} (V1 the exact $K^{-1}$;
V2 the transfer form vs full brute deconvolution; V3 $Q$ irreducible and its
smallest root $=$ the alternating-word threshold; V4 binding word alternating;
V5 the $p^2/8$ vs binary $p^2/4$ ratio $\tfrac12$; V6 the large-$p$
constant-word crossover --- all PASS).

\textbf{Lemma 7.34k (the A-ary symmetric Gray threshold: the general alphabet
law).} \emph{Let $X$ be the stationary A-ary symmetric Markov chain on
$\{0,\dots,A-1\}$ ($A\ge2$) with $T(i,i)=1-p$, $T(i,j)=p/(A-1)$ ($j\ne i$),
$\pi$ uniform, under Hamming distortion, and let the backward test channel be
the A-ary symmetric channel $K=(1-D-b_0)I+b_0J$, $b_0=D/(A-1)$, so
$K^{-1}=\tfrac1{1-D-b_0}\,(I-b_0J)$. With $D_c^{(A)}(p)=\sup\{D:(K^{-1})^{\otimes
n}P_X\ge0\ \forall n\}$:}
\begin{itemize}
\item[(i)] \emph{for every $A\ge2$ the binding word is the period-$2$ two-symbol
\emph{alternating} word; under the residual $S_{A-2}$ symmetry the $A-2$
non-participating symbols split into one coupled average mode and $A-3$
decoupled \emph{spectator} modes, so the period product $M=B_1B_0$
($B_y[a,b]=K^{-1}[y,a]\,T(a,b)$) has characteristic polynomial $(\text{relevant
cubic})\cdot(\text{spectator linear})^{A-3}$ for $A\ge3$ (a quadratic for
$A=2$);}
\item[(ii)] \emph{$D_c^{(A)}(p)$ is the complex-onset of the relevant block's
dominant eigenpair --- a clean radical ($\sqrt{\ }$) only for $A=2$ (binary,
Lemma 7.34f), and the smallest positive root of an $A$-dependent
\emph{irreducible quartic} (the cubic-discriminant) for all $A\ge3$ (the case
$A=3$ is Lemma 7.34j); the spectator modes add no degree, so the threshold stays
a quartic for \emph{all} $A\ge3$;}
\item[(iii)] \emph{small-$p$, $D_c^{(A)}(p)=\dfrac{p^2}{4(A-1)}+O(p^3)$: the Gray
region is strictly decreasing in $A$ and shrinks as $1/(A-1)$ in the alphabet
size.}
\end{itemize}
\emph{Proof grade.} parts (i)--(ii) are \emph{proven for every $A\ge3$} by the centralizer-algebra
argument: $B_y=D_yT$ with $D_y=\mathrm{diag}(K^{-1}[y,\cdot])$ uniform on the
spectator block, so $M=B_1B_0$ commutes with the $S_{A-2}$-action permuting
$\{2,\dots,A-1\}$; under $S_{A-2}$, $\mathbb R^A=\mathrm{span}\{e_0\}\oplus
\mathrm{span}\{e_1\}\oplus\mathrm{span}\{\sum_{i\ge2}e_i\}\oplus V_{\mathrm{std}}$
($V_{\mathrm{std}}$ the $(A-3)$-dimensional standard irrep), and by Schur's lemma
$M$ acts as a scalar on $V_{\mathrm{std}}$ --- the $(A-3)$-fold spectator
eigenvalue --- while the three trivial-isotypic lines form the $3$-dimensional
relevant block (cubic) for \emph{all} $A\ge3$, so the binding factor is a quartic
\emph{independent of $A$}; (the case checks are exact symbolic for $A=3,4,5$). The
small-$p$ coefficient $1/(4(A-1))$ is exact-symbolic for $A\le4$ and numeric
($D_c^{(A)}/p^2\to1/(4(A-1))$) for $A\le8$; the binding word is the alternating
one by exhaustive scan ($A\le5$, all $A^n$ words $n\le7$: no word binds below the
alternating-onset). \emph{Open} (as for Lemma 7.34j): the
word-uniform all-$n$ certificate; the large-$p$ constant-word crossover; and
exact $A$-uniform symbolic forms beyond the small-$p$ coefficient.
\emph{Validated} in \texttt{scripts/verify/probe\_7\_34k\_aary\_gray.py} (K1 the
A-ary $K^{-1}$; K2 the spectator decoupling; K3 the quartic binding factor for
$A=3,4$; K4 the $1/(4(A-1))$ small-$p$ law; K5 the alternating-onset as the
all-$n$ threshold; K6 monotonicity in $A$ --- all PASS).

\textbf{Lemma 7.34l (the A-ary symmetric replica spectral inequality: the
source-side transfer operator, the $S_A$ pattern-quotient, and the
alphabet-uniform curvature floor).} \emph{Let $X$ be the stationary A-ary
symmetric Markov chain of Lemma 7.34k ($T(i,i)=1-p$, $T(i,j)=p/(A-1)$, $\pi$
uniform) under Hamming distortion, let $0<D\le D_c^{(A)}(p)$ be in the A-ary
Gray region, write $\lambda^*=\ln\tfrac{(1-D)(A-1)}{D}$, $\nu:=e^{-\lambda^*}
=\tfrac{D}{(1-D)(A-1)}$, and let $\Phi(s;x^n):=\mathbb E_{Y\sim Q(\cdot\mid
x^n)}[e^{is\,d_H(x^n,Y)}]$ be the posterior characteristic function of the
distortion under the matched backward channel $Q(y\mid x)\propto
P_{Y^*}(y)e^{-\lambda^* d_H(x,y)}$.}

\emph{\textbf{(i) (The dual identity / structural cancellation.)} On the Gray
region the posterior normaliser factorises exactly:
$M(x^n):=\sum_y P_{Y^*}(y)e^{-\lambda^* d_H(x,y)}=\gamma^n P_X(x^n)$ with
$\gamma=\gamma(p,D)$ \emph{independent of $x^n$} (to $10^{-15}$), the A-ary
analogue of the binary $\Phi(0;x)=1$ cancellation. Consequently the annealed
second moment is a \emph{source-side} transfer-matrix product on triples
$(x,x',x'')\in\{0,\dots,A-1\}^3$:}
\[
  \mathbb E_x\big|\Phi(s;x^n)\big|^2=\gamma^{-2n}\,u(s)^{T}W_A(s)^{\,n-1}\mathbf 1,
\qquad
  W_A\big[(x_a,x_b,x_c),(x,x',x'')\big]=\frac{T_{x_bx'}T_{x_cx''}}{T_{x_ax}}\,
  G_s(x,x')\,\overline{G_s(x,x'')},
\]
\emph{with the dressed single-site kernel $G_s(x,x')=\sum_y K^{-1}[y,x']\,
f_s(x,y)$, $f_s(x,y)=1$ ($y=x$), $\nu e^{is}$ ($y\ne x$), and
$u(s)[(x,x',x'')]=A^{-1}G_s(x,x')\overline{G_s(x,x'')}$. \emph{(The naive
single-site \emph{output}-side operator --- which reproduces the second moment
for $A=2$ --- does not extend to $A\ge3$; the source-side form above is exact for
every $A$, validated against the brute $A^n$ second moment to $10^{-15}$.)} Set
$g_A(s):=\rho(W_A(s))/\gamma^2$ (the per-site replica rate; $g_A(0)=1$ by
Perron normalisation, since $\mathbb E_x|\Phi(0;x)|^2=1$ identically).}

\emph{\textbf{(ii) ($S_A$ pattern-quotient: an $A$-independent finite spectral
problem.)} $W_A$ commutes with the diagonal $S_A$-action on triples;
restricted to the trivial isotypic it collapses to the \emph{equality-pattern
quotient} $Q(s)$ of dimension $5$ for $A\ge3$ ($4$ for $A=2$) --- the orbits
being all-equal, $x{=}x'{\ne}x''$, $x{=}x''{\ne}x'$, $x'{=}x''{\ne}x$, and
all-distinct --- and $\rho(W_A(s))=\rho(Q(s))$ exactly (Perron sits in the
pattern-averaged block). Thus $A$ enters only as a \emph{parameter} of a fixed
$\le5\times5$ matrix, not as the dimension --- the A-ary analogue of the
binary global-flip reduction $W_8\to W_4$.}

\emph{\textbf{(iii) (Spectral gap, curvature, alphabet-uniform floor.)} For
every $A\ge2$, every $p$ with $V_{\mathrm{lossless}}^{(A)}>0$ (i.e.\
$p\ne\tfrac{A-1}A$), every $0<D\le D_c^{(A)}(p)$ and every $s\in(0,\pi]$,}
\[
  g_A(s)\;\le\;1-\tfrac{11}{16}\,D(1-D)(1-\cos s)\;<\;1,
\]
\emph{and the binding corner (small $p$, $D=D_c$, $s=\pi$) has ratio
$(1-g_A)/\bigl(D(1-D)(1-\cos s)\bigr)\approx\tfrac32$ --- the binary small-$p$
limit ($\approx1.50$ at $A=2$, rising slightly to $\approx1.53$ with $A$) ---
comfortably above the conservative alphabet-uniform floor $\tfrac{11}{16}$ that
the achievability machinery actually consumes. The curvature is positive
and equals the per-site posterior variance of the Hamming distortion,
$g_A(s)=1-\bar v_A s^2+O(s^4)$, $\bar v_A=D(1-D)\bigl[1-\kappa_A^2
D(1-D)/(1-2D)^2\bigr]>0$, where $\kappa_A^2$ is a positive rational of
$(A,p)$ recovering the binary $\kappa_2^2=(1-2p)^2/(p^2(1-p)^2)$ at $A=2$
(known in closed form only at $A=2$; numerically $\kappa_A^2$ is increasing
in $A$ at fixed $p$, $D$); $s=\pi$ is the \emph{easiest} point (smallest
$g_A$, largest margin), exactly as in the binary case.}

\emph{Proof grade.} (i) is exact: the dual factorisation $M=\gamma^nP_X$ is
machine-checked $x$-independent (std $4\times10^{-15}$), and the source-side
operator reproduces the brute $A^n$-word second moment to $9\times10^{-15}$ for
$A=2,3,4,5$, $n\le4$. (ii) is the exact $S_A$-equivariance / Schur reduction
(the same centralizer-algebra argument as Lemma 7.34k(i)), verified
$|\rho(W_A)-\rho(Q)|<3\times10^{-15}$ for $A=2{,}\dots,6$. (iii): the floor
inequality and the $\to\tfrac32$ binding-corner limit are established by a
hostile grid (geometric near $0$ plus a $\pi$-dense tail) across $A\le6$,
$p\in\{0.02,\dots,0.45\tfrac{A-1}A\}$, $D=D_c$; the curvature$=$variance
identity matches the brute per-site posterior variance to $\sim8\%$ at $n=5$
(the $O(1/n)$ finite-$n$ bias). What is \emph{proven uniformly in $A$} (not
merely tested) is the reduction (ii) to a single $\le5$-dimensional spectral
problem, the exact dual cancellation (i), and positivity of the curvature; the
remaining residual --- a fully symbolic Sturm-style certification of
$\sup_{[\varepsilon,\pi]}g_A<1$ on the $\le5\times5$ quotient over
$\mathbb Q(p,D)$, exactly as Lemma 7.34e supplied for the asymmetric binary
$8\times8$ --- is mechanical per $(A,p,D)$ and discharged numerically here.
(The $A=2$ instance of the floor is now \emph{sharp}: $R_2\ge\tfrac32$ in
closed form --- the upgraded replica theorem of Remark 7.34b; the general-$A$
sharp-floor campaign, including its regional scope map, is Remarks
7.34m$''''$/7.34m$'''''$.)
\emph{Validated} in \texttt{scripts/verify/probe\_7\_34l\_aary\_replica.py}
(V0 dual identity; V1 operator$=$brute; V2 $\sup g_A<1$ on the Gray region;
V3 curvature$=$variance; V4 $s=\pi$ binding; V5 off-Gray control;
V6 the $\le5$-dim $S_A$ quotient; V7 the alphabet-uniform $\tfrac{11}{16}$
floor; V8 the A-ary Markov--McDiarmid Dobrushin input; V9 the converse ratio
$\rho=1$; V10 the closed-form $V_{\mathrm{lossless}}^{(A)}$ --- all PASS),
cross-validated by an \emph{independent} operator construction (the closed-form
disagreement-tilt $W_{A^3}$, vs.\ the source-side $G_s$ form above) in
\texttt{scripts/verify/probe\_7\_34l\_aary\_dispersion.py} (C0--C4: the dual
identity, $W_A=$ brute, $g_A<1$ with curvature, the Green--Kubo varentropy
$=$ closed form, off-Gray validity --- all PASS).

\textbf{Theorem 7.34m (the A-ary symmetric Gray-region rate--distortion
dispersion).} \emph{Let $X$ be the stationary A-ary symmetric Markov chain
($T(i,i)=1-p$, $T(i,j)=p/(A-1)$, $\pi$ uniform) under Hamming distortion, with
$p\in(0,\tfrac{A-1}A)$ \emph{(}so $V_{\mathrm{lossless}}^{(A)}>0$\emph{)}, and
let $0<D\le D_c^{(A)}(p)$ be in the A-ary Gray region (Lemma 7.34k). Then the
operational fixed-$\varepsilon$ excess-distortion dispersion satisfies the}
\textbf{converse}
\[
  V_{\mathrm{op}}(D)\;\ge\;V_{\mathrm{lossless}}^{(A)}
  \;=\;p(1-p)\log_2^2\frac{(1-p)(A-1)}{p}
\]
\emph{\textbf{unconditionally}, and the matching} \textbf{achievability}
$V_{\mathrm{op}}(D)\le V_{\mathrm{lossless}}^{(A)}$ \emph{holds modulo the same
single residual as the asymmetric binary instance --- a symbolic Sturm
certification of $\sup_{[\varepsilon,\pi]}g_A(s)<1$ on the $\le5\times5$
$S_A$-quotient of Lemma 7.34l(ii), discharged numerically throughout the
region --- so that $V_{\mathrm{op}}(D)=V_{\mathrm{lossless}}^{(A)}$.}

\emph{Proof.} \emph{Converse.} Remark 7.34d is source-agnostic and names the
$q$-ary case explicitly: its hypotheses are (i) a varentropy CLT for
$\imath_n=-\log_2P(X^n)$ and (ii) an all-$n$ SLB-tight (balanced-kernel
deconvolution) region. Here (i) holds because the increments
$\xi_i=-\log T(X_{i-1},X_i)$ are \emph{i.i.d.}\ for the symmetric chain (the
increment law is state-independent: $-\log(1-p)$ w.p.\ $1-p$,
$-\log\tfrac p{A-1}$ w.p.\ $p$), so $\imath_n$ is a sum of i.i.d.\ terms with
$\mathrm{Var}(\imath_n)/n\to V_{\mathrm{lossless}}^{(A)}=p(1-p)\log_2^2
\tfrac{(1-p)(A-1)}p>0$ (no Gordin cross-terms; binary is $A=2$); and (ii) is
precisely the A-ary Gray region of Lemma 7.34k, with the $q$-ary symmetric
backward kernel of noise entropy $h(D)+D\log(A-1)$. On the Gray region the SLB
shift is deterministic ($\jmath_n=\imath_n-n[h(D)+D\log_2(A-1)]+c_n$, $c_n$
$x$-independent), so $\mathrm{Var}(\jmath_n)=\mathrm{Var}(\imath_n)$ exactly
(the converse ratio $\rho=\mathrm{Var}(\jmath_n)/\mathrm{Var}(\imath_n)=1$ to
$12$ digits, V9), and the one-shot Kostina--Verd\'u $d$-tilted converse with
the $\jmath_n$ CLT gives $V_{\mathrm{op}}\ge V_{\mathrm{lossless}}^{(A)}$ verbatim.

\emph{Achievability.} The chain-agnostic machinery (the quenched tail lemma,
the KP/polymer window, GAP-1, the (R2) variance split, the Abel handoff)
consumes only three inputs, all now A-ary: (a) the replica spectral gap with
positive curvature, $\sup_{[\varepsilon,\pi]}g_A<1$ and
$g_A(s)=1-\bar v_As^2+O(s^4)$, $\bar v_A>0$ --- supplied by Lemma 7.34l, with
the $S_A$ pattern-quotient making it an $A$-independent $\le5$-dimensional
problem and the alphabet-uniform floor $g_A\le1-\tfrac{11}{16}D(1-D)
(1-\cos s)$ (the binding-corner ratio is $\approx\tfrac32$, the binary small-$p$
limit); (b) the A-ary Markov bounded-differences
replacement, supplied by Paulin (2015, \emph{EJP} 20(79), Thm.~2.1; Marton
coupling) with one-step Dobrushin coefficient $\delta=\lambda_2(T)=1-pA/(A-1)<1$,
giving McDiarmid with finite factor $1/(1-\delta)^2=((A-1)/(pA))^2$ (V8); and
(c) the all-$n$ deconvolution validity, supplied by Lemma 7.34k. The single
residual is item (a)'s symbolic certification on the $\le5\times5$ quotient ---
the exact A-ary analogue of the asymmetric binary residual closed by Lemma
7.34e --- which is numerically discharged across the region. Modulo it,
$V_{\mathrm{op}}\le V_{\mathrm{lossless}}^{(A)}$, whence equality. \hfill$\square$

\emph{Scope (honest).} The \emph{converse} is unconditional for every $A\ge2$
and is the first general-alphabet discrete-Markov operational RD-dispersion
\emph{lower bound}; it specialises to the binary Remark 7.34b and to Kostina
2017 in the i.i.d.\ plateau ($p\to$ any, $D$ in the single-letter range). The
\emph{achievability} is at the same proof grade as the asymmetric binary
instance (2) immediately before its closure: every structural ingredient is in
place (the exact source-side operator, the $A$-independent $\le5$-dim
$S_A$-quotient, the uniform $\tfrac{11}{16}$ floor, the Paulin McDiarmid, the
all-$n$ deconvolution), with the lone residual being the mechanical symbolic
Sturm step on a fixed small matrix. The binary closed-form $\tfrac{11}{16}$
chain ($W_8\to W_4$, the $\mathbf 1m^T+q\Pi M$ determinant factorisation) does
not transplant verbatim --- the quotient is $5\times5$, not $2\times2$ --- but
the variational Russo--Dye certificate of Lemma 7.34e does (it is dimension-
and alphabet-agnostic). Beyond-Gray ($D>D_c^{(A)}$), the third-order
$\tfrac12\log n$ prefactor, and the general (non-symmetric, larger-alphabet
\emph{and} non-balanced) discrete-Markov dispersion remain open
(Krishnamachari 2026); Gray (1971, \emph{IEEE Trans.\ IT} 17(2):127--134)
established that for finite-state finite-alphabet Markov sources under balanced
(Hamming/Lee) distortion the SLB is tight on a positive $0\le D\le D_c$ and
\emph{equals the memoryless bound for the same entropy and alphabet} --- the
classical first-order statement this theorem upgrades to second order.
\emph{Validated} in \texttt{scripts/verify/probe\_7\_34l\_aary\_replica.py}
(V0--V10, all PASS).

\textbf{Remark 7.34m$'$ (the A-ary binding endpoint $s=\pi$ and the residual
convexity conjecture --- the A-ary analogue of Remark 7.34e$'$).} The
achievability residual of Theorem 7.34m --- a symbolic certification of
$\sup_{[\varepsilon,\pi]}g_A(s)<1$ on the $\le5\times5$ $S_A$-quotient ---
admits the same sharp structural reduction the asymmetric binary instance
received in Remark 7.34e$'$, with one new feature (the irreducible \emph{quartic}
endpoint block, in place of the binary degree-$6$ one).

\emph{The binding endpoint.} Numerically (a geometric-near-$0$ plus $\pi$-dense
$s$-grid, every $A\le6$, all $(p,D)$ on the Gray region) the floor ratio
$R_A(s):=[1-g_A(s)]/[D(1-D)(1-\cos s)]$ is \emph{minimised at $s=\pi$};
equivalently, in $u=\cos s$ the Perron branch $G(u):=g_A(\arccos u)$ is convex
on $[-1,1]$ with $G(1)=1$, so the chord slope $[1-G(u)]/(1-u)=D(1-D)\,R_A$ is
nondecreasing in $u$ and $R_A$ attains its minimum $R_A(\pi)$ at $u=-1$. The
$s\to0$ end is a closed-form anchor: $R_A(0^+)=2\bar v_A/[D(1-D)]=2[1-\kappa_A^2
D(1-D)/(1-2D)^2]$ (with $\bar v_A$ the per-site posterior variance of
Lemma 7.34l(iii)), which is $\approx2$ for the small $D=O(p^2)$ of the Gray
region, while $R_A(\pi)$ descends from $\approx2$ at $D\to0$ to its $\approx\tfrac32$
infimum at the small-$p$, $D\to D_c$ corner --- so $R_A(0^+)\ge R_A(\pi)$, and the
endpoint \emph{ordering} is robust; only the \emph{interior} convexity is
conjectural. (Because $R_A$ is a $0/0$ limit as $s\to0$, float-$64$ fine grids
are noisy in the $s,D\to0$ corner; $60$-dps recomputation there gives
$R_A(s)-R_A(\pi)>0$ for all $s$, e.g.\ $+2\times10^{-7}$ at $A{=}3,p{=}0.02,
D{=}2\times10^{-6}$, confirming the minimiser is $s=\pi$.) The endpoint is exactly solvable: at $s=\pi$ the disagreement tilt $\nu e^{is}=-\nu$ (with
$\nu=e^{-\lambda^*}=D/((1-D)(A-1))$) is real, so with
\[
  N:=AD-2D^2-(A-1),\qquad \mathrm{Den}:=AD-(A-1)=N+2D^2
  \qquad(\text{both }<0\text{ on the Gray region}),
\]
the contour data is real and rational, \emph{proven uniformly in $A$}:
\[
  \eta_\pi=\frac{2D(1-D)}{N}<0,\qquad
  \Big|\frac{C_\pi}{C_0}\Big|^2=\Big(\frac{N}{\mathrm{Den}}\Big)^2
  =\Big(1-\frac{2D^2}{\mathrm{Den}}\Big)^2,
\]
reducing at $A=2$ to the committed binary $\eta_\pi=-2D(1-D)/((1-D)^2+D^2)$ and
$|C_\pi/C_0|^2=(c^2+4D^2(1-D)^2)/c^2$ ($c=1-2D$) of Remark 7.34e$'$ (the
identity $((1-D)^2+D^2)^2=(1-2D)^2+4D^2(1-D)^2$ matches the two forms).

\emph{The endpoint spectrum (the wall).} The $5\times5$ $S_A$-quotient
characteristic polynomial at $s=\pi$ \emph{factors} as
$(\text{linear})\times(\text{quartic})$, the quartic irreducible over
$\mathbb Q(p,D)$ for $A=3,4$ (symbolically). The linear factor has a clean root
(e.g.\ at $A=3$, $\lambda_{\mathrm{lin}}=D(1-D)(2-3p)^2/[\,p(1-p)(2D^2-3D+2)\,]$)
but is \emph{subdominant}; the spectral radius $\rho(Q_A(\pi))$ is the
\emph{positive} dominant root of the irreducible \emph{quartic} factor --- a
genuine degree-$4$ algebraic number (solvable by radicals via Ferrari, but with
\emph{no clean closed form}: in particular no reduction to a quadratic-in-$\lambda$,
i.e.\ no single-$\sqrt{\,}$ Perron root as in the binary \emph{symmetric} $W_4$
block). This is precisely the wall the binary $W_8\to W_4$ determinant
factorisation $\mathbf1m^T+q\Pi M$ does not cross; the binary \emph{asymmetric}
endpoint of Remark 7.34e$'$ sits one step further, in a degree-$6$
symmetric-pair block (genuinely radical-free).
Because the dominant root is \emph{positive} --- so
$g_A(\pi)=|C_\pi/C_0|^2\,\rho(Q_A(\pi))>0$ is a genuine Perron value --- the
signed-operator Kre\u{\i}n--Rutman/Collatz--Wielandt PSD-cone (Russo--Dye)
certificate of Lemma 7.34e, which is \emph{dimension- and alphabet-agnostic}
($\eta_\pi<0$ makes $Q_A$ signed, forbidding the M-matrix minor test), is the
route to $\rho(Q_A(\pi))<1/|C_\pi/C_0|^2=(\mathrm{Den}/N)^2$, i.e.\
$g_A(\pi)<1$ on the open Gray interior; numerically the gap
$\rho(Q_A(\pi))-(\mathrm{Den}/N)^2$ is strictly negative throughout, growing in
magnitude toward $D_c$ (e.g.\ $\approx-2\times10^{-2}$ at $D\to D_c$, $p=0.2$,
$A=3$) and vanishing only at the $D\to0$ (equivalently $s=0$) Perron anchor where
$g_A\to1$, $A\in\{3,4,5\}$. A uniform-in-$A$ \emph{closed-form} cone witness (the
explicit once-iterated dressed pair, as written out for binary in Lemma 7.34e)
remains open.

\emph{The residual.} Were the convexity of $u\mapsto g_A(\arccos u)$ proven,
Theorem 7.34m's achievability would collapse to this single endpoint,
\[
  g_A(s)\le1-\tilde c_A(p,D)\,D(1-D)(1-\cos s)\ \text{ on }(0,\pi],\qquad
  \tilde c_A(p,D):=R_A(\pi)=\frac{1-g_A(\pi)}{2D(1-D)}>0,
\]
with $\tilde c_A$ explicit (an algebraic number of degree $4$ over
$\mathbb Q(p,D)$ for $A\ge3$), eliminating the $s$-Sturm; the convexity is then
the sole residual, exactly as in the binary case. The endpoint floor is uniform
and \emph{comfortable}: $\inf_{\text{Gray}}\tilde c_A(p,D)\approx\tfrac32$,
attained in the small-$p$, $D\to D_c$ corner and \emph{increasing} mildly in $A$
($1.5000$ at $A=2$, $1.5049$ at $A=3$, rising to $1.5085$ at $A=8$) --- well
above the conservative alphabet-uniform $\tfrac{11}{16}$ the achievability
machinery actually consumes (Lemma 7.34l(iii)). Thus the A-ary residual is the
\emph{same shape} as the binary one: a single endpoint convexity, with the
endpoint value now a degree-$4$ algebraic number and $g_A(\pi)<1$ cone-certified
(numerically here; the explicit uniform-in-$A$ witness, as in Lemma 7.34e,
open).
(Scope note, post hoc: the interior convexity --- and with it the $s=\pi$
binding --- turns out to be \emph{regional}, failing for $A\ge16$ within $1\%$
of the uniform corner $p=\tfrac{A-1}A$ while the floor target itself survives
there by a wide direct margin; Remark 7.34m$'''''$.)
\emph{Validated} in
\texttt{scripts/verify/probe\_7\_34m\_endpoint\_pi.py} (M1 the closed forms and
$A=2$ binary reduction; M2 the linear$\times$irreducible-quartic split; M3 the
positive dominant root and $g_A(\pi)=|C_\pi/C_0|^2\rho$; M4 the $\approx\tfrac32$
endpoint floor; M5 the signed-cone $g_A(\pi)<1$; M6 the $s=\pi$ binding /
convexity --- all PASS).

\textbf{Remark 7.34m$''$ (the convexity residual at leading order in $D$: an
affine-at-$D{=}0$ law, a cubic dispersion correction, and full uniform-in-$u$
convexity).} The convexity conjecture left open by Remarks 7.34m$'$ and
7.34e$'$ --- that $G(u):=g_A(\arccos u)$ is convex on $[-1,1]$ (so that the floor
ratio $R_A$ is minimised at $s=\pi$) --- is here \emph{resolved at leading order
in $D$}. We first reduce it at its \emph{tightest} point $u\to1$ (the $s\to0$
curvature anchor, where $\min_u G''$ sits) to a scalar Taylor-coefficient
inequality, then show the leading coefficient is uniform across all $u$.

\emph{The reduction (exact).} Write the even expansion $g_A(s)=1-v\,s^2
+g_4\,s^4+O(s^6)$, with $v$ the per-site posterior variance of Lemma 7.34l(iii)
and $g_4$ the quartic coefficient. The substitution $u=\cos s$ (using
$1-\cos s=\tfrac{s^2}2-\tfrac{s^4}{24}+\cdots$) gives, as $u\to1$,
\[
  G''(1^-)=8g_4-\tfrac23 v,\qquad\text{so}\qquad
  G''(1^-)\ge0\iff g_4\ge\tfrac{v}{12}.
\]

\emph{The structure (the new content).} At $D=0$ the disagreement tilt is
$\eta(s)=-D(1-e^{is})/(A-1)+O(D^2)=O(D)$, whence $g_A(s)=1-2v(1-\cos s)+O(D^3)$:
$G(u)$ is \emph{affine} in $u$ at leading order ($G''\equiv0$), i.e.\
$g_4=\tfrac{v}{12}$ \emph{exactly} at $D=0$. The binding-endpoint convexity is
therefore a pure \emph{cubic-in-$D$} effect, and the inequality $g_4\ge v/12$ is
asymptotically tight. Its leading coefficient
$c_A(p):=\lim_{D\to0}(g_4-v/12)/D^3$ is positive for every $A$, and at $A=2$ it
coincides, \emph{in closed form}, with the second-order dispersion constant of
Lemma 7.34l(iii):
\[
  c_2(p)=\kappa_2^2=\frac{(1-2p)^2}{p^2(1-p)^2},\qquad
  G''(1^-)=8c_2\,D^3+O(D^4)>0 .
\]
Thus the binding-endpoint convexity holds at leading order in $D$, governed by
the positive cubic constant $c_A$. This $A=2$ identity $c_A=\kappa_A^2$ is,
however, a \emph{closed-form coincidence that does not persist to $A\ge3$}: $c_A$
is a genuine \emph{third}-order spectral object --- the $\cos2s$ harmonic of the
$O(D^3)$ term of $\rho(W_A)$, the $O(D^2)$ term carrying \emph{zero} second
harmonic (the double cancellation) --- whereas the dispersion constant
$\kappa_A^2$ (which sets the second-order dispersion
$\bar v_A=D(1-D)[1-\kappa_A^2 D(1-D)/(1-2D)^2]$) is a \emph{second}-order
(variance) object; the two agree at $A=2$ only through the binary closed-form
Perron quadratic. For $A\ge3$ they differ: with two independent high-precision
extractors (finite-difference Richardson and Chebyshev quadrature, in agreement)
the ratio $c_A/\kappa_A^2$ at $p=0.2$ is $1.0000,\,0.984,\,1.116,\,1.235$ for
$A=2,3,4,5$ (not even sign-definite), so the leading curvature is set by the
cubic constant $c_A$, \emph{not} by the dispersion constant. The double
cancellation that makes $g_4-v/12=O(D^3)$ (agreement of $g_4$ and $v/12$ through
$O(D^2)$) is the \emph{net} of the $O(D^2)$ near-cancelling terms --- the concave
curvature $\langle l|M''|r\rangle$ against the positive level-repulsion
$2\langle l|M'SM'|r\rangle$ of Remark 7.34e$'$ --- which cancel to one higher
order; directly, $G''(1^-)=8(g_4-v/12)\sim D^3$ (numerically scaling order
$2.99$).

\emph{Uniform in $u$ (the full-interval upgrade).} The leading coefficient is, in
fact, \emph{independent of $u$}: $G''(u)=8c_A D^3+O(D^4)$ for \emph{all}
$u\in[-1,1]$, not only at the endpoint. The $u$-dependence is pushed to $O(D^4)$:
the $u$-spread of $G''(u)/D^3$ shrinks as $O(D)$ (e.g.\ at $A=2,p=0.2$ it is
$2.0\to0.53\to0.13$ as $D=2{\times}10^{-3}\to5{\times}10^{-4}\to1.25{\times}10^{-4}$),
its common value matching $8c_A$ (at $A=2$, $112.41$ vs $8\kappa_2^2=112.5$, the
identity point). Equivalently, the only leading non-affine content of
$g_A(s)=G(\cos s)$ is the \emph{second harmonic} $\cos2s=2u^2-1$ --- a pure $u^2$
curvature, hence a \emph{constant} $G''$ --- with positive coefficient
$4c_A D^3$; the
affine ($T_0,T_1$) part is the $D=0$ law $1-2v(1-u)$.

\emph{The prefactor carries no curvature (the convexity is purely spectral).}
The contour prefactor is \emph{exactly affine} in $u$: since $C_s$ is linear in
$e^{is}$ for every $A$ (the dual identity gives $C_s/C_0$ linear in $E=E_0e^{is}$),
$|C_s/C_0|^2$ is affine in $u=\cos s$; at $A=2$ in closed form
\[
  |C_s/C_0|^2=1+\frac{2D^2(1-D)^2}{(1-2D)^2}\,(1-\cos s)
\]
(verified: its Chebyshev modes $k\ge2$ vanish to $10^{-16}$, $A=3,4,5$). Hence
$|C_s/C_0|^2$ contributes \emph{zero} curvature, and with $G=|C_s/C_0|^2\rho$,
$\rho=\rho(W_A)$,
\[
  G''(u)=|C_s/C_0|^2\,\rho''(u)+2\,(|C_s/C_0|^2)'\,\rho'(u),
\]
so the convexity of $G$ --- the sole achievability residual --- is \emph{purely}
the spectral convexity of the replica Perron eigenvalue $\rho(W_A)$, modulated by
the known affine prefactor; the curvature constant $4c_A D^3$ above is
$\rho$'s.

\emph{Asymmetric binary (the residual of Remark 7.34e$'$ too).} The whole
analysis transfers verbatim to the asymmetric binary instance (source
$P(0\!\to\!1)=a$, $P(1\!\to\!0)=b$, $a\ne b$, BSC$(D)$ channel; the $8\times8$
operator $W_8$ of Lemma 7.34e): $g(s)$ is even there as well, the reduction
$G''(1^-)=8g_4-\tfrac23 v$ is unchanged, $g_4\ge v/12$ holds throughout the
asymmetric Gray region $0<D\le D_c(a,b)$ (high precision), $G''(1^-)\sim D^3$
(scaling order $2.99$ at $a=0.05,b=0.3$ and $a=0.1,b=0.2$ --- not $D^2$; the
asymmetry does \emph{not} cost an order), and the leading coefficient of
$G''(u)$ is again $u$-independent (the $u$-spread of $G''/D^3$ shrinks $\propto D$,
common value $\approx570.5$ at $a=0.05,b=0.3$). So the binding-endpoint
convexity, the uniform-in-$u$ full convexity at leading order, and hence the
unconditional-for-small-$D$ achievability all hold for the asymmetric instance
too --- closing, at leading order in $D$, the convexity residual of \emph{both}
Remark 7.34m$'$ (A-ary symmetric, Theorem 7.34m) and Remark 7.34e$'$ (binary
asymmetric, Theorem 7.34i).

\emph{Status.} Thus the \emph{full} convexity of $G$ on $[-1,1]$ holds at leading
order in $D$, \emph{uniformly in $u$}, with the positive constant $8c_A$
($=8\kappa_A^2$ only at $A=2$): for sufficiently small $D$, $G''(u)>0$ for every $u$, so
$s=\pi$ is the binding point. At such $D$ the other two inputs of Theorem 7.34m's
achievability are automatic from the same affine law $g_A(s)=1-2v(1-\cos s)
+O(D^3)$ with $v=\bar v_A>0$ --- the spectral gap $\sup_{s\in[\varepsilon,\pi]}
g_A(s)<1$ (since $1-\cos s>0$) and the positive curvature --- so all three hold
and the achievability is \emph{unconditional for sufficiently small $D$}.
Numerically the convexity ---
hence the $s=\pi$ binding, hence the whole achievability --- holds on the
\emph{entire} Gray region for every tested $A$. The residual is therefore no
longer ``prove convexity'' (an opaque spectral inequality) but the narrower
quantitative passage from ``sufficiently small $D$'' to all $D\in(0,D_c^{(A)}]$:
a bound on the $O(D^4)$ remainder showing it cannot overturn the $8c_A D^3$
leading term before $D_c$ --- which, since $D_c^{(A)}=O(p^2)$ is itself small, the
numerics confirm throughout.

\emph{All-$D$ via a Chebyshev tail (the bulk closed).} This remainder bound has
an explicit form. Expanding $g_A(s)=\sum_{k\ge0}a_k(D)\,T_k(\cos s)$ in
Chebyshev modes, the coefficients decay as $a_k=O(D^{k+1})$ for $k\ge2$ with
\emph{alternating signs} and a geometric ratio $|a_{k+1}/a_k|\approx0.9\,D$
(numerically; $a_1=O(D)$ is the affine mode, $a_2=2c_A D^3+\cdots>0$ the
curvature). Since $T_0''=T_1''=0$, $T_2''\equiv4$, and $|T_k''(u)|\le
T_k''(1)=\tfrac13 k^2(k^2-1)$ on $[-1,1]$,
\[
  G''(u)\;\ge\;\mathrm{LB}(D):=4a_2-\sum_{k\ge3}|a_k|\,\tfrac13 k^2(k^2-1),
\]
a \emph{$u$-independent} lower bound. $\mathrm{LB}(D)>0$ throughout the Gray
region for $p$ in the \emph{memory regime} (verified for $p\lesssim0.3$ at $A=2$,
$p\lesssim0.4$ at $A=3$; tail/$4a_2$ ratio $0.001$--$0.21$), \emph{proving the
full all-$D$ convexity there} (modulo making the $O(D^{k+1})$ coefficient bounds
rigorous --- discharged in closed form for $A=2$ below, $A\ge3$ reducing to the
degree-2 coefficient recursion). The analyticity route is concrete: $g(s)$'s nearest singularity is the
Perron branch point of $W(s)$, numerically at $|\mathrm{Im}\,s|=\tau^*\approx
\ln(1/D)$ (the top-two eigenvalue gap stays $>0.1$ up to $\tau=5\approx\ln(1/D)$
at $p=0.2$), and $\tau^*\ge\ln(1/D)-O(1)$ follows \emph{rigorously} from
gap-stability (Bauer--Fike) of the frozen operator $W_0$ (gap $\Delta_0$) against
$\|W(s)-W_0\|\le B\,D\,e^{|\mathrm{Im}\,s|}$. The one subtlety is that the bare
Bernstein bound $|a_k|\le2M\,e^{-k\tau^*}=O((cD)^k)$ is one $D$-power looser than
the observed $O(D^{k+1})$ decay (the harmonic cancellation), so a sharp tail
closure must also capture that extra cancellation --- a genuine
several-step analyticity task, not a one-shot. It
degrades only as $p\to\tfrac{A-1}{A}$ (the near-uniform edge, where $D_c$ is
largest and the geometric decay slows), where the crude endpoint bound
$T_k''(1)$ over-counts and a sharper interior estimate is needed --- though for moderate alphabets ($A\le12$) the true $G''(u)$ stays positive
there (the alternating $a_k$ cancel); for $A\ge16$ deep in the near-uniform
pocket the convexity itself fails, and the target is held by direct margin
instead (Remark 7.34m$'''''$). Since the
near-uniform point is the memoryless limit ($V_{\mathrm{lossless}}\to0$), the
all-$D$ closure already covers the entire memory regime of interest.

\emph{The $A=2$ case in closed form (the cancellation made rigorous).} For $A=2$
the residual is dischargeable in closed form. The $4\times4$ pattern-quotient
operator's characteristic polynomial factors (the $W_8\to W_4$ reduction), and the
Perron branch is the larger root of a \emph{quadratic},
\[
  \rho(s)=\tfrac12\bigl[\,B+\sqrt{B^2+4\kappa^2\,\eta\bar\eta}\,\bigr],\qquad
  B=(1+\eta)(1+\bar\eta),\quad \kappa^2=\tfrac{(1-2p)^2}{p^2(1-p)^2}
\]
(checked $<10^{-30}$ against the matrix Perron eigenvalue throughout the Gray
region). Both $B$ and $W:=\eta\bar\eta$ are invariant under $s\to-s$ (the
$z\to1/z$ symmetry), hence real functions of $u=\cos s$: with $a_0=D-1$,
$\gamma=D^2/(1-D)$, $\Delta(u)=a_0^2+\gamma^2+2a_0\gamma u$,
\[
  \eta+\bar\eta=\frac{2D(a_0-\gamma)(1-u)}{\Delta(u)},\qquad
  \eta\bar\eta=\frac{2D^2(1-u)}{\Delta(u)} .
\]
The discriminant is a \emph{perfect square plus an $O(D^2)$ deviation}:
$f:=B^2+4\kappa^2 W=P^2+R$ with the affine root $P=1+2\alpha(1-u)$
($\alpha=D/(D-1)$) and $R=4(1{+}2\kappa^2)\alpha^2(1-u)+O(D^3)=O(D^2)(1-u)$
uniformly in $u$ (the worst point $u\to1$ is finite, with leading constant
$4+8\kappa^2$; $B,W$ are affine in $u$ \emph{only through $O(D^2)$}, their
non-affine modes themselves $O(D^{k+1})$). Hence
$\sqrt f=P+R/(2P)-R^2/(8P^3)+\cdots$ is affine through $O(D^2)$, and its non-affine
Chebyshev harmonics inherit the deviation order times the geometric per-harmonic
ratio $2|\alpha|+O(D)=2D/(1-D)+O(D)<1$, giving $|a_k|\le 8\kappa^2 D^{k+1}$ for
$k\ge2$. This is the rigorous form of the cancellation the bare Bernstein
majorant misses --- $f$ being a perfect square through $O(D^2)$ is exactly why
$\sqrt f$ is affine through $O(D^2)$ and its non-affine tail is one $D$-power
below $O(D^k)$. The uniform ratio then closes the tail by a \emph{convergent
geometric majorant} (not a truncation),
$\sum_{k\ge3}|a_k|\tfrac13 k^2(k^2-1)\le|a_3|\sum_{k\ge3}r_0^{\,k-3}\tfrac13
k^2(k^2-1)$ with $r_0=2|\alpha|+O(D)<1$, so $\mathrm{LB}(D)>0$ throughout
$(0,D_c]$ on the memory regime ($p\lesssim0.3$; the surviving fraction
$\mathrm{LB}/4a_2$ runs from $0.99$ at small $p$ to $0.84$ at $p=0.3$). Thus the
all-$D$ convexity --- hence the $s=\pi$ binding, hence the achievability --- is
closed in closed form for $A=2$ on the memory regime. The exact Perron root and
the perfect-square decomposition are symbolic identities; the explicit
uniform-in-$u$ remainder bound (the analyticity/Bernstein step) is the narrow
elementary task that upgrades the certified geometric majorant to a one-line
theorem. For $A=2$ the achievability was already unconditional via the GAP-1
floor route of Remark 7.34b; this discharges the distinct \emph{convexity}
route's residual at $A=2$ and exhibits the cancellation mechanism. For $A\ge3$
the Perron branch is a root of an irreducible \emph{degree-five} polynomial (no
single radical), so the closed-form route does not transfer; but the
characteristic polynomial's coefficients are trigonometric polynomials of degree
$\le2$ in $s$ ($\eta,\bar\eta$ each inject a single harmonic), so the
$a_k=O(D^{k+1})$ bound reduces to a cancellation-aware coefficient recursion from
that degree-2 characteristic equation, with the $A=2$ perfect square as its
template --- the remaining open content of the convexity capstone.

\emph{Validated} in
\texttt{scripts/verify/lemma\_7\_34m\_ak\_bound\_A2.py} (C1 the exact Perron
quadratic; C2 the perfect-square deviation $R=O(D^2)(1-u)$; C3 the geometric
majorant $|a_k|\le8\kappa^2 D^{k+1}$; C6 the all-$k$ closure via the convergent
geometric majorant; C4 $\mathrm{LB}(D)>0$ on $(0,D_c]$; C5 the $O(D^4)$ tail ---
all PASS) and in
\texttt{scripts/verify/probe\_7\_34m\_convexity\_endpoint.py} (E1 the reduction
$G''(1^-)=8g_4-\tfrac23 v$; E2 $g_4\ge v/12$ on the Gray region; E3 the
affine-at-$D{=}0$ law; E4 $c_A>0$ for all $A$ and $c_2=\kappa_2^2$ in closed
form; E5 the assembled $G''(1^-)=8\kappa_2^2 D^3>0$; E6 the uniform-in-$u$
$G''(u)=8c_A D^3+O(D^4)$; E7 the asymmetric-binary transfer; E8 the
Chebyshev-tail all-$D$ closure on the memory bulk --- all PASS).

\emph{The representation-free engine (general $A$): boundary affineness.} The
perfect square above is the $A=2$ shadow of a structure that holds for \emph{every}
$A$ and needs no closed-form root. At $\bar\eta=0$ the replica operator
$Q_A(\eta,0)$ has \emph{rank two}, trace $1+(A-1)\eta$, and all other eigenvalues
zero, so its characteristic polynomial is $\lambda^{\,n-1}\bigl(\lambda-(1+(A-1)
\eta)\bigr)$ and its \emph{unique nonzero} --- hence Perron --- eigenvalue is
exactly affine:
\[
  \rho(\eta,0)=1+(A-1)\eta,\qquad \rho(0,\bar\eta)=1+(A-1)\bar\eta .
\]
This is a one-line trace/rank identity (the unique nonzero eigenvalue of a
trace-$\tau$ operator whose every other eigenvalue vanishes is $\tau$), uniform in
$A$ and $p$ --- verified symbolically for $A=2,\dots,6$ (at $A=2$ it is the
boundary trace $\tfrac12[(1+\eta)+\sqrt{(1+\eta)^2}\,]=1+\eta$ of the perfect
square). Boundary affineness forces the \emph{pure}-$\eta^k$ and pure-$\bar\eta^k$
coefficients of $\rho(\eta,\bar\eta)$ to vanish for $k\ge2$, so the only non-affine
content is the mixed coupling; since $\eta$'s harmonic-$m$ content is $O(D^{m+1})$
(the leading injection is $\alpha(1-z)$, $\alpha=O(D)$, $z=e^{is}$; higher
harmonics are more suppressed) and likewise for $\bar\eta$ in $z^{-1}$, the $k$-th
harmonic of $\rho$ is assembled from mixed monomials $\eta^a\bar\eta^b$ with $a\ge
k,\ b\ge1$, the minimal $\eta^k\bar\eta$ being $O(D^{k+1})$. Hence the \emph{decay
order} $a_k=O(D^{k+1})$ ($k\ge2$) is established for \emph{all} $A$ (the mixed
coefficients are finite). What remains regime-dependent is only the
\emph{quantitative closure} $\mathrm{LB}(D)>0$: for $A=2$ the explicit
$|a_k|\le8\kappa^2 D^{k+1}$ and the ratio $2|\alpha|<1$ certify it on the whole
Gray region (above); for $A\ge3$ the convexity $\mathrm{LB}(D)>0$ is certified
numerically with an even \emph{larger} margin ($\mathrm{LB}/4a_2\approx0.96$--$0.98$,
vs.\ $0.84$ at $A=2$, since $|\eta|$ shrinks with $A$), and rests on the curvature
constant $c_A$ being positive. Crucially the perturbation is \emph{benign}: the
frozen operator $Q_A(0,0)$ has spectral gap $1$ (Perron eigenvalue $1$, every other
eigenvalue $0$) and the physical $|\eta|\le0.03$ shrinks with $A$ ($0.032$ at $A=2$
down to $0.001$ at $A=6$), so neither convergence nor the geometric tail is the
obstruction --- the large mixed coefficients (e.g.\ $c_{2,2}\approx-1.1\times10^4$
at $A=3$) are offset by the tiny $\eta^{a+b}$. The convexity capstone's residual is
thereby reduced, from ``prove $a_k=O(D^{k+1})$'' (now done, all $A$, via boundary
affineness), to the \emph{single positivity} $c_A>0$ (the second-harmonic curvature
$=$ dispersion constant; $c_2=\kappa_2^2>0$ in closed form, $c_A>0$ verified and
growing --- $26.4,42.5,60.4$ at $A=3,4,5$) --- a sharply narrower analytic task.
\emph{Validated} in
\texttt{scripts/verify/probe\_7\_34m\_boundary\_affineness.py} (A1 the rank-two/
trace identity and $\rho(\eta,0)=1+(A-1)\eta$, symbolic $p$, $A=2,\dots,6$; A2 its
Perron slope $A{-}1$ at all $p$; A3 the vanishing pure-$\eta^k$ coefficients; A4
the resulting $a_k=O(D^{k+1})$ order; A5 the LB-closure diagnostic ($Q_0$ gap $1$,
$c_A>0$, $\mathrm{LB}/4a_2\gtrsim0.96$, $A=3,4,5$) --- all PASS).

\textbf{Remark 7.34m$'''$ (the operational RD-dispersion residual ledger:
converse vs.\ achievability, sub-question by sub-question).} The §7.34 grid is a
\emph{converse} that is unconditional and fully general, plus an
\emph{achievability} that is closed in two stated regimes (BSMS on its whole
Gray region; general-A at leading order in $D$ and on the memory bulk) with a
single, precisely-localised 2nd-order residual --- the rigorous
cancellation-aware $a_k$ bound --- and the beyond-Gray exact value still open.
This is a \emph{how-fast-to-$nR(D)$} (2nd-order) residual; the 1st-order rate
$R(D)$ and the converse dispersion bound do \emph{not} depend on it. The
following ledger records, per sub-question, its status, paper location, the
named obstacle (if OPEN) or proof input (if CLOSED), and the witnessing verify
script.

\begingroup\footnotesize\sloppy
\setlength{\tabcolsep}{4pt}
\noindent\begin{tabular}{@{}p{3.0cm}p{1.5cm}p{2.6cm}p{2.9cm}p{3.0cm}@{}}
\hline
\textbf{sub-question} & \textbf{status} & \textbf{paper location} &
\textbf{obstacle (OPEN) / proof input (CLOSED)} & \textbf{witnessing script}\\
\hline
\emph{Converse} $V_{\mathrm{op}}\ge V_{\mathrm{lossless}}$ --- all source types,
all balanced / group-difference distortions on each distortion's own Gray region
& \textbf{CLOSED} (unconditional) & Rem 7.34d/n/q; Thm 7.34i/m, Thm 7.34n
converse & source-agnostic dual identity $K_\nu K^{-1}=\gamma I$ scalar
$\Rightarrow j_n=i_n-nc$ deterministic, $\rho=1$, $i_n$-CLT $+$ KV $d$-tilted
one-shot & \texttt{probe\_7\_34n\_aary\_asym} (N5), \texttt{probe\_7\_34q\_lee\_general\_distortion} (Q1--Q4), \texttt{remark\_7\_34d\_general\_converse\_discrete\_markov}\\
\hline
\emph{Achievability (a1)} --- BSMS (binary symmetric, Hamming), whole closed Gray
region $0<D\le D_c(p)$ & \textbf{CLOSED} & Rem 7.34b concluding theorem; Lem
7.34e (asymmetric binary) & full proof: replica spectral inequality
$g(s)\le1-\tfrac{11}{16}D(1-D)(1-\cos s)$, quenched tail, GAP-1 steepest descent,
M\"obius-orbit $D=D_c$ endpoint & \texttt{probe\_7\_34\_gap1\_writeout\_v2}, \texttt{probe\_7\_34\_dc\_endpoint}, \texttt{lemma\_7\_34e\_spi\_closed\_form}\\
\hline
\emph{Achievability (a2)} --- general A-ary / asymmetric, leading order in $D$,
\emph{uniform in $u$} & \textbf{CLOSED} (leading order) & Rem 7.34m$''$
(reduction, structure, uniform-in-$u$) & affine-at-$D{=}0$ law
$g_A=1-2v(1-\cos s)+O(D^3)$; second-harmonic curvature $G''(u)=8c_A
D^3+O(D^4)>0$ ($c_A=\kappa_A^2$ only at $A=2$); prefactor $|C_s/C_0|^2$ affine in $u$ (no curvature) & \texttt{probe\_7\_34m\_convexity\_endpoint} (E1--E7), \texttt{probe\_7\_34m\_endpoint\_pi}\\
\hline
\emph{Achievability (a3)} --- general A-ary, \emph{all} $D\le D_c$ on the
\emph{memory bulk} & \textbf{CLOSED} ($A=2$; $\mathrm{LB}>0$ certified) /
\textbf{$a_k=O(D^{k+1})$ order proven all-$A$, $\mathrm{LB}$-closure modulo the
explicit coefficient bound} ($A\ge3$) & Rem 7.34m$''$ (``$A=2$ in closed form''
/ ``boundary affineness'') & boundary affineness $\rho(\eta,0)=1+(A-1)\eta$ (rank-
two/trace, all $A$) $\Rightarrow$ pure-$\eta^k$ vanish $\Rightarrow$ the order
$|a_k|=O(D^{k+1})$ (all $A$). $A=2$: explicit $|a_k|\le8\kappa^2D^{k+1}$ $+$ ratio
$2|\alpha|<1$ $\Rightarrow\mathrm{LB}(D)>0$. $A\ge3$ \emph{obstacle:} the mixed
coefficients are large; an explicit uniform bound is needed to carry the majorant
to $D_c$ &
\texttt{lemma\_7\_34m\_ak\_bound\_A2} (C1--C6), \texttt{probe\_7\_34m\_boundary\_affineness} (A1--A4), \texttt{bug\_009\_ak\_decay\_convexity} (B1--B5)\\
\hline
\emph{Achievability (a4)} --- near-uniform / memoryless edge $p\to\tfrac{A-1}A$
& \textbf{OPEN} & Rem 7.34m$''$ (status paragraph); Rem 7.34m$'''''$ (the
pocket) & \emph{obstacle:} the crude endpoint bound $T_k''(1)$ over-counts
where the geometric decay slows; and deep in the pocket ($p\gtrsim0.98\,
\tfrac{A-1}A$, $A\ge16$) the convexity itself fails --- there the target holds
by direct margin $R_A\ge1.74$. Outside the memory regime: this is the
$V_{\mathrm{lossless}}\to0$ limit & \texttt{probe\_7\_34m\_convexity\_endpoint}
(E8), \texttt{probe\_7\_34\_corner\_pocket\_refutation} (H1--H4)\\
\hline
\emph{Sharp replica floor} $R_A\ge\tfrac32$ (strengthens the consumed
$\tfrac{11}{16}$; not a hypothesis of any row above) & \textbf{CLOSED} ($A=2$,
sharp, closed form) / \textbf{regional certificates} ($A\ge3$) & upgraded
Rem 7.34b theorem (steps (v)--(vi)); Rem 7.34m$''''$/7.34m$'''''$ & $A=2$: the
monotone discriminant $4C-B^2=8am(1-(3+\kappa^2)a)$ $+$ the endpoint
$P(a,p)\ge0$ via $f$-monotonicity. $A\ge3$: effective endpoint theorem
($A=3..6$, all $p\le p_0$) $+$ the universal coefficient ladder; Schur
certificate to $p\le0.40$--$0.45\,\tfrac{A-1}A$ (discrete $p$); pocket direct
margin. \emph{obstacle:} the mid-band certification program (real-rootedness
lemma $+$ five sign certificates) &
\texttt{probe\_7\_34m\_monotone\_a2\_proof}, \texttt{probe\_7\_34m\_endpoint\_sharp\_a2}, \texttt{lemma\_7\_34\_endpoint\_effective\_bound}, \texttt{lemma\_7\_34\_schur\_smallp\_certificate}, \texttt{probe\_7\_34\_midband\_charpoly\_route}\\
\hline
\emph{Exact $V_{\mathrm{op}}$, beyond Gray} $D>D_c$ & \textbf{OPEN} & Rem 7.34o &
\emph{obstacle:} needs the optimal non-SLB-tight, memory-aware test channel
(optimal output of growing Markov order, Eswaran--Gastpar 2022); $V_{\mathrm{op}}$
flat $=V_{\mathrm{lossless}}$ on $[0,D_c]$, decreasing below for $D>D_c$,
continuous at $D_c$ (finite-$n$ map established; exact value open) & \texttt{probe\_7\_34o\_beyond\_gray} (O1--O3), \texttt{probe\_7\_34\_offgray\_rho}\\
\hline
\end{tabular}
\par\endgroup

\noindent Thus the headline is neither a flat ``achievability open'' nor a flat
``closed'': the converse is unconditional and fully general; the BSMS Gray
achievability is fully proven; the general-A achievability is closed at leading
order in $D$ and (modulo the $a_k$ bound) on the memory bulk; and exactly two
items remain genuinely open --- the rigorous cancellation-aware $a_k=O(D^{k+1})$
bound (closing row a3 all the way and removing the ``modulo'') and the
beyond-Gray exact $V_{\mathrm{op}}(D)$. The $a_k$-decay order itself (the
$O(D^{k+1})$ scaling, alternating signs, geometric ratio $\approx0.9D$) is
\emph{observed and reproducible} (\texttt{bug\_009\_ak\_decay\_convexity}); what
stays open is the rigorous \emph{bound}, not the measured decay.
In parallel, the sharp-constant campaign of Remarks 7.34m$''''$/7.34m$'''''$
upgrades the \emph{binary} replica floor to its sharp $\tfrac32$ in closed form
and assembles regional certificates of $R_A\ge\tfrac32$ for $A\ge3$ (corner:
effective theorem; small/moderate $p$: Schur, discrete $p$; near-uniform
pocket: direct margin --- where the monotonicity conjectures themselves
regionally fail), leaving the mid-band certification program as the named
residual on that axis. No row above \emph{consumes} the sharp constant --- the
achievability machinery runs on the $\tfrac{11}{16}$ floor --- so the campaign
strengthens constants and scope-corrects conjectures without touching the
ledger's hypotheses.

\textbf{Remark 7.34m$''''$ (the sharp-floor campaign I: the monotone-reduction,
order-by-order $s$-monotonicity, and the universal endpoint coefficient
ladder).} The sharp $A=2$ constant (the upgraded replica theorem of Remark
7.34b: $R_2(s)\ge\tfrac32$, all $s$, all $D\le D_c$, closed form) makes the
general-$A$ target precise:
\[
  R_A(s;D)\;\ge\;\tfrac32\qquad\text{on the Gray region, for every }A
\]
--- the sharp form of the $\approx\tfrac32$ binding-corner limit of Lemma
7.34l(iii), strictly stronger than the conservative $\tfrac{11}{16}$ the
achievability machinery consumes. The route that closed $A=2$ is a
\emph{monotone-reduction}: split the target into (M) $R_A(\cdot\,;D)$
non-increasing in $s$, and (E) the single endpoint $R_A(\pi)\ge\tfrac32$ ---
no convexity of the Perron branch, no joint $(s,D)$ bound. At $A=2$ both
halves are proven in closed form (steps (v)--(vi) of the theorem). For
$A\ge3$ the two halves stand as follows.

\emph{(M) at its leading orders --- proven for every alphabet.} Expanding
$1-g_A(s;D)=e_1(t,p)\,D+e_2(t,p)\,D^2+e_3(t,p)\,D^3+O(D^4)$, $t=1-\cos s$:
\emph{exactly} $e_1=2t$ and $e_2=t\,q_A(p)$ for every $A$ (symbolic;
$q_A(p)=-2(A-1)/p^2+(8A-12)/p+O(1)$ at small $p$), so
\[
  R_A(s;D)=2+\bigl[2+q_A(p)\bigr]D+O(D^2)
\]
is \emph{exactly flat in $s$ through $O(D)$} --- an honest flatness lemma: at
this order the whole-interval inequality carries no information beyond the
endpoint (any single $s$ determines the $O(D)$ level), and genuine
$s$-dependence starts at $O(D^2)$. There the $t$-ladder continues,
$e_3=\alpha_A(p)\,t+\beta_A(p)\,t^2$ exactly, whence
$dR_A/dt=\beta_A(p)\,D^2/(1-D)+O(D^3)$, and the sign is settled
\emph{universally}:
\[
  \beta(A,p)=-\frac{4\,\bigl(1-A(1-p)\bigr)^2\,P_3(A,p)}{p^2\,(A-1)^5\,(1-p)^2}
  \;\le\;0,
\]
with $P_3(A,p)=2A^4(1-p)^4-A^3(4p^4-20p^3+39p^2-32p+9)-A^2(8p^3-33p^2+40p-15)
-A(6p^2-16p+11)+3$ (at $A=2$, $P_3\equiv1$, recovering the manifestly
nonpositive $\beta_2=-4(1-2p)^2/(p^2(1-p)^2)$). ($\beta$'s closed form: a symbolic-$A$ series-ring derivation, machine-anchored
exactly at $A=2,\dots,6$ including the out-of-sample $A=6$.) The positivity $P_3\ge0$ on
$\{A\ge2,\ 0<p<\tfrac{A-1}A\}$ is \emph{proven} by a two-region P\'olya-shift
certificate: on $p\in(0,\tfrac12]$ every coefficient of $P_3(2+B,p)$ as a
polynomial in $B\ge0$ is Sturm-root-free and positive on the interval; on
$p\in[\tfrac12,1)$ likewise for $(1-p)^4P_3(1/(1-p)+B,p)$, with the
domain-boundary value $2p^3(1-p)^2\ge0$. Hence $R_A$ is monotone
non-increasing in $s$ at its first nontrivial order $O(D^2)$, for \emph{every}
alphabet. (All-orders-in-$D$ monotonicity is \emph{regional}, not global ---
the corner pocket of Remark 7.34m$'''''$.)

\emph{(E) the universal endpoint coefficient ladder.} At the binding corner
$D=D_c^{(A)}(p)$ the threshold and the endpoint margin expand with exact
rational coefficients,
\[
  D_c^{(A)}=\frac{p^2}{4(A-1)}\Bigl(1+2p+\delta_2p^2+\delta_3p^3+\delta_4p^4
  +O(p^5)\Bigr),\qquad
  R_A(\pi)\Big|_{D_c}-\tfrac32=c_1p+c_2p^2+c_3p^3+c_4p^4+O(p^5),
\]
\[
  \delta_2=\frac{12A^2-28A+21}{4(A-1)^2},\quad
  \delta_3=\frac{8A^3-32A^2+53A-32}{2(A-1)^3},\quad
  \delta_4=\frac{40A^4-240A^3+652A^2-848A+429}{8(A-1)^4},
\]
\[
  c_1=\frac{A-2}{A-1},\quad
  c_2=-\frac{4A^2+24A-65}{8(A-1)^2},\quad
  c_3=-\frac{18A^2-153A+232}{8(A-1)^3},\quad
  c_4=-\frac{60A^3-890A^2+3171A-3272}{32(A-1)^4}
\]
($A=2$ anchors from the closed forms: $\delta_2=\tfrac{13}4$, $\delta_3=5$,
$\delta_4=\tfrac{61}8$; $c_1=0$, $c_2=\tfrac18$, $c_3=\tfrac14$,
$c_4=\tfrac5{16}$ --- every coefficient on-family). The mechanisms: the
$\delta_k$ come from the discriminant of the $3\times3$ symmetry-reduced
alternating-word cubic ($D_c$ is where its class-$0/1$ eigenvalue pair turns
complex --- the collision mechanism, Remark 7.34o$'$), the ansatz
$D=Wp^2(1+d_1p+\cdots)$ series-expanded in exact arithmetic; the $c_k$ from a
rank-one Rayleigh--Schr\"odinger ladder at the degenerate $D=0$ quotient
($Q(0)=c\,e_0^T$ has $\lambda=1$ simple, so perturbation theory collapses to
pure matrix algebra --- no determinants, no root-solving of the dead $A\ge3$
endpoint quartic), composed with the exact $\eta_\pi(D)$, $|C_\pi/C_0|^2(D)$
of Remark 7.34m$'$ and the $D_c$ series; the even-graded poles
$\lambda_k\sim p^{-2\lfloor k/2\rfloor}$ put ladder term $k$ into the margin at
order $p^{2\lceil k/2\rceil-2}$, so $c_2,c_3$ are carried entirely by
$k\le4$. The leading coefficient splits as $c_1=-(\varepsilon_1+\delta_1)/2$:
$\delta_1=2$ (the threshold's relative $p$-slope, from the discriminant's
$p^{12}/p^{13}$ coefficients) and $\varepsilon_1=-2-2(A-2)/(A-1)$ (the
endpoint's $D$-slope, $R_A(\pi;D)=2-E_A(p)D+O(D^2)$,
$E_A=[2(A-1)/p^2](1+\varepsilon_1p+\cdots)$).

\emph{Evidence grading (per coefficient, honest).} $\delta_1$,
$\varepsilon_1$ (hence $c_1$), $\delta_2$, and $c_2$ are theorems for
\emph{every} $A$ --- derived with $A$ a symbolic parameter (exact arithmetic
throughout). $c_3$'s assembly is likewise symbolic-$A$, but it consumes the
$\delta_3$ input, and $\delta_3$ is exact per-alphabet ($A=2,\dots,10$: nine
exact rational points overdetermining the stated degree-$3$ numerator) with no
symbolic-$A$ derivation claimed --- so the third order is
exact-at-listed-alphabets. $\delta_4$ and $c_4$ are exact at $A=2,\dots,7$
($c_4$'s numerator overdetermined already by $A\le6$; $\delta_4$'s determined
by $A\le6$ --- both confirmed \emph{out of sample} at $A=7$); symbolic-$A$ likewise not claimed. Consequences: $c_1>0$
for every $A\ge3$ (the corner margin opens \emph{linearly}; the $A=2$ corner,
where $c_1=0$ and the margin is $\tfrac18p^2$, is the hardest); the two-term
margin stays positive to $p^*=c_1/|c_2|$ ($\approx0.37$ at $A=3$ rising to
$\approx0.72$ at $A=6$); and $c_3>0$ for $A\le6$, $c_4>0$ for $A\le10$ ---
the higher orders \emph{help} at small alphabets.

\emph{(E) made effective --- a non-asymptotic theorem at $A=3,\dots,6$.} For
each $A\in\{3,4,5,6\}$ there are explicit rational constants $(p_0,\mu,K_4)$
--- $p_0=\tfrac1{16}$ at $A=3$, $p_0=\tfrac1{12}$ at $A=4,5,6$;
$K_4\approx598.5,\ 128.8,\ 108.8,\ 109.7$ --- such that for \emph{all}
$p\in(0,p_0]$: (a) the reduced-cubic discriminant provably changes sign across
the bracket $B(p)=[\hat D(p)-\mu p^6,\ \hat D(p)+\mu p^6]$,
$\hat D=[p^2/(4(A-1))](1+2p+\delta_2p^2+\delta_3p^3)$ (Sturm at the bracket
ends); \emph{and} --- the no-root-below-bracket lemma, closing the one gap an
adversarial review found in the bracket argument --- the discriminant
numerator factors as $D^4\,C(D,p)$ and $C((1-v)D_{\mathrm{lo}}(p),p)$, as a
polynomial in $v$, has all coefficients positive and $p$-root-free on
$(0,p_0]$ ($A=3,4$; at $A=5,6$ all but one, with the paired sum of the $v^4$/$v^5$
coefficients root-free positive, whence their combined contribution stays
$\ge0$ on $[0,1]$) --- so the discriminant is positive on $(0,D_{\mathrm{lo}}(p)]$ and
the bracketed crossing \emph{is} the first crossing $D_c$; (b) for all $D\in
B(p)$, $|R_A(\pi)-\tfrac32-c_1p-c_2p^2-c_3p^3|\le K_4p^4$ (and
$|R_A(\pi)-\tfrac32-c_1p-c_2p^2|\le K_{\mathrm{eff}}p^3$,
$K_{\mathrm{eff}}\approx38.4,\ 11.2,\ 9.2,\ 9.2$), whence
$R_A(\pi)>\tfrac32$ for \emph{all} $p\in(0,p_0]$. The proof is exact rational
arithmetic end to end: Sturm bracket capture; the exact perturbation ladder to
$k\le11$ with bivariate coefficient bounds and remainder $p$-valuation $4$;
and the $k\ge12$ tail by a diagonally-rescaled resolvent (Kato) bound --- for
$|\eta|\le\gamma p$ the four non-branch eigenvalues stay in $|z|\le\tfrac25$
and the branch in $|z-1|\le\tfrac25$, so the branch \emph{is} the spectral
radius --- plus a Cauchy geometric tail. This is the corner program's first
fully certified non-asymptotic statement: the sharp endpoint (E) holds at the corner
bracket (the $D$-interval containing $D_c$) for all $p\in(0,p_0]$ at
$A=3,\dots,6$, not merely as $p\to0$. ($A=2$ needs no
effective theorem: the sharp endpoint is proven in closed form.)

\emph{Validated} in
\texttt{scripts/verify/probe\_7\_34\_monotone\_s\_leading\_D.py} (W1--W4 the
flatness lemma, exact symbolic),
\texttt{lemma\_7\_34\_e3\_beta\_sign.py} (X1--X8; X8 the two-region
P\'olya-shift positivity, exact --- the $O(D^2)$ monotonicity theorem for
every alphabet),
\texttt{lemma\_7\_34\_dc\_smallp\_delta1.py} (Y1--Y5; Y5 symbolic-$A$),
\texttt{lemma\_7\_34\_eps1\_slope\_coefficient.py} (Z1--Z6; Z6 symbolic-$A$
rank-one perturbation),
\texttt{lemma\_7\_34\_endpoint\_c2\_exact.py} (E1--E4; E4 symbolic-$A$),
\texttt{lemma\_7\_34\_endpoint\_c3\_exact.py} (Y1--Y4),
\texttt{lemma\_7\_34\_endpoint\_c4\_exact.py} (W1--W3; W3 the out-of-sample
$A=7$), and
\texttt{lemma\_7\_34\_endpoint\_effective\_bound.py} (C1/C1b/C2/C3, exact
rational) --- all PASS.

\textbf{Remark 7.34m$'''''$ (the sharp-floor campaign II: regional scope ---
the corner pocket, the Schur certificate, the threshold upper bound, and the
mid-band instrument).} The campaign's honest map of where each mechanism
works, and where it provably does not.

\emph{The corner pocket (a scope correction).} The two monotonicity statements
that organize the campaign are \emph{false as global conjectures}; they fail
in a bounded pocket at the near-uniform corner, found by a hostile stress
sweep ($\sim2\times10^4$ points, $A\le32$, $p$ up to
$(1-10^{-6})\tfrac{A-1}A$, $D$ up to $0.99999\,D_c$, $50$--$100$ dps) and
independently reproduced by a separate implementation to $9$ significant
figures. Endpoint $D$-monotonicity ($dR_A(\pi;D)/dD\le0$ on $(0,D_c]$) fails
for \emph{every} $A\ge3$ with $p$ close enough to $\tfrac{A-1}A$ --- onset
$p\big/\tfrac{A-1}A\approx0.9999,\ 0.999,\ 0.997,\ 0.995,\ 0.99,\ 0.98$ at
$A=3,4,5,6,8,{\ge}12$, all clean at $\le0.97$; witness $A=16$,
$p=0.99\cdot\tfrac{15}{16}$: $R(\pi;0.8D_c)=1.9088\ldots<R(\pi;0.9D_c)
=1.9320\ldots$. The $s$-monotonicity (M) fails for large alphabets in the same
pocket ($A=16,24$ from $0.995$, $A=32$ from $0.99$; $A\le12$ clean to
$0.9999$), the $s$-argmin flipping to $s\to0^+$ (the anchor
$R_A(0^+)=2\bar v_A/[D(1-D)]$ of Remark 7.34m$'$) --- so deep in the pocket
the interior convexity of Remarks 7.34m$'$/7.34m$''$ itself fails, not merely
its proof. What survives --- the point of the sweep --- is the \emph{target}:
$R_A(s;D)\ge\tfrac32$ holds at every sweep point, with global minimum
$1.50000999\ldots$ at the usual tight corner and \emph{pocket} minimum
$1.7404$ (at $A=32$) --- an absolute margin $+0.2404$, four orders of
magnitude above the $\sim10^{-5}$ tight-corner margin. The
monotone-reduction is therefore \emph{regional}: on the clean region
$p\le0.97\,\tfrac{A-1}A$ (where every committed order-by-order theorem lives)
it is monotone $+$ endpoint; in the pocket the target holds by direct margin.
(Evidence level: hostile-sweep numeric, independently reproduced.)

\emph{The Schur small/moderate-$p$ certificate (whole Gray region, per
$(A,p)$).} A twice-peeled Schur fixed-point localization of the replica Perron
branch certifies $R_A(s;D)\ge\tfrac32$ for all $0<D\le D_c$ and all $s$ at
once, per $(A,p)$: the structural identities (the perturbation blocks have
zero all-equal column, so the coupling column is exact; $q_{00}=1$;
$P_1=(A-1)(\eta+\bar\eta)+a_{11}\eta\bar\eta$) reduce the branch to
\[
  \lambda=1+\frac{P_1}{\lambda}+\frac{P_2}{\lambda^2}
  +\frac{w^TQ_{rr}(\lambda-Q_{rr})^{-1}Q_{rr}\,c_r}{\lambda^2},
\]
with explicit weighted norms giving a self-consistent polydisc radius $r$ and
sup-bound $|\lambda-1|\le\rho(r)$; Cauchy bounds $|m_{jk}|\le\rho\,r^{-(j+k)}$
then dominate the two-variable perturbation tail geometrically, the exact part
being the two-variable ladder to total degree $K\le26$. Certified coverage
(margins at $K=26$): $p\big/\tfrac{A-1}A\in\{0.35,0.40\}$ at $A=3,4$ and
$\{0.35,0.40,0.45\}$ at $A=5,6$ (margins $+0.056$ to $+0.149$), and the
smaller $p\big/\tfrac{A-1}A\in\{0.05,0.10,0.20,0.30\}$ certified \emph{per
$p$} (margins $+0.016$ to $+0.114$ at $A=3,5$; an a-fortiori inheritance from
larger $p$ is \emph{not} claimed --- adversarial-review note). Honest grading:
the $p$-coverage is \emph{discrete} (certified at the listed $p$; the
continuum between them, and the $(D,s)$ evaluation inside each certificate,
are grid evidence pending interval reasoning --- mechanical but unfinished);
and the frontier is mapped, not hidden: $p\big/\tfrac{A-1}A\ge0.45$--$0.50$
fails at $K=26$ (would need $K\approx50$), while $\ge0.55$ admits \emph{no}
norm-certifiable polydisc radius at all --- the branch moves $\approx0.6$ from
$1$, genuinely non-perturbative.

\emph{An elementary threshold upper bound (proven at $A=3,\dots,8$; equality
at $A=2$).} On the whole parameter range,
\[
  D_c(A,p)\;<\;\bar D(A,p):=\frac{2(A-1)}A
  \Bigl[\tfrac12-\frac{\sqrt{1-q}}{2-q}\Bigr],\qquad
  q=\frac{pA}{A-1},\qquad p\in\bigl(0,\tfrac{A-1}A\bigr),
\]
proven per-$A$ by one-variable Sturm certificates after the
radical-eliminating substitution $s=\sqrt{1-q}$: the reduced-cubic
discriminant is positive as $D\to0^+$ (its numerator $D^4\,C(D,p)$ with
$C(0,p)$ root-free positive), negative at $\bar D(p)$ (the substituted
numerator root-free in $s$ on $(0,1)$ with negative sign), and its first sign
change is $D_c$ (the collision mechanism, Remark 7.34o$'$); at $A=2$, $\bar D$
\emph{equals} the exact closed-form $D_c$ identically. An explicit elementary
domain bound for the mid-band program and a two-sided companion to the
small-$p$ ladder.

\emph{The mid-band instrument (an instrument and a program --- not a finished
certificate).} Between the Schur frontier and the pocket lies the mid-band
$p\big/\tfrac{A-1}A\in(0.45,0.98)$, where no perturbation-about-$1$
architecture can work (the mapped obstruction above). The instrument: on the
Gray locus the $5\times5$ pattern-quotient spectrum appears \emph{exactly
real} --- worst $|\mathrm{Im}\,\lambda|\sim10^{-41}$ at $40$ dps across a
$288$-point band grid ($A\in\{3,4,6\}$, six values of $p\big/\tfrac{A-1}A$
spanning $0.5$--$0.95$, $D/D_c\in\{0.3,0.6,0.9,0.99\}$,
$s\in\{30^\circ,90^\circ,150^\circ,180^\circ\}$) --- and for a real-rooted
monic quintic $\chi(\lambda)=\det(\lambda I-Q)$ the target localizes with
\emph{no eigenvalue extraction}: with
$\Lambda:=L(s)/|C_s/C_0|^2$, $L(s)=1-\tfrac32D(1-D)(1-\cos s)$,
\[
  \chi^{(k)}(\Lambda)>0,\ k=0,\dots,4
  \quad\Longleftrightarrow\quad\text{all roots}<\Lambda
  \quad\Longleftrightarrow\quad\rho(Q)<\Lambda
  \quad\Longleftrightarrow\quad R_A(s;D)>\tfrac32,
\]
and all five derivative signs hold at \emph{every} grid point (verified
numeric). The structural lead behind the realness: the $\eta$-grading of the
quotient is column-determined, $Q=M\cdot\mathrm{diag}(1,\bar\eta,\eta,
|\eta|^2,|\eta|^2)$ with $M$ real and $\eta$-free (exact symbolic). The
program --- stated, not claimed: a structural real-rootedness lemma for this
column grading, plus Sturm/P\'olya sign-certification of the five rational
functions $\chi^{(k)}(\Lambda)$ of $(A,p,D,s)$ over the band --- the same
certification class discharged repeatedly in this arc.

\emph{The assembled map.} The sharp floor $R_A\ge\tfrac32$ now stands: proven
in closed form at $A=2$ (all $s$, all $D\le D_c$); and at $A\ge3$ --- tight
corner proven effective ($A=3,\dots,6$, all $p\le p_0$; Remark 7.34m$''''$),
small/moderate $p$ Schur-certified at the listed $(A,p)$ over the whole
$(D,s)$-range, pocket held by direct margin ($\ge1.74$), mid-band held at grid
level by the instrument with its certification program stated. What the
achievability machinery consumes remains the softer
$\tfrac{11}{16}$-plus-curvature package of Lemma 7.34l --- the campaign
strengthens the constant and \emph{localizes} what remains open on this axis
(the mid-band certification; the $p$-continuum interpolation), adding no
hypothesis anywhere.

\emph{Validated} in
\texttt{scripts/verify/probe\_7\_34\_corner\_pocket\_refutation.py} (H1 the
$D$-monotonicity witness; H4 the $s$-monotonicity witness; H2 the pocket floor
$\ge1.74$; H3 the clean-region boundary),
\texttt{lemma\_7\_34\_schur\_smallp\_certificate.py} (G1 the $K=26$ coverage
with margins; G2 the honest frontier),
\texttt{lemma\_7\_34\_dbar\_upper\_bound.py} (V1 the $A=2$ equality, symbolic;
V2/V2b the two Sturm legs, $A=3..8$; V3 slack sanity), and
\texttt{probe\_7\_34\_midband\_charpoly\_route.py} (S1 the $288$-point
real-spectrum $+$ derivative-certificate scout, numeric; S2 the column-grading
factorization, exact symbolic) --- all PASS.

\textbf{Theorem 7.34n (the general A-ary \emph{asymmetric} Gray-region
rate--distortion dispersion --- the grid capstone).} \emph{Let $X$ be a
stationary irreducible aperiodic A-ary Markov chain with a general (possibly
non-symmetric) transition matrix $T$, stationary $\pi$, under Hamming distortion,
with the A-ary symmetric backward channel $K(D)$ of Lemma 7.34k. There is a
positive Gray threshold $D_c(T)>0$ --- the all-$n$ SLB-tight
(valid-deconvolution) region --- whose binding word is a $2$-symbol alternating
word. On $0<D\le D_c(T)$ the operational fixed-$\varepsilon$ excess-distortion
dispersion satisfies the} \textbf{converse}
\[
  V_{\mathrm{op}}(D)\;\ge\;V_{\mathrm{lossless}}
  \;=\;\mathrm{Var}_\pi(g)+2\sum_{k\ge1}\mathrm{Cov}(g_0,g_k),\qquad
  g(x,y):=-\log_2 T(x,y),
\]
\emph{\textbf{unconditionally}} \emph{(the Gordin/Green--Kubo varentropy rate of
$-\log_2 P(X^n)$, carrying the Markov memory term), and the matching}
\textbf{achievability} \emph{$V_{\mathrm{op}}(D)\le V_{\mathrm{lossless}}$ modulo
the convexity residual of Remarks 7.34m$'$/7.34m$''$ (established at leading order
in $D$, all-$D$ on the memory bulk), so that $V_{\mathrm{op}}=V_{\mathrm{lossless}}$.}

\emph{Proof. Converse.} Remark 7.34d is source-agnostic; its two hypotheses hold:
(i) the varentropy CLT for $i_n=-\log_2 P(X^n)$ (Theorem 7.28, any finite-state
Markov), and (ii) the all-$n$ SLB-tight region $D_c(T)>0$ (non-empty: at $D=0$,
$K=I$ and $P_{Y^*}=P_X\ge0$; $D_c(T)$ is the first $D$ at which the alternating
binding word's deconvolution goes negative). On the region the \emph{dual
identity} holds for \emph{any} source:
\[
  M(x^n):=\sum_y P_{Y^*}(y)\,\nu^{\,d_H(x^n,y)}
  =(K_\nu K^{-1})^{\otimes n}P_X(x^n)=\gamma^n P_X(x^n),
  \qquad \nu=\tfrac{D}{(A-1)(1-D)},
\]
with $\gamma$ \emph{$x$-independent}, because $K_\nu K^{-1}=\gamma I$ is scalar:
writing $K_\nu=(1-\nu)I+\nu J$ (the $\nu^{d_H}$ kernel) and
$K^{-1}=\tfrac1{a_0}(I-b_0 J)$ ($a_0=1-D-b_0$, $b_0=\tfrac D{A-1}$), the
$J$-component of $K_\nu K^{-1}$ vanishes \emph{exactly} at $\nu(1-D)=b_0$, i.e.\
the SLB tilt --- a property of the A-ary symmetric \emph{channel}, independent of
the source. Hence the SLB shift $j_n-i_n$ is deterministic,
$\mathrm{Var}(j_n)=\mathrm{Var}(i_n)$ (converse ratio $\rho=1$, machine-checked to
$10^{-15}$), and the one-shot Kostina--Verd\'u $d$-tilted converse with the $i_n$
CLT gives $V_{\mathrm{op}}\ge V_{\mathrm{lossless}}$. The varentropy rate is the
\emph{Gordin} variance: for a non-symmetric $T$ the increments
$g_t=-\log_2 T(X_{t-1},X_t)$ are \emph{not} i.i.d.\ (unlike the symmetric chain of
Theorem 7.34m, whose increment law is state-independent), so the memory
cross-terms $2\sum_k\mathrm{Cov}(g_0,g_k)$ are present and nonzero (e.g.\ $+0.11$
to $+1.13$ of $V$ in the tested chains).

\emph{Achievability.} The chain-agnostic machinery consumes (a) the replica
spectral floor $\sup g<1$ with the $s=\pi$ binding, (b) Markov bounded-differences
(Paulin 2015), (c) deconvolution validity $D_c(T)$. Item (a)'s binding is exactly
the convexity residual of Remarks 7.34m$'$/7.34m$''$, established at leading order
in $D$ and all-$D$ on the memory bulk. \hfill$\square$

\emph{Scope (honest).} The \emph{converse} is the first fully-general operational
RD-dispersion lower bound for discrete Markov sources (any finite-state $T$,
A-ary symmetric channel, Gray region); it specialises to Remarks 7.34b (binary
symmetric), 7.34i (binary asymmetric), Theorem 7.34m (A-ary symmetric), and to
Kostina (2017) in the i.i.d.\ limit. The Markov memory term distinguishes it from
the symmetric i.i.d.-increment cases. The threshold \emph{mechanism} also
differs: for a non-symmetric $T$ the binding pair's period-2 product $B_iB_j$
($B_y[a,b]=K^{-1}[y,a]T[a,b]$) has \emph{complex} dominant eigenpair already
\emph{below} $D_c$ (complex-onset $D_{\mathrm{complex}}<D_c$), so $D_c$ is the
oscillation-driven \emph{negativity} onset, not the complexity onset (the two
coincide only in the symmetric case). Achievability inherits the convexity
residual. \emph{Validated} in \texttt{scripts/verify/probe\_7\_34n\_aary\_asym.py}
(N1 the threshold $D_c(T)>0$; N2 the $2$-symbol alternating binding word; N3 the
complex/negativity onset $D_{\mathrm{complex}}<D_c$; N4 the Gordin
$V_{\mathrm{lossless}}$ with nonzero memory term; N5 the source-agnostic dual
identity $\Rightarrow\rho=1$ --- all PASS).

\textbf{Remark 7.34n$'$ (the all-order dual identity: lossy-on-Gray $=$ lossless,
shifted --- and the third-order $\tfrac12\log n$ prefactor).} The converses of
Theorems 7.34m and 7.34n use only $\mathrm{Var}(j_n)=\mathrm{Var}(i_n)$, but the
dual identity $M(x^n)=\gamma^n P_X(x^n)$ gives far more. The $n$-block d-tilted
information is, by definition (Kostina--Verd\'u),
\[
  j_n(x^n,D)=-\lambda^* nD-\log_2 M(x^n)=i_n(x^n)-n\,c,\qquad
  c:=\lambda^*D\log_2 e+\log_2\gamma,
\]
a \emph{deterministic} shift of the source surprisal $i_n=-\log_2 P(X^n)$
(machine-checked $x$-independent to $10^{-14}$, N6) --- because $M=\gamma^n P_X$
exactly. Hence on the Gray region the \emph{entire distribution} of the d-tilted
information, not merely its variance, equals that of $i_n$ shifted by $-nc$: all
cumulants and all Edgeworth terms coincide. Therefore \emph{every} term of the
finite-blocklength expansion that the d-tilted information governs transfers
verbatim from the lossless source coding of $X$ to the lossy excess-distortion
problem on Gray, shifted by $nR(D)$ (with $R(D)=h-c$):
\[
  L_{\mathrm{lossy}}(n,\varepsilon,D)=L_{\mathrm{lossless}}(n,\varepsilon)-n\,c
  \qquad\text{(to the order the Kostina--Verd\'u bounds resolve).}
\]
In particular the dispersion ($V_{\mathrm{op}}=V_{\mathrm{lossless}}$, recovering
Theorems 7.34m/7.34n) \emph{and} the third-order prefactor coincide: the
long-open lossy third-order on Gray is reduced to the lossless source-coding
third-order, which is $\tfrac12\log n$ in the generic non-lattice case (the
$i_n$-lattice being the only obstruction) by Kontoyiannis--Verd\'u (2014). Thus
\[
  L_{\mathrm{lossy}}(n,\varepsilon,D)=nR(D)+\sqrt{nV_{\mathrm{lossless}}}\,
  Q^{-1}(\varepsilon)+\tfrac12\log n+O(1)\quad\text{on the Gray region (non-lattice $i_n$),}
\]
the same $\tfrac12\log n$ as lossless. The dual identity delivers this for the
\emph{converse} (lower-bound) third-order for free, being a functional of the
$j_n$ marginal alone; the matching \emph{achievability} third-order, like the
dispersion term, inherits the convexity residual of Remarks 7.34m$'$/7.34m$''$
(fully proven for BSMS on its whole Gray region, closed modulo the $a_k$ bound on
the general-A memory bulk per the ledger of Remark 7.34m$'''$), so the two-sided
expansion holds wherever that achievability does, not from the dual identity
alone. \emph{Validated} in
\texttt{scripts/verify/probe\_7\_34n\_aary\_asym.py} (N6: $j_n-i_n$ deterministic
to $10^{-14}$).

\textbf{Remark 7.34o (the beyond-Gray regime $D>D_c$: where the dispersion
plateau ends).} Theorems 7.34b/i/m/n give $V_{\mathrm{op}}=V_{\mathrm{lossless}}$
on the Gray region $0<D\le D_c$. The following maps what happens beyond it (the
exact value is open).

\emph{The plateau ends at $D_c$.} The per-$n$ deconvolution threshold $D_c^{(n)}$
--- the smallest $D$ at which the $n$-block $(K^{-1})^{\otimes n}P_X$ acquires a
negative entry (on the alternating binding word) --- \emph{decreases
monotonically} to the all-$n$ threshold $D_c=\inf_n D_c^{(n)}$ (for $p=0.2$:
$D_c^{(2)}=0.113\to D_c^{(10)}=0.0174\to D_c=0.0159$). Hence at finite $n$ the
dispersion plateau extends to $D_c^{(n)}>D_c$, but in the operational
($n\to\infty$) limit it ends exactly at $D_c$.

\emph{Why the Gray machinery stops.} For $D>D_c$ the SLB is no longer tight, the
deconvolution $P_{Y^*}$ ceases to be a valid distribution, the dual identity
$M=\gamma^n P_X$ --- hence $j_n=i_n-nc$ (Remark 7.34n$'$) --- fails, and the
converse's deterministic-shift argument no longer applies.

\emph{What replaces it.} With the optimal (Blahut--Arimoto) output, the
dispersion $V_{\mathrm{op}}(D)=\lim_n\mathrm{Var}(j_n)/n$ equals
$V_{\mathrm{lossless}}$ on $[0,D_c]$ and \emph{decreases} below it for $D>D_c$,
continuously at $D_c$ (the $d$-tilted information being smoother than the
surprisal once distortion is permitted). The finite-$n$ descent is
\emph{illustrative of the plateau-then-descent shape}, not the operational value:
e.g.\ at $p=0.2$, $n=7$, $V_{\mathrm{op}}/n\approx0.549$ is flat through $D_c$,
falling to $0.524$ at $3D_c$ and $0.347$ at $6D_c$ --- a finite-block snapshot,
whereas the operational $V_{\mathrm{op}}(D)$ is the $n\to\infty$ limit.

\emph{Why the exact value is open, and the shared obstruction.} The exact
$V_{\mathrm{op}}(D)$ for $D>D_c$ requires the optimal non-SLB-tight, memory-aware
test channel and \emph{remains open} --- the last of the three directions named
open after Theorem 7.34m (the general-asymmetric \emph{converse} and the
converse $\tfrac12\log n$ prefactor having been closed by Theorem 7.34n and
Remark 7.34n$'$; the matching general-A achievability remaining modulo the $a_k$
residual of Remark 7.34m$'''$). Its clean
characterization is the Green--Kubo variance rate
$V_{\mathrm{conv}}(D)=\lim_n\mathrm{Var}(j_n)/n=\chi''(0)$ (the second
cumulant-rate of the $\lambda^*$-tilted joint (source, optimal-output) transfer
operator), equivalently the converse ratio $\rho(D)=\lim_n\mathrm{Var}(j_n)/
\mathrm{Var}(i_n)$: elementary in Gray ($\rho\equiv1$, the plateau), but off Gray
the variance rate of an \emph{infinite-dimensional} operator, because the
$n$-letter optimal output requires a \emph{growing Markov order}
(Eswaran--Gastpar 2022). This is the \emph{same growing-order obstruction} that
walls the in-Gray general-A achievability capstone (the $a_k$ convexity bound of
Remark 7.34m$'''$, row a3) and the BSMS beyond-Gray analysis (Remark 7.34b's
``Beyond the Gray region'' discussion): a
single structural wall governs both the in-Gray 2nd-order achievability \emph{and}
the off-Gray exact value. Consequently $\rho(D)$ (and hence the exact
$V_{\mathrm{op}}$ for $D>D_c$) has \emph{no elementary closed form} --- the best
one-parameter fit $((1-2D)/(1-2D_c))^a$ carries a $p$-dependent exponent with
max-error far above what a closed form would need (the same beyond-Gray
discussion) --- so the
residual is the optimal memory-aware test channel itself, not a missing formula;
re-hunting an elementary $\rho(D)$ is provably futile.

\emph{Per-distortion.} The boundary is per-distortion: each balanced /
group-difference distortion has its \emph{own} Gray region with its own
threshold $D_c$ (Remark 7.34q; e.g.\ $D_c^{\mathrm{Lee}}<D_c^{\mathrm{Hamming}}$),
and the plateau-then-descent map above holds beyond each distortion's own $D_c$.
\emph{Validated} in
\texttt{scripts/verify/probe\_7\_34o\_beyond\_gray.py} (O1 $D_c^{(n)}\!\downarrow
D_c$; O2 the identity failure for $D>D_c$; O3 the plateau-and-descent of
$V_{\mathrm{op}}$ --- all PASS) and \texttt{probe\_7\_34\_offgray\_rho.py} (the
off-Gray $\rho(D)$: $n$-stable, $\rho=1$ at $D_c$, $\to0$ as $D\to\tfrac12$, no
elementary closed form, the Green--Kubo $\chi''(0)$ characterization with the
growing-Markov-order obstruction).

\textbf{Remark 7.34o$'$ (beyond-Gray structure: the collision mechanism, the
$1/n^2$ finite-$n$ law, and the essential-singularity conjecture).} Three
structural facts sharpen the beyond-Gray map of Remark 7.34o.

\emph{The threshold is an eigenvalue collision (proven, symbolic).} For the
BSMS$(p)$ deconvolution the period-$2$ alternating transfer product $C=B_0B_1$
($B_y[a,b]=K^{-1}[y,a]\,T[a,b]$, as in Theorem 7.34n) has
\[
  \mathrm{disc}(C)=\frac{p^2\bigl[p^2-4D(1-D)(1-p)^2\bigr]}{(1-2D)^2},
\]
vanishing \emph{exactly} at the Gray threshold ($p^2=4D(1-D)(1-p)^2$ is
precisely $D_c=\tfrac12-\sqrt{1-2p}/(2(1-p))$): $D_c$ \emph{is} the point
where the binding word's eigenvalue pair collides and turns complex --- the
mechanism behind the whole Gray family (its general-$A$ form is the $3\times3$
reduced-cubic discriminant driving the coefficient ladder and the $\bar D$
bound of Remarks 7.34m$''''$/7.34m$'''''$; for non-symmetric $T$ the collision
and negativity onsets split, Theorem 7.34n).

\emph{The finite-$n$ law (verified numeric, closed-form constant).} The
per-$n$ thresholds of Remark 7.34o approach $D_c$ at a universal $1/n^2$ rate:
\[
  D_c^{(n)}-D_c=\frac{K(p)^2}{n^2}\,\bigl(1+o(1)\bigr),\qquad
  K(p)=\frac{2\pi}{c(p)}=\frac{\pi\,p\,(1-2p)^{1/4}}{2\,(1-p)^{3/2}},
\]
where $c(p)=4(1-p)^{3/2}/\bigl(p(1-2p)^{1/4}\bigr)$ is the collision-phase
slope: past the collision the eigenpair rotates with phase
$\phi_2(D)\approx c(p)\sqrt{D-D_c}$ per period, and the $n$-block
deconvolution first goes negative when the accumulated phase completes a
cycle, $n\,\phi_2\approx2\pi$ --- whence the law. At $p=0.2$: $K^2=0.1493$
against measured $0.1485$--$0.1535$ ($n\le16$, $80$-bit control). So the
finite-$n$ Gray region of Remark 7.34o extends strictly beyond $D_c$ by an
explicitly quantified $\Theta(1/n^2)$ margin.

\emph{Beyond $D_c$ (each item at its stated evidence level).} The optimal
reverse channel remains \emph{exactly} product BSC$(D)$ throughout the
per-$n$ region $D\le D_c^{(n)}$ (solver evidence: total correlation and KL to
the product both $0$ within tolerance) --- the product structure survives
beyond $D_c$; the converse ratio $\rho(D)$ crosses $D_c$ continuously, with no
jump (Richardson, $n\le16$ --- consistent with Remark 7.34o's continuity
claim); and the support collapse beyond threshold cascades word family by word
family, each even-period defected-alternating family at its own exact
threshold from the same $2\times2$ discriminant mechanism. Finally, a
\emph{conjecture, labeled as such}: beyond Gray the rate gap closes with an
essential singularity,
\[
  R(D)-R_{\mathrm{SLB}}(D)
  =\exp\Bigl(-\,\frac{b_R(p)\,\bigl(1+o(1)\bigr)}{\sqrt{D-D_c}}\Bigr)
  \qquad(D\downarrow D_c)\qquad\text{(conjecture)},
\]
supported over $\approx8$ decades of residual, with no closed form for
$b_R(p)$ (the naive rare-alternating-run candidates are excluded); if true, it
explains why the plateau's end is invisible at every polynomial order --- the
two regimes are $C^\infty$-tangent at $D_c$. \emph{Validated} in
\texttt{scripts/verify/probe\_7\_34\_beyond\_gray\_structure.py} (B1 the
discriminant identity and its root at $D_c$, symbolic exact; B2 the $1/n^2$
law, drift $<15\%$ to $n=14$ with $80$-bit control at $n=16$; B3 the
WHT/XOR-convolution $O(N\log N)$ exact block-RD solver against the dense
reference --- all PASS; the product-channel survival, the $\rho$-continuity,
the cascade, and the singularity conjecture at the artifact evidence levels
stated above).

\textbf{Remark 7.34p (the lossy deviation hierarchy on the Gray region: §7.34
$\otimes$ §7.31--7.36).} The all-order identity $j_m=i_m-mc$ of Remark 7.34n$'$
holds for \emph{every} block length $m$ (the dual identity is per-length; the
per-symbol constant $c=\lambda^*D\log_2 e+\log_2\gamma$ is $m$-independent ---
machine-checked constant across $m$ to $10^{-16}$). The shift is thus a
\emph{deterministic linear drift}, so the \emph{centered} $d$-tilted-information
process equals the centered surprisal process \emph{pathwise}:
\[
  j_m-mR(D)=i_m-mh\qquad(\text{every }m\text{ on the Gray region},\ R(D)=h-c).
\]
Consequently the \emph{entire} second- and higher-order deviation hierarchy of §7
--- established there for the lossless surprisal/archive process --- transfers
verbatim to the lossy rate--distortion setting on the Gray region under
$i_m\mapsto j_m=i_m-mc$:
\begin{itemize}
\item \emph{Functional CLT / streaming buffer (Theorem 7.36 of [RNR-II]).} The centred
process is pathwise identical, so the lossy repair/write-buffer process converges
to the \emph{same} Brownian motion with the \emph{same} $V=V_{\mathrm{lossless}}$;
the half-normal high-water-mark and Lindley draw-up laws carry over unchanged.
\item \emph{Large-deviation overflow exponent (Theorem 7.31 of [RNR-II]).}
$E^{\mathrm{lossy}}(R)=E^{\mathrm{lossless}}(R+c)$: the lossy archive-overflow
exponent is the lossless one evaluated at the rate shifted by $c$.
\item \emph{Moderate deviations / RA fade (Theorem 7.33 of [RNR-II]).} The standardized
moderate-deviation law transfers (same $V$, mean shifted by $c$ per symbol).
\item \emph{Extreme-value worst-case query (Theorem 7.35 of [RNR-II]).} The maximal block
$d$-tilted information equals the maximal block surprisal minus $Kc$, so the
Gumbel law and the read-buffer provisioning $M_m=KR(D)+\sqrt{2KV\ln(N/K)}+\cdots$
follow with $h\mapsto R(D)$.
\end{itemize}
Thus a \emph{single} varentropy rate $V_{\mathrm{lossless}}$ governs the lossy
dispersion, streaming/read/write buffers, overflow exponent, and query-cost
extremes on the entire Gray region, exactly as in the lossless case: to all these
orders the rate--distortion problem on Gray \emph{is} the lossless source-coding
problem, shifted by $nR(D)$. This unifies the rate--distortion results of §7.34
with the deviation hierarchy of §7.31--7.36. \emph{Validated} in
\texttt{scripts/verify/probe\_7\_34p\_lossy\_devhier.py} (P1 the per-$m$ identity
with $m$-independent $c$; P2 the pathwise centered-process identity; P3 the
distribution shift $\Rightarrow$ tail/LDP/EVT transfer --- all PASS).

\textbf{Remark 7.34q (the distortion-type axis: the converse extends from
Hamming to any group-difference distortion --- Lee, etc.).} Theorems
7.34b/i/m/n are stated for Hamming distortion, but the source-agnostic converse
depends only on the backward channel and extends verbatim. For \emph{any}
difference distortion $d(x,y)=\rho(x\ominus y)$ on a finite abelian group $G$
(Hamming on $\mathbb Z_A$; \emph{Lee} $d=\min(|x-y|,A-|x-y|)$ on $\mathbb Z_A$;
$\ell_1$/circular variants; products), the SLB-achieving backward channel is the
\emph{Gibbs kernel} $K=K_\nu/Z$ with $K_\nu[x,y]=\nu^{\,d(x,y)}$,
$\nu=e^{-\lambda^*}$ and $Z=\sum_y\nu^{\,d(x,y)}$ ($x$-independent by group
symmetry). Hence
\[
  K_\nu K^{-1}=(ZK)K^{-1}=Z\,I=\gamma I\qquad(\gamma=Z),
\]
\emph{trivially} scalar, so the dual identity $M(x^n)=\sum_y P_{Y^*}(y)\,
\nu^{\,d(x^n,y)}=(K_\nu K^{-1})^{\otimes n}P_X=\gamma^n P_X$ holds for any source,
$j_n=i_n-nc$ is a deterministic shift, the converse ratio $\rho=1$, and
\[
  V_{\mathrm{op}}(D)\ge V_{\mathrm{lossless}}\quad\text{on the distortion's own
  Gray region, unconditionally.}
\]
For Hamming $\gamma=Z=1/(1-D)$, recovering Theorem 7.34n. The concrete second
instance is \emph{Lee} distortion: the dual identity holds ($A=3,\dots,6$;
$A=3$ Lee $=$ Hamming since all $d=1$), and the Lee Gray region is nonempty with a \emph{positive} all-$n$ threshold,
though \emph{smaller} than Hamming's --- the graded Lee circulant $K$ has a
sign-oscillating inverse, so $D_c^{\mathrm{Lee}}$ is small (e.g.\ $A=4,p=0.2$: the
per-$n$ thresholds $D_c^{(n)}$ decrease and stabilise to $D_c^{\mathrm{Lee}}
\approx3.9\times10^{-3}$, the binding word an alternating \emph{Lee-adjacent}
(distance-$1$) pair --- the same alternating mechanism as Hamming, on the
minimum-distortion pairs). $V_{\mathrm{lossless}}$
is unchanged (a source quantity, distortion-independent), and achievability
inherits the same convexity residual. So the entire RD-dispersion grid's
\emph{converse} is a statement about group-difference distortions, not Hamming
specifically; and by Remark 7.34n$'$ the all-order shift $j_n=i_n-nc$ likewise
carries the third-order $\tfrac12\log n$ and the lossy deviation hierarchy
(Remark 7.34p) to every such distortion on its Gray region. Indeed the
\emph{only} property used is that $Z=\sum_y\nu^{\,d(x,y)}$ is $x$-independent,
which holds for every \emph{balanced} distortion in Gray's (1971) sense --- each
row of the distortion matrix $[d(x,y)]$ a permutation of the others, so the
distance multiset $\{d(x,y):y\}$ is common to all $x$ --- the group-difference
(circulant) distortions being the concrete computable sub-case. Hence the
unconditional converse $V_{\mathrm{op}}\ge V_{\mathrm{lossless}}$ holds for the
\emph{entire balanced-distortion class} on its Gray region, the natural setting
of Gray's tight-SLB theorem. \emph{Validated} in
\texttt{scripts/verify/probe\_7\_34q\_lee\_general\_distortion.py} (Q1 the dual
identity $K_\nu K^{-1}=ZI$; Q2 the nonempty Lee Gray region; Q3 the operational
$j_n=i_n-nc$; Q4 $V_{\mathrm{lossless}}$ distortion-independent --- all PASS).

\textbf{References}

\textbf{{[}1{]}} Bernstein, J., Wang, Y.-X., Azizzadenesheli, K., and
Anandkumar, A. (2018). signSGD: Compressed Optimisation for Non-Convex
Problems. ICML.
\textbf{{[}2{]}} Burrows, M., and Wheeler, D. J. (1994). A Block-Sorting
Lossless Data Compression Algorithm. SRC Research Report 124, Digital
Equipment Corporation.
\textbf{{[}3{]}} Carlsson, G., Ishkhanov, T., de Silva, V., and
Zomorodian, A. (2008). On the Local Behavior of Spaces of Natural
Images. International Journal of Computer Vision, 76(1), 1-12.
\textbf{{[}4{]}} Chen, X., Mishra, N., Rohaninejad, M., and Abbeel, P.
(2018). PixelSNAIL: An Improved Autoregressive Generative Model.
International Conference on Machine Learning.
\textbf{{[}5{]}} Charikar, M., Lehman, E., Liu, D., Panigrahy, R.,
Prabhakaran, M., Sahai, A., and Shelat, A. (2005). The Smallest Grammar
Problem. IEEE Transactions on Information Theory, 51(7), 2554-2576.

\textbf{{[}8a{]}} Casel, K., Fernau, H., Gaspers, S., Gras, B., and
Schmid, M. L. (2021). On the Complexity of the Smallest Grammar Problem
over Fixed Alphabets. Theory of Computing Systems, 65(1), 158-203.

\textbf{{[}8b{]}} Ga\'nczorz, M. (2018). Entropy bounds for grammar
compression. arXiv:1804.08547 (also presented at DCC 2018/2020).

\textbf{{[}8c{]}} Kozma, L., and Voderholzer, J. (2024). The
Computational Complexity of Tokenisation. arXiv:2411.08671.

\textbf{{[}8d{]}} Je\.z, A. (2015). A really simple approximation of
smallest grammar. Theoretical Computer Science, 616, 141-150.

\textbf{{[}8e{]}} Karzanov, A. V. (1998). Minimum 0-Extensions of Graph
Metrics. European Journal of Combinatorics, 19(1), 71-101.

\textbf{{[}8f{]}} Schwartz, R., and Tur, S. (2023). Improved hardness
of approximation for graph problems via Boolean PCPs.
arXiv:2104.11670 (also conference version, ICALP 2023).
\textbf{{[}6{]}} Cover, T. M. (1973). Enumerative Source Encoding. IEEE
Transactions on Information Theory, 19(1), 73-77.
\textbf{{[}7{]}} Cover, T. M., and Thomas, J. A. (2006). Elements of
Information Theory, 2nd edition. Wiley.
\textbf{{[}8{]}} Donoho, D. L. (2000). High-Dimensional Data Analysis:
The Curses and Blessings of Dimensionality. AMS Math Challenges Lecture.
\textbf{{[}9{]}} Duda, J. (2009). Asymmetric Numeral Systems.
arXiv:0902.0271.
\textbf{{[}10{]}} Frey, B. J. (1997). Bayesian Networks for Pattern
Classification, Data Compression, and Channel Coding. Ph.D. thesis,
University of Toronto.
\textbf{{[}11{]}} Frey, B. J. (1998). Graphical Models for Machine
Learning and Digital Communication. MIT Press.
\textbf{{[}12{]}} Hinton, G. E., and van Camp, D. (1993). Keeping Neural
Networks Simple by Minimizing the Description Length of the Weights.
Conference on Learning Theory.
\textbf{{[}13{]}} Hoogeboom, E., Peters, J. W. T., van den Berg, R., and
Welling, M. (2019). Integer Discrete Flows and Lossless Compression.
NeurIPS.
\textbf{{[}14{]}} Huffman, D. A. (1952). A Method for the Construction
of Minimum-Redundancy Codes. Proceedings of the IRE, 40(9), 1098-1101.
\textbf{{[}15{]}} Knoll, B. (ongoing). Cmix: a lossless data compression
program. https://www.byronknoll.com/cmix.html.
\textbf{{[}16{]}} Levina, E., and Bickel, P. J. (2004). Maximum
Likelihood Estimation of Intrinsic Dimension. NIPS.
\textbf{{[}17{]}} Liveris, A. D., Xiong, Z., and Georghiades, C. N.
(2002). Compression of Binary Sources with Side Information at the
Decoder Using LDPC Codes. IEEE Communications Letters, 6(10), 440-442.
\textbf{{[}18{]}} Mentzer, F., Agustsson, E., Tschannen, M., Timofte,
R., and Van Gool, L. (2019). Practical Full Resolution Learned Lossless
Image Compression. CVPR.
\textbf{{[}19{]}} Pradhan, S. S., and Ramchandran, K. (2003).
Distributed Source Coding Using Syndromes (DISCUS): Design and
Construction. IEEE Transactions on Information Theory, 49(3), 626-643.
\textbf{{[}20{]}} Reed, I. S., and Solomon, G. (1960). Polynomial Codes
over Certain Finite Fields. Journal of SIAM, 8(2), 300-304.
\textbf{{[}21{]}} Rissanen, J. (1978). Modeling by Shortest Data
Description. Automatica, 14(5), 465-471.
\textbf{{[}22{]}} Rissanen, J., and Langdon, G. G. (1979). Arithmetic
Coding. IBM Journal of Research and Development, 23(2), 149-162.
\textbf{{[}23{]}} Salimans, T., Karpathy, A., Chen, X., and Kingma, D.
P. (2017). PixelCNN++: Improving the PixelCNN with Discretized Logistic
Mixture Likelihood and Other Modifications. ICLR.

\textbf{{[}29a{]}} Crutchfield, J. P., and Feldman, D. P. (2003).
Regularities Unseen, Randomness Observed: Levels of Entropy
Convergence. Chaos, 13(1), 25-54. DOI: 10.1063/1.1530990.
\textbf{{[}24{]}} Shannon, C. E. (1948). A Mathematical Theory of
Communication. Bell System Technical Journal, 27, 379-423 and 623-656.

\textbf{{[}31a{]}} Shannon, C. E. (1951). Prediction and Entropy of
Printed English. Bell System Technical Journal, 30(1), 50-64.
\textbf{{[}25{]}} Slepian, D., and Wolf, J. K. (1973). Noiseless Coding
of Correlated Information Sources. IEEE Transactions on Information
Theory, 19(4), 471-480.
\textbf{{[}26{]}} Tenenbaum, J. B., de Silva, V., and Langford, J. C.
(2000). A Global Geometric Framework for Nonlinear Dimensionality
Reduction. Science, 290(5500), 2319-2323.
\textbf{{[}27{]}} Townsend, J., Bird, T., Kunze, J., and Barber, D.
(2020). HiLLoC: Lossless Image Compression with Hierarchical Latent
Variable Models. ICLR.
\textbf{{[}28{]}} van den Berg, R., Gritsenko, A. A., Dehghani, M.,
Sonderby, C. K., and Salimans, T. (2021). IDF++: Analyzing and Improving
Integer Discrete Flows for Lossless Compression. ICLR.
\textbf{{[}29{]}} van den Oord, A., Kalchbrenner, N., Vinyals, O.,
Espeholt, L., Graves, A., and Kavukcuoglu, K. (2016). Conditional Image
Generation with PixelCNN Decoders. NeurIPS.
\textbf{{[}30{]}} Rhee, H., Jang, Y. I., Kim, S., and Cho, N. I. (2022).
LC-FDNet: Learned Lossless Image Compression with Frequency
Decomposition Network. CVPR, 6023-6032.
\textbf{{[}31{]}} Witten, I. H., Neal, R. M., and Cleary, J. G. (1987).
Arithmetic Coding for Data Compression. Communications of the ACM,
30(6), 520-540.
\textbf{{[}32{]}} Wyner, A. D., and Ziv, J. (1976). The Rate-Distortion
Function for Source Coding with Side Information at the Decoder. IEEE
Transactions on Information Theory, 22(1), 1-10.
\textbf{{[}33{]}} Ziv, J., and Lempel, A. (1977). A Universal Algorithm
for Sequential Data Compression. IEEE Transactions on Information
Theory, 23(3), 337-343.
\textbf{{[}34{]}} Ziv, J., and Lempel, A. (1978). Compression of
Individual Sequences via Variable-Rate Coding. IEEE Transactions on
Information Theory, 24(5), 530-536.
\textbf{{[}35{]}} Watanabe, S. (1960). Information Theoretical
Analysis of Multivariate Correlation. IBM Journal of Research and
Development, 4(1), 66-82.
\textbf{{[}36{]}} Studen\'y, M., and Vejnarov\'a, J. (1998). The
Multiinformation Function as a Tool for Measuring Stochastic
Dependence. In: \emph{Learning in Graphical Models}, Jordan, M. I.
(ed.), MIT Press, pp. 261-297.
\textbf{{[}37{]}} Lund, C., and Yannakakis, M. (1994). On the
Hardness of Approximating Minimization Problems. Journal of the ACM,
41(5), 960-981.
\textbf{{[}38{]}} Lehman, E., and Shelat, A. (2002). Approximation
Algorithms for Grammar-Based Compression. Proceedings of the
Thirteenth Annual ACM-SIAM Symposium on Discrete Algorithms (SODA),
205-212.
\textbf{{[}39{]}} Berman, P., and Karpinski, M. (1999). On Some
Tighter Inapproximability Results. Proceedings of the 26th
International Colloquium on Automata, Languages and Programming
(ICALP), Lecture Notes in Computer Science 1644, Springer, 200-209.
\textbf{{[}40{]}} Storer, J. A., and Szymanski, T. G. (1982). Data
Compression via Textual Substitution. Journal of the ACM, 29(4),
928-951.
\textbf{{[}41{]}} Ar{\i}kan, E. (2009). Channel Polarization: A Method
for Constructing Capacity-Achieving Codes for Symmetric Binary-Input
Memoryless Channels. IEEE Transactions on Information Theory, 55(7),
3051-3073.
\textbf{{[}42{]}} Azuma, K. (1967). Weighted Sums of Certain Dependent
Random Variables. T\^ohoku Mathematical Journal, 19(3), 357-367.
\textbf{{[}43{]}} Bellare, M. (2006). New Proofs for NMAC and HMAC:
Security Without Collision-Resistance. In: Advances in Cryptology ---
CRYPTO 2006, Lecture Notes in Computer Science 4117, Springer, 602-619.
\textbf{{[}44{]}} Bellare, M., Canetti, R., and Krawczyk, H. (1996).
Keying Hash Functions for Message Authentication. In: Advances in
Cryptology --- CRYPTO '96, Lecture Notes in Computer Science 1109,
Springer, 1-15.
\textbf{{[}45{]}} Bellare, M., and Namprempre, C. (2000). Authenticated
Encryption: Relations Among Notions and Analysis of the Generic
Composition Paradigm. In: Advances in Cryptology --- ASIACRYPT 2000,
Lecture Notes in Computer Science 1976, Springer, 531-545.
\textbf{{[}46{]}} Berger, T. (1971). Rate Distortion Theory: A
Mathematical Basis for Data Compression. Prentice-Hall, Englewood
Cliffs, NJ.
\textbf{{[}47{]}} Breslau, L., Cao, P., Fan, L., Phillips, G., and
Shenker, S. (1999). Web Caching and Zipf-like Distributions: Evidence
and Implications. Proceedings of IEEE INFOCOM, 126-134.
\textbf{{[}48{]}} Cover, T. M. (1975). A Proof of the Data Compression
Theorem of Slepian and Wolf for Ergodic Sources. IEEE Transactions on
Information Theory, 21(2), 226-228.
\textbf{{[}49{]}} Cover, T. M., and King, R. C. (1978). A Convergent
Gambling Estimate of the Entropy of English. IEEE Transactions on
Information Theory, 24(4), 413-421.
\textbf{{[}50{]}} Csisz\'ar, I., and K\"orner, J. (2011). Information
Theory: Coding Theorems for Discrete Memoryless Systems, 2nd edition.
Cambridge University Press. (First edition: Akad\'emiai Kiad\'o, 1981.)
\textbf{{[}51{]}} Cybenko, G. (1989). Approximation by Superpositions
of a Sigmoidal Function. Mathematics of Control, Signals, and Systems,
2(4), 303-314.
\textbf{{[}52{]}} Del\'etang, G., Ruoss, A., Duquenne, P.-A., Catt, E.,
Genewein, T., Mattern, C., Grau-Moya, J., Wenliang, L. K.,
Aitchison, M., Orseau, L., Legg, S., and Veness, J. (2024). Language
Modeling Is Compression. Proceedings of ICLR. arXiv:2309.10668.
\textbf{{[}53{]}} Doshi, V., Shah, D., and M\'edard, M. (2007). Source
Coding with Distortion through Graph Coloring. Proceedings of the IEEE
International Symposium on Information Theory (ISIT), 1501-1505.
\textbf{{[}54{]}} Doshi, V., Shah, D., M\'edard, M., and Effros, M.
(2010). Functional Compression through Graph Coloring. IEEE
Transactions on Information Theory, 56(8), 3901-3917.
\textbf{{[}55{]}} Equitz, W. H. R., and Cover, T. M. (1991). Successive
Refinement of Information. IEEE Transactions on Information Theory,
37(2), 269-275.
\textbf{{[}56{]}} Fagin, R. (1977). Generalized First-Order Spectra and
Polynomial-Time Recognizable Sets. In: Complexity of Computation,
SIAM-AMS Proceedings 7, 43-73.
\textbf{{[}57{]}} Ferragut, A., Rodr{\'\i}guez, I., and Paganini, F.
(2018). Optimizing TTL Caches under Heavy-Tailed Demand. Queueing
Systems, 88(3-4), 207-241.
\textbf{{[}58{]}} Grossi, R., and Vitter, J. S. (2005). Compressed
Suffix Arrays and Suffix Trees with Applications to Text Indexing and
String Matching. SIAM Journal on Computing, 35(2), 378-407.
\textbf{{[}59{]}} Hoeffding, W. (1963). Probability Inequalities for
Sums of Bounded Random Variables. Journal of the American Statistical
Association, 58(301), 13-30.
\textbf{{[}60{]}} Hoffmann, J., Borgeaud, S., Mensch, A., Buchatskaya,
E., Cai, T., Rutherford, E., et al. (2022). Training Compute-Optimal
Large Language Models. arXiv:2203.15556.
\textbf{{[}61{]}} Hornik, K., Stinchcombe, M., and White, H. (1989).
Multilayer Feedforward Networks Are Universal Approximators. Neural
Networks, 2(5), 359-366.
\textbf{{[}62{]}} Kaplan, J., McCandlish, S., Henighan, T., Brown,
T. B., Chess, B., Child, R., Gray, S., Radford, A., Wu, J., and
Amodei, D. (2020). Scaling Laws for Neural Language Models.
arXiv:2001.08361.
\textbf{{[}63{]}} Korada, S. B., and Urbanke, R. L. (2010). Polar
Codes Are Optimal for Lossy Source Coding. IEEE Transactions on
Information Theory, 56(4), 1751-1768.
\textbf{{[}64{]}} McDiarmid, C. (1989). On the Method of Bounded
Differences. In: Surveys in Combinatorics 1989, London Mathematical
Society Lecture Notes 141, Cambridge University Press, 148-188.
\textbf{{[}65{]}} McLeish, D. L. (1974). Dependent Central Limit
Theorems and Invariance Principles. Annals of Probability, 2(4),
620-628.
\textbf{{[}66{]}} Muthitacharoen, A., Chen, B., and Mazi\`eres, D.
(2001). A Low-Bandwidth Network File System. Proceedings of the 18th
ACM Symposium on Operating Systems Principles (SOSP), 174-187.
\textbf{{[}67{]}} Polyanskiy, Y., Poor, H. V., and Verd\'u, S. (2010).
Channel Coding Rate in the Finite Blocklength Regime. IEEE Transactions
on Information Theory, 56(5), 2307-2359.
\textbf{{[}68{]}} Rabin, M. O. (1981). Fingerprinting by Random
Polynomials. Technical Report TR-15-81, Center for Research in Computing
Technology, Harvard University.
\textbf{{[}69{]}} Richardson, T., and Urbanke, R. (2008). Modern Coding
Theory. Cambridge University Press.
\textbf{{[}70{]}} Rimoldi, B. (1994). Successive Refinement of
Information: Characterization of the Achievable Rates. IEEE Transactions
on Information Theory, 40(1), 253-259.
\textbf{{[}71{]}} Rissanen, J. (1996). Fisher Information and
Stochastic Complexity. IEEE Transactions on Information Theory, 42(1),
40-47.
\textbf{{[}72{]}} Said, A. (2003). Arithmetic Coding for Data
Compression. Technical Report, Hewlett-Packard Laboratories.
\textbf{{[}73{]}} Samson, P.-M. (2000). Concentration of Measure
Inequalities for Markov Chains and $\Phi$-Mixing Processes. Annals of
Probability, 28(1), 416-461.
\textbf{{[}74{]}} Shannon, C. E. (1959). Coding Theorems for a Discrete
Source with a Fidelity Criterion. IRE National Convention Record, Part 4,
142-163.
\textbf{{[}75{]}} Steinberg, Y., and Merhav, N. (2004). On Successive
Refinement for the Wyner-Ziv Problem. IEEE Transactions on Information
Theory, 50(8), 1636-1654.
\textbf{{[}76{]}} Witsenhausen, H. S. (1976). The Zero-Error Side
Information Problem and Chromatic Numbers. IEEE Transactions on
Information Theory, 22(5), 592-593.
\textbf{{[}77{]}} Krichevsky, R. E., and Trofimov, V. K. (1981). The
Performance of Universal Encoding. IEEE Transactions on Information
Theory, 27(2), 199-207.
\textbf{{[}78{]}} Baum, L. E., and Petrie, T. (1966). Statistical
Inference for Probabilistic Functions of Finite State Markov Chains.
Annals of Mathematical Statistics, 37(6), 1554-1563.
\textbf{{[}79{]}} Lauritzen, S. L., and Spiegelhalter, D. J. (1988).
Local Computations with Probabilities on Graphical Structures and
Their Application to Expert Systems. Journal of the Royal Statistical
Society, Series B, 50(2), 157-194.
\textbf{{[}80{]}} Dinur, I. (2007). The PCP Theorem by Gap
Amplification. Journal of the ACM, 54(3), Article 12.
\textbf{{[}81{]}} Dinur, I., Kindler, G., Raz, R., and Safra, S.
(2003). Approximating CVP to Within Almost-Polynomial Factors is
NP-Hard. Combinatorica, 23(2), 205-243.
\textbf{{[}82{]}} Barron, A. R. (1985). The Strong Ergodic Theorem
for Densities: Generalized Shannon-McMillan-Breiman Theorem. Annals
of Probability, 13(4), 1292-1303.
\textbf{{[}83{]}} Gray, R. M. (1990). Entropy and Information
Theory. Springer.
\textbf{{[}84{]}} Pradhan, S. S., and Ramchandran, K. (1999).
Distributed Source Coding Using Syndromes (DISCUS): Design and
Construction. Proceedings of the Data Compression Conference (DCC),
158-167.
\textbf{{[}85{]}} Stratonovich, R. L. (1960). Conditional Markov
Processes. Theory of Probability and Its Applications, 5(2), 156-178.
\textbf{{[}86{]}} Wainwright, M. J., and Jordan, M. I. (2008).
Graphical Models, Exponential Families, and Variational Inference.
Foundations and Trends in Machine Learning, 1(1-2), 1-305.
\textbf{{[}87{]}} Arora, S., Babai, L., Stern, J., and Sweedyk, Z.
(1997). The Hardness of Approximate Optima in Lattices, Codes, and
Systems of Linear Equations. Journal of Computer and System Sciences,
54(2), 317-331.
\textbf{{[}88{]}} Birch, B. J. (1962). Approximations for the Entropy
for Functions of Markov Chains. Annals of Mathematical Statistics,
33(3), 930-938.
\textbf{{[}89{]}} Pakzad, P., and Anantharam, V. (2002). Estimation
and Marginalization Using Kikuchi Approximation Methods. Neural
Computation, 17(8), 1836-1873.
\textbf{{[}90{]}} Kieffer, J. C., and Yang, E.-H. (2000). Grammar-Based
Codes: A New Class of Universal Lossless Source Codes. IEEE Transactions
on Information Theory, 46(3), 737-754.
\textbf{{[}91{]}} Merlev\`ede, F., Peligrad, M., and Rio, E. (2009).
Bernstein Inequality and Moderate Deviations under Strong Mixing
Conditions. In: High Dimensional Probability V: The Luminy Volume,
IMS Collections 5, 273-292.
\textbf{{[}92{]}} Bradley, R. C. (2005). Basic Properties of Strong
Mixing Conditions: A Survey and Some Open Questions. Probability
Surveys, 2, 107-144.
\textbf{{[}93{]}} Rio, E. (2017). Asymptotic Theory of Weakly
Dependent Random Processes. Probability Theory and Stochastic
Modelling 80, Springer.
\textbf{{[}94{]}} Kontorovich, L., and Ramanan, K. (2008).
Concentration Inequalities for Dependent Random Variables via the
Martingale Method. Annals of Probability, 36(6), 2126-2158.
\textbf{{[}95{]}} Davydov, Yu.\ A. (1968). Convergence of
Distributions Generated by Stationary Stochastic Processes. Theory
of Probability and Its Applications, 13(4), 691-696.
\textbf{{[}96{]}} Bradley, R. C. (2007). Introduction to Strong
Mixing Conditions, Volumes 1--3. Kendrick Press.
\textbf{{[}97{]}} Brown, P. F., Della Pietra, V. J., Mercer,
R. L., Della Pietra, S. A., and Lai, J. C. (1992). An Estimate of
an Upper Bound for the Entropy of English. Computational Linguistics,
18(1), 31-40.
\textbf{{[}98{]}} Lezaud, P. (1998). Chernoff-Type Bound for Finite
Markov Chains. Annals of Applied Probability, 8(3), 849-867.
\textbf{{[}99{]}} Che, H., Wang, Z., and Tung, Y. (2002). Analysis
and Design of Hierarchical Web Caching Systems. Proceedings of IEEE
INFOCOM, 1416-1424.
\textbf{{[}100{]}} Aldous, D., and Diaconis, P. (1987). Shuffling
Cards and Stopping Times. American Mathematical Monthly, 93(5),
333-348.
\textbf{{[}101{]}} Knudsen, J. B., and Tsitsiklis, J. N. (1990).
On the Time Required to Drain a Network. Mathematics of Operations
Research, 15(4), 685-708.
\textbf{{[}102{]}} Paulin, D. (2015). Concentration Inequalities
for Markov Chains by Marton Couplings and Spectral Methods.
Electronic Journal of Probability, 20(79), 1-32.
\textbf{{[}103{]}} Aldous, D., and Fill, J. A. (2002).
Reversible Markov Chains and Random Walks on Graphs.
Unfinished monograph, available at
https://www.stat.berkeley.edu/\textasciitilde aldous/RWG/book.html.
\textbf{{[}104{]}} Diaconis, P., and Stroock, D. (1991). Geometric
Bounds for Eigenvalues of Markov Chains. Annals of Applied
Probability, 1(1), 36-61.
\textbf{{[}105{]}} King, W. F. (1971). Analysis of Demand Paging
Algorithms. Proceedings of the IFIP Congress, 485-490.
\textbf{{[}106{]}} Kesidis, G., Shan, Y., Fumagalli, A., Gulati,
A., and Singh, A. (2017). An Analysis of LRU Replacement Caches
under Independent Reference Models. arXiv:1704.04849.
\textbf{{[}107{]}} Niyogi, P., Smale, S., and Weinberger, S. (2008).
Finding the Homology of Submanifolds with High Confidence from
Random Samples. Discrete and Computational Geometry, 39(1-3),
419-441. (Introduces the ``reach'' parameter of a sub-manifold and
gives quantitative manifold-learning sampling bounds.)
\textbf{{[}108{]}} Diaconis, P., and Freedman, D. (1987).
A dozen de~Finetti-style results in search of a theory. Annales de
l'Institut Henri Poincar\'e (B) Probabilities et Statistiques, 23(S2),
397-423.
\textbf{{[}109{]}} Gersho, A., and Gray, R. M. (1992). Vector
Quantization and Signal Compression. Kluwer Academic Publishers.
(Bennett high-resolution quantization theory, Chapter 5.)
\textbf{{[}110{]}} Talagrand, M. (1995). Concentration of Measure
and Isoperimetric Inequalities in Product Spaces. Publications
Math\'ematiques de l'IH\'ES, 81, 73-205.
\textbf{{[}111{]}} Bennett, W. R. (1948). Spectra of Quantized
Signals. Bell System Technical Journal, 27(3), 446-472.
\textbf{{[}112{]}} Willems, F. M. J., Shtarkov, Y. M., and Tjalkens,
T. J. (1995). The Context-Tree Weighting Method: Basic Properties.
IEEE Transactions on Information Theory, 41(3), 653-664.
\textbf{{[}113{]}} Kostina, V., and Verd\'u, S. (2012). Fixed-Length
Lossy Compression in the Finite Blocklength Regime. IEEE Transactions on
Information Theory, 58(6), 3309-3338.
\textbf{{[}114{]}} Kostina, V. (2017). Data Compression with Low
Distortion and Finite Blocklength. IEEE Transactions on Information
Theory, 63(7), 4268-4285.
\textbf{{[}115{]}} Kontoyiannis, I., and Verd\'u, S. (2014). Optimal
Lossless Data Compression: Non-Asymptotics and Asymptotics. IEEE
Transactions on Information Theory, 60(2), 777-795.
\textbf{{[}116{]}} Tian, C., and Kostina, V. (2019). The Dispersion of
the Gauss--Markov Source. IEEE Transactions on Information Theory,
65(10), 6355-6384.
\textbf{{[}117{]}} Kostina, V., and Tian, C. (2020). Nonasymptotic and
Second-Order Achievability Bounds for Coding With Side Information.
IEEE Transactions on Information Theory (Gauss--Markov rate-distortion
dispersion, achievability side).
\textbf{{[}118{]}} Zhu, J., and Alajaji, F. Rate-distortion analysis of
the binary memoryless source under the Hamming distortion measure
(second-order / dispersion characterisation).
\textbf{{[}119{]}} Hafouta, Y. (2020). Explicit Berry--Esseen and
Edgeworth expansions for inhomogeneous Markov chains and
locally-dependent random variables. Nonlinearity, 33; arXiv:1903.04018.
\textbf{{[}120{]}} Hafouta, Y. (2025). Sharper sequential Berry--Esseen
and Edgeworth expansions for uniformly elliptic inhomogeneous Markov
chains. Preprint, arXiv:2511.06414.
\textbf{{[}121{]}} Dolgopyat, D., and Sarig, O. (2023). Local Limit
Theorems for Inhomogeneous Markov Chains. Lecture Notes in Mathematics,
vol.\ 2331, Springer.
\textbf{{[}122{]}} Dolgopyat, D., and Hafouta, Y. (2022). Edgeworth
expansions for independent bounded integer valued random variables.
Stochastic Processes and their Applications, 152, 486-532.
\textbf{{[}123{]}} Krishnamachari, R. T. (2026). Second-order and
dispersion analysis of finite-state Markov sources; the closed-form
long-run varentropy of the binary Markov chain. Preprint.
\textbf{{[}124{]}} Tasci, M. C., and Kostina, V. (2026). On the lossy
dispersion of sources with memory. Preprint.
\textbf{{[}125{]}} Eswaran, K., and Gastpar, M. (2022). Remote Source
Coding under Gaussian Noise: Dueling Roles of Power and Distortion (and
the growing-Markov-order requirement for universal coding of sources
with memory). IEEE Transactions on Information Theory.
\textbf{{[}126{]}} Elshafiy, A., and Rose, K. (2022). Fundamental limits
and achievability for order-$k$ Markov codebooks in lossy source coding.
IEEE International Symposium on Information Theory (ISIT).
\textbf{{[}127{]}} Diaconis, P., and Saloff-Coste, L. (1996). Logarithmic
Sobolev Inequalities for Finite Markov Chains. Annals of Applied
Probability, 6(3), 695-750.
\textbf{{[}128{]}} Kingman, J. F. C. (1961). A Convexity Property of
Positive Matrices. Quarterly Journal of Mathematics, 12(1), 283-284.
\textbf{{[}129{]}} Nussbaum, R. D. (1986). Convexity and log convexity
for the spectral radius. Linear Algebra and its Applications, 73,
59-122.
\textbf{{[}130{]}} Hennion, H. (1997). Limit theorems for products of
positive random matrices. Annals of Probability, 25(4), 1545-1587.
\textbf{{[}131{]}} Lalley, S. P. (1989). Renewal theorems in symbolic
dynamics, with applications to geodesic flows, non-Euclidean tessellations
and their fractal limits. Acta Mathematica, 163, 1-55.
\textbf{{[}132{]}} G\"otze, F., and Hipp, C. (1983). Asymptotic
Expansions for Sums of Weakly Dependent Random Vectors. Zeitschrift f\"ur
Wahrscheinlichkeitstheorie und verwandte Gebiete, 64(2), 211-239.
\textbf{{[}133{]}} Lahiri, S. N. (1993). Refinements in asymptotic
expansions for sums of weakly dependent random vectors. Annals of
Probability, 21(2), 791-799.
\textbf{{[}134{]}} Avram, F., and Berger, T. (1985). On critical
distortion for Markov sources and the rate-distortion function. IEEE
Transactions on Information Theory.
\textbf{{[}135{]}} Aaronson, J., and Denker, M. (2001). Local limit
theorems for partial sums of stationary sequences generated by
Gibbs--Markov maps. Stochastics and Dynamics, 1(2), 193-237.
\textbf{{[}136{]}} Cohen, J. E. (1981). Convexity of the dominant
eigenvalue of an essentially nonnegative matrix. Proceedings of the
American Mathematical Society, 81(4), 657-658.
\textbf{{[}137{]}} Kato, T. (1982). Superconvexity of the spectral radius,
and convexity of the spectral bound and the type. Mathematische
Zeitschrift, 180(3), 265-273.

\textbf{{[}RNR-I{]}} Ratushnyi, P. (2026). Ratushnyi Neural Repair Coding,
Part I: Framework, Type Taxonomy, Achievability--Converse Matrix, and
Hardness. Companion paper (RNR series, Part I).

\textbf{{[}RNR-II{]}} Ratushnyi, P. (2026). Ratushnyi Neural Repair Coding,
Part II: Random Access and Deviation Theory for RNR Archives. Companion
paper (RNR series, Part II).

\end{document}
